% This is samplepaper.tex, a sample chapter demonstrating the
% LLNCS macro package for Springer Computer Science proceedings;
% Version 2.25 of 2026/09/03
%
\documentclass[runningheads,orivec]{llncs}
%
\usepackage[T1]{fontenc}
\usepackage[utf8]{inputenc}
\let\origvec\vec
\let\vec\relax
\usepackage{amsmath}
\usepackage{amssymb}
\let\vec\origvec
\usepackage{mathptmx}
\makeatletter
\def\boldmath{\relax}
\def\unboldmath{\relax}
\makeatother
\AtBeginDocument{%
  \let\boldmath\relax
  \let\unboldmath\relax
}
\usepackage{graphicx}
\usepackage{url}
\usepackage{longtable}
\makeatletter
% Pure LaTeX2e float barrier (avoids missing placeins.sty in local MiKTeX)
\providecommand{\FloatBarrier}{\par\ifx\@deferlist\@empty\else\clearpage\fi}
\makeatother
\usepackage{color}
\usepackage[colorlinks=true,allcolors=blue,bookmarks=true,bookmarksnumbered=true,bookmarksopen=true,bookmarksopenlevel=2]{hyperref}
\setcounter{tocdepth}{2}
\renewcommand\UrlFont{\color{blue}\rmfamily}
\urlstyle{rm}
\emergencystretch 3em
\hfuzz=100pt
\vfuzz=100pt
\hbadness=10000
\vbadness=10000
\raggedbottom
\setlength\textfloatsep{8mm plus 0.5pt minus 0.5pt}
\setlength\intextsep{8mm plus 0.5pt minus 0.5pt}


\renewcommand{\topfraction}{0.9}
\renewcommand{\bottomfraction}{0.8}
\renewcommand{\textfraction}{0.1}
\renewcommand{\floatpagefraction}{0.85}

\begin{document}
%
\title{How Does Equity Return Predictability Vary Across Horizons and Market Regimes? An Out-of-Sample Analysis}
%
\titlerunning{How Does Equity Return Predictability Vary?}
\author{Doan Duc Hieu Nguyen\orcidID{0009-0008-5897-8147} \and
Thi Phuong Hoa Nguyen\orcidID{0009-0002-7378-5091} \and
Van Phat Nguyen\orcidID{0009-0009-4400-0794} \and
Anh Thu Tran\orcidID{0009-0000-7168-2123} \and
Ngoc Thao Tran\orcidID{0009-0009-6563-9394} \and
Phi Long Duong\orcidID{0000-0000-0000-0000}}
%
\authorrunning{D. D. H. Nguyen et al.}
%
\institute{University of Information Technology, Vietnam National University Ho Chi Minh City,\\
Ho Chi Minh City, Vietnam\\
\email{\{24520500, 24520524, 24521307, 24521643, 24521609\}@gm.uit.edu.vn, longdp@uit.edu.vn}}
\maketitle              % typeset the header of the contribution
%
\begin{abstract}
This study examines how out-of-sample (OOS) predictability of S\&P 500 price-index returns varies across forecast horizons, volatility regimes, and chronological subperiods. Using daily S\&P 500 and VIX data from 1990 to 2026, we evaluate twelve econometric, conventional machine-learning, and deep-learning models against the Historical Mean (HM) benchmark for cumulative log returns over 1, 5, 20, 60, and 120 trading days. The final holdout evaluation contains 1,362 forecast origins. Descriptive results reveal non-monotonic and model-specific performance. TSMixer records positive OOS $R^2$ values at 20, 60, and 120 days, reaching 17.53\% at the longest horizon. However, no model significantly outperforms HM after Holm correction across the 60 model--horizon comparisons. For H1, the cross-horizon equality test rejects only for ARIMA under the baseline HAC specification and alternative bandwidth checks, but this result is not corroborated by the synchronized moving-block bootstrap. For H2, the Normal-volatility regime contains the largest number of positive model--horizon OOS $R^2$ outcomes, while nine Giacomini--White conditional predictive ability tests reject their joint null after Holm correction under the baseline specification, all with negative reported full-sample mean differentials. The horizon--regime interaction is significant only for TCN under the baseline 20-day classification, although its significance varies under alternative volatility-window definitions. For H3, holdout block-level and rolling-window analyses identify model- and horizon-specific temporal variation, but the statistical evidence is sensitive to HAC bandwidth and bootstrap block length. Finally, multi-seed evaluation reveals substantial initialization sensitivity in TSMixer: its mean OOS $R^2$ is negative across all five horizons, reaching $-$18.93\% at 120 days, with only one of ten runs achieving positive OOS $R^2$ at that horizon. Overall, the findings distinguish descriptive forecasting gains from statistically robust and reproducible predictive ability.

\keywords{Equity return predictability \and Multi-horizon forecasting \and Volatility regimes \and Temporal stability \and Machine learning \and Out-of-sample evaluation.}
\end{abstract}


\section{Introduction}\label{sec:intro}

The predictability of equity returns has long been a central question in financial economics, with important implications for asset pricing, portfolio allocation, and risk management. Despite extensive research, forecasting equity returns remains challenging because expected returns are small relative to return volatility, and predictive relationships may vary across investment horizons, market conditions, and historical periods~\cite{campbell1998,welch2008}. Although machine learning provides flexible approaches to capturing complex and potentially nonlinear relationships in financial data~\cite{gu2020}, greater model complexity does not necessarily translate into reliable out-of-sample (OOS) predictive performance. Consequently, understanding the conditions under which equity returns can be predicted remains as important as developing more sophisticated forecasting models.

Previous research suggests that equity return predictability may vary along three related dimensions: forecast horizons, market volatility, and time. Existing studies have documented horizon-dependent predictive performance~\cite{nechvatalova2024,xu2020}, differences in predictive relationships across market volatility conditions~\cite{baltas2018,cakici2025}, and temporal instability in return prediction models~\cite{harvey2021,paye2006}. Related research has also examined the relationship between market volatility, implied volatility, and expected stock returns~\cite{banerjee2007,bekaert2014,french1987}. Collectively, this literature indicates that equity return predictability should not be assessed under a single forecasting horizon, market condition, or evaluation period.

Reliable out-of-sample (OOS) evaluation is essential for assessing whether predictive models provide information beyond conventional historical benchmarks~\cite{campbell2008,welch2008}. Recent research further emphasizes the importance of controlling information leakage and maintaining chronological consistency in financial forecasting, particularly when evaluating multi-period prediction targets~\cite{hewamalage2023,perez2026,tashman2000}. These considerations motivate a unified evaluation framework that applies consistent forecasting and validation procedures across different investment horizons.

Despite these contributions, evidence on equity return predictability remains fragmented across different forecasting horizons, market conditions, data settings, and evaluation procedures. In particular, the joint relationship between forecast horizons and volatility regimes, as well as the temporal stability of these conditional patterns, has not been sufficiently examined within a unified out-of-sample framework for aggregate equity index returns. This motivates an integrated empirical investigation of whether predictive performance varies across horizons, whether such variation depends on market volatility, and whether the resulting patterns persist across chronological sub-periods.

To address this gap, this study investigates OOS equity return predictability using daily S\&P~500 index returns and CBOE Volatility Index (VIX) data spanning 1990--2026~\cite{fredvix,yahoosandp}. The empirical analysis evaluates 12 econometric and machine learning forecasting models across five forecast horizons: 1, 5, 20, 60, and 120 trading days. Predictive performance is examined under three market volatility regimes---low, normal, and high---and across chronological sub-periods to assess the temporal stability of the observed patterns. All forecasting models are evaluated against an expanding historical mean benchmark, following the importance of benchmark-based OOS assessment in equity return prediction~\cite{campbell2008}. The evaluation design also considers the importance of chronological forecasting procedures and appropriate treatment of multi-period forecasting targets~\cite{hewamalage2023,tashman2000}.

Specifically, this study addresses three research questions:
\begin{enumerate}
\item How does OOS equity return predictability vary across forecast horizons?
\item Does the relationship between forecast horizon and predictive performance differ across market volatility regimes?
\item Do the observed horizon- and regime-dependent patterns remain stable across chronological sub-periods?
\end{enumerate}

This study contributes to the empirical literature on equity return predictability in three ways. First, it provides a unified multi-model assessment of return predictability across five forecast horizons, comparing 12 forecasting models from different methodological families against a common historical mean benchmark. Second, it examines the heterogeneity of predictive performance across low-, normal-, and high-volatility regimes, providing evidence on whether the relationship between forecast horizon and predictive performance depends on prevailing market conditions. Third, it investigates the temporal stability of these horizon- and regime-dependent patterns across chronological sub-periods, allowing the analysis to distinguish predictive gains that are relatively persistent from those concentrated in particular periods. Rather than proposing a new forecasting algorithm, this study provides an integrated empirical evaluation of the horizon, regime, and temporal dimensions of equity return predictability within a common OOS framework.

The remainder of this paper is organized as follows. Section~\ref{sec:lit_review} reviews the related literature and develops the research hypotheses. Section~\ref{sec:data} describes the data and research design. Section~\ref{sec:methodology} presents the forecasting models and evaluation methodology. Section~\ref{sec:results} reports and discusses the empirical results, and Section~\ref{sec:conclusion} concludes the study.


\section{Related Work and Hypothesis Development}\label{sec:lit_review}

\subsection{Equity Return Predictability and Model Complexity}

Achieving reliable out-of-sample (OOS) equity return predictability beyond historical benchmarks remains a central challenge in empirical finance. Welch and Goyal~\cite{welch2008} document the limited and time-varying performance of conventional equity premium prediction models, highlighting the difficulty of identifying stable predictive relationships in aggregate stock returns. Campbell and Thompson~\cite{campbell2008} show that imposing economically motivated constraints can improve the OOS performance of predictive regressions. More recently, Gu et al.~\cite{gu2020} demonstrate that machine learning methods can capture nonlinear relationships and interactions in asset pricing, offering an alternative to conventional linear specifications.

However, greater model flexibility does not necessarily translate into superior OOS performance. When predictive signals are weak relative to return volatility, complex models may be more exposed to estimation uncertainty and overfitting. Recent empirical evidence further suggests that increasing model size does not necessarily improve forecasting robustness, as predictive performance can vary substantially across random initializations~\cite{coakley2026}. This creates a trade-off between model flexibility and generalization, motivating the comparison of models with different levels of complexity under a common OOS evaluation framework.

Beyond model specification, the evaluation of predictive performance must also account for the forecast horizon. Quaedvlieg~\cite{quaedvlieg2021} develops multi-horizon forecast comparison procedures that jointly evaluate predictive performance across a forecast path, rather than relying exclusively on separate single-horizon comparisons. This is particularly relevant because cumulative returns over overlapping horizons induce dependence between forecast errors. Hansen and Hodrick~\cite{hansen1980} demonstrate that when the forecast horizon exceeds the sampling interval, overlapping multi-period forecast errors can exhibit serial dependence extending up to $h-1$ lags under the corresponding forecasting framework. This dependence arises because adjacent forecast errors share some of the same unobserved innovations over their overlapping forecast periods. Although OLS estimation may remain consistent under the relevant orthogonality conditions, conventional standard errors that ignore this dependence can lead to unreliable statistical inference. Hodrick~\cite{hodrick1992} further investigates this issue in the context of long-horizon stock return predictability and documents substantial finite-sample distortions in conventional inference procedures, particularly as the forecast horizon increases. These findings motivate the use of dependence-robust covariance estimators and additional finite-sample robustness checks in multi-horizon return prediction. Xu~\cite{xu2020} further examines the statistical challenges of testing predictability across multiple horizons and shows that conventional inference procedures may become unreliable as the horizon increases. These considerations suggest that equity return predictability should be evaluated not only across competing model specifications but also across forecast horizons using procedures that account for horizon-related dependence.

\subsection{Horizon and Volatility-Regime Dependence}

The predictive performance of equity return models may vary across investment horizons because the relevance of financial information can differ between short- and long-term forecasting periods. Nechv{\'a}talov{\'a}~\cite{nechvatalova2024} documents substantial differences in predictive performance across monthly horizons in a machine learning framework for cross-sectional stock returns. Cakici and Zaremba~\cite{cakici2025} also investigate how the relative importance of accounting and technical information varies with the forecasting horizon. Although these studies provide evidence that predictive relationships are horizon-dependent, their settings differ from aggregate time-series forecasting of the equity premium. In particular, cross-sectional predictability does not necessarily imply the same horizon-related patterns in aggregate market returns.

The comparison of predictive performance across horizons also requires careful interpretation. A decline in statistical forecast accuracy at longer horizons does not necessarily imply a corresponding decline in economic value, since portfolio turnover and transaction costs may also change with the investment horizon. Cenesizoglu and Timmermann~\cite{cenesizoglu2012} demonstrate that statistical forecast accuracy and economic performance can differ substantially, with models exhibiting weaker RMSE performance sometimes generating higher certainty equivalent returns, Sharpe ratios, or risk-adjusted returns. Similarly, Balduzzi and Lynch~\cite{balduzzi1999} show that transaction costs can substantially reduce the economic benefits of return predictability, particularly when frequent portfolio rebalancing is required. Nechv{\'a}talov{\'a}~\cite{nechvatalova2024}, for example, illustrates that statistical predictability and risk-adjusted portfolio performance need not exhibit identical patterns across horizons. Accordingly, the present study focuses on OOS predictive performance rather than assuming that longer horizons must produce uniformly weaker forecasting results or necessarily translate into lower investment value.

Market volatility provides another important dimension of conditional predictability. Baltas and Karyampas~\cite{baltas2018} show that equity risk premium forecasting performance varies across market volatility regimes, with model performance differing substantially between low- and high-volatility conditions. Related evidence from Cakici and Zaremba~\cite{cakici2025} suggests that the predictive importance of information sources may also depend on volatility conditions. Other studies examine conditional predictability through alternative definitions of market states. For example, Li and Tsiakas~\cite{li2017} investigate differences between recessionary and expansionary periods, while Hammerschmid and Lohre~\cite{hammerschmid2018} construct macroeconomic regime factors to explain variation in aggregate stock return predictability. These distinctions are important because business-cycle regimes, macroeconomic regimes, and volatility regimes capture related but conceptually different dimensions of market conditions.

Although prior research establishes that predictive performance can vary across volatility or economic regimes, much of the regime-dependent equity premium literature focuses on a single short-term forecasting horizon. Conversely, multi-horizon studies generally examine how predictive performance changes with the horizon without explicitly conditioning the horizon--performance relationship on market volatility. This separation motivates a joint evaluation of forecast horizons and volatility regimes. In particular, it remains important to examine whether the pattern of predictive performance across horizons differs between low- and high-volatility market conditions.

\subsection{Temporal Stability and Hypothesis Development}

Equity return predictability may also evolve over time as economic conditions, market structures, monetary policies, and investor behavior change. Paye and Timmermann~\cite{paye2006} document substantial instability in return prediction models and identify structural changes in the relationship between predictors and future returns. Harvey et al.~\cite{harvey2021} further examine the real-time detection of regimes of predictability in the US equity premium, highlighting that periods of predictive ability may emerge and subsequently disappear. These findings suggest that aggregate OOS performance can conceal substantial variation across chronological subperiods.

Structural instability may also have implications for investors operating over different investment horizons. Pettenuzzo and Timmermann~\cite{pettenuzzo2011} show that accounting for structural breaks in return prediction models can materially alter long-horizon asset allocation decisions. Their findings emphasize that model instability is relevant not only to the estimation of predictive relationships but also to the interpretation of long-term investment outcomes. Nevertheless, structural-break analysis and chronological OOS evaluation address different aspects of temporal variation. The former identifies changes in model parameters or the underlying data-generating process, whereas the latter assesses whether predictive performance remains consistent across successive evaluation periods.

More importantly, temporal instability may affect not only overall predictive performance but also the relationship between forecast horizons and market volatility. For example, a forecasting model may perform relatively well at short horizons during low-volatility periods but lose its advantage during high-volatility periods. Such a pattern may also change across chronological subperiods. Therefore, examining temporal stability requires evaluating whether horizon-dependent and volatility-conditioned predictive patterns persist across different OOS evaluation windows, rather than relying solely on full-sample performance measures.

Building on the above literature, this study develops three hypotheses to examine the horizon dependence, volatility conditioning, and temporal stability of equity return predictability.

\begin{quote}
\noindent\textbf{H1:} \textit{Out-of-sample predictive performance differs significantly across forecast horizons.}

\medskip
\noindent\textbf{H2:} \textit{Out-of-sample predictive performance relative to the Historical Mean benchmark exhibits regime-conditional predictability across market volatility states, and the relative performance contrasts across volatility regimes vary with the forecast horizon.}

\medskip
\noindent\textbf{H3:} \textit{Out-of-sample predictive performance and its horizon--volatility regime patterns exhibit temporal instability across chronological subperiods.}
\end{quote}

Collectively, these hypotheses provide an empirical framework for investigating whether equity return predictability varies across investment horizons, whether such variation depends on prevailing market volatility, and whether the resulting patterns remain stable over time.


\section{Research Design and Methodology}\label{sec:data}

\subsection{Data Sources and Sample Construction}\label{subsec:data_sources}

This study uses daily financial data from the S\&P~500 index and the CBOE Volatility Index (VIX). Historical S\&P~500 data are obtained from Yahoo Finance~\cite{yahoosandp}, while VIX data are retrieved from the Federal Reserve Economic Data (FRED) database under the series identifier \texttt{VIXCLS}~\cite{fredvix}. The raw datasets contain 9,253 S\&P~500 observations (January 2, 1990--September 29, 2026) and 9,585 VIX observations (January 2, 1990--September 28, 2026), respectively.

The predictor variables are organized into six groups: lagged returns, momentum, historical volatility, VIX-related indicators, intraday indicators, and trading activity. These groups draw on prior research on return dynamics and volatility~\cite{french1987}, momentum, including cross-sectional momentum~\cite{jegadeesh1993} and time-series momentum~\cite{moskowitz2012}, implied volatility and market uncertainty~\cite{banerjee2007,bekaert2014}, trading volume~\cite{campbell1993}, and range-based volatility estimation~\cite{alizadeh2002}. In total, the study employs 16 predictors, whose group-level composition is summarized in Table~\ref{tab:predictors}. Their exact definitions and construction procedures are provided in Section~\ref{subsec:targets_predictors}.

\begin{table}[htbp]
\caption{Summary of predictor variables across the six feature groups.}\label{tab:predictors}
\centering
\begin{tabular}{lcl}
\hline
Feature group & Number of predictors & Information captured \\
\hline
Lagged returns & 3 & Past return observations \\
Momentum & 3 & Cumulative return trends \\
Historical volatility & 3 & Recent return dispersion \\
VIX-related indicators & 4 & Implied market uncertainty \\
Intraday indicators & 2 & Intraday price movements \\
Trading activity & 1 & Changes in reported volume \\
\hline
Total & 16 & \\
\hline
\end{tabular}
\end{table}

Following Campbell, Lo, and MacKinlay~\cite{campbell1998}, continuously compounded returns (log returns) are defined as the natural logarithm of the ratio of consecutive asset prices. Since the S\&P~500 price index is used without incorporating dividend payments, the one-day log return is calculated as:
\begin{equation}\label{eq:daily_log_return}
R_t^{(1)} = \ln\left(\frac{P_t}{P_{t-1}}\right) = \ln(P_t) - \ln(P_{t-1}),
\end{equation}
where $P_t$ and $P_{t-1}$ denote the S\&P~500 closing index levels on trading days $t$ and $t-1$, respectively. The ``Close'' and ``Adj Close'' series are identical in the downloaded dataset. The first return observation is computed using the available open-to-close price movement but is subsequently excluded during feature construction.

The two datasets are aligned using the S\&P~500 trading calendar through a left join, with missing VIX observations forward-filled, resulting in 9,252 observations. Following the construction of forecasting targets and historical features, and the removal of observations with unavailable values, the final modeling dataset contains 9,073 observations from March~27, 1990, to April~7, 2026. The first 60 observations are excluded due to rolling-feature requirements, while the final 119 observations are unavailable for constructing the 120-day forward-return target.

\subsection{Forecasting Targets and Predictor Variables}\label{subsec:targets_predictors}

The study examines five cumulative forward-return horizons: 1, 5, 20, 60, and 120 trading days. For each horizon $h$, the forecasting target is defined as:
\begin{equation}\label{eq:target_return}
Y_t^{(h)} = \ln\left(\frac{P_{t+h}}{P_t}\right) = \sum_{i=1}^h R_{t+i}^{(1)}, \qquad h \in \{1, 5, 20, 60, 120\},
\end{equation}
where $P_t$ and $P_{t+h}$ denote the S\&P~500 closing index levels on trading days $t$ and $t+h$, respectively, and $R_{t+i}^{(1)}$ represents the one-day log return between trading days $t+i-1$ and $t+i$. Thus, the target measures the cumulative log return from the close of day $t$ to the close of day $t+h$.

\noindent\textbf{Predictor construction.} The 16 predictors are constructed from six groups of market information: lagged returns, momentum, historical volatility, VIX-related indicators, intraday price movements, and trading activity. Table~\ref{tab:predictor_definitions} presents the exact definitions and calculation procedures. All rolling statistics are calculated using observations available up to and including date $t$.

\begin{table}[htbp]
\centering
\caption{Definition and construction of predictor variables.}\label{tab:predictor_definitions}
\setlength{\tabcolsep}{4pt}
\begin{tabular}{lll}
\hline
Predictor group & Predictor & Definition \\
\hline
Lagged returns & \texttt{Return\_lag\_1} & $R_{t-1}^{(1)}$ \\
& \texttt{Return\_lag\_5} & $R_{t-5}^{(1)}$ \\
& \texttt{Return\_lag\_20} & $R_{t-20}^{(1)}$ \\
\hline
Momentum & \texttt{Momentum\_5D} & $\sum_{j=0}^{4} R_{t-j}^{(1)} = \ln(P_t/P_{t-5})$ \\
& \texttt{Momentum\_20D} & $\sum_{j=0}^{19} R_{t-j}^{(1)} = \ln(P_t/P_{t-20})$ \\
& \texttt{Momentum\_60D} & $\sum_{j=0}^{59} R_{t-j}^{(1)} = \ln(P_t/P_{t-60})$ \\
\hline
Historical volatility & \texttt{Volatility\_5D} & Sample standard deviation of $R_{t-4}^{(1)},\dots,R_t^{(1)}$ \\
& \texttt{Volatility\_20D} & Sample standard deviation of $R_{t-19}^{(1)},\dots,R_t^{(1)}$ \\
& \texttt{Volatility\_60D} & Sample standard deviation of $R_{t-59}^{(1)},\dots,R_t^{(1)}$ \\
\hline
VIX-related indicators & \texttt{VIXCLS} & $\text{VIX}_t$ \\
& \texttt{VIX\_lag\_1} & $\text{VIX}_{t-1}$ \\
& \texttt{VIX\_change\_1D} & $\text{VIX}_t - \text{VIX}_{t-1}$ \\
& \texttt{VIX\_pct\_change\_1D} & $(\text{VIX}_t - \text{VIX}_{t-1})/\text{VIX}_{t-1}$ \\
\hline
Intraday indicators & \texttt{HL\_range} & $(\text{High}_t - \text{Low}_t)/\text{Close}_t$ \\
& \texttt{OC\_return} & $(\text{Close}_t - \text{Open}_t)/\text{Open}_t$ \\
\hline
Trading activity & \texttt{Volume\_change} & $(\text{Volume}_t - \text{Volume}_{t-1})/\text{Volume}_{t-1}$ \\
\hline
\end{tabular}
\par\vspace{4pt}\parbox{\linewidth}{\scriptsize\textit{Note:} The rolling standard deviations used for the three historical-volatility predictors are calculated using the sample standard deviation (denominator $k-1$, consistent with the default pandas implementation). They are daily, non-annualized measures.}
\end{table}

The lagged-return predictors capture returns observed in previous trading sessions. Momentum predictors measure cumulative log returns over the corresponding rolling windows, reflecting the time-series momentum perspective, which focuses on an asset's own past returns rather than its relative performance against other assets~\cite{moskowitz2012}. Accordingly, the three momentum variables capture the historical return trends of the S\&P 500 index itself rather than cross-sectional rankings of individual securities, as examined in~\cite{jegadeesh1993}. Historical-volatility predictors measure the dispersion of daily log returns within the corresponding rolling windows. In particular, the historical-volatility predictors are calculated using daily observations and are not annualized. Consistent with this convention, the 20-day rolling sample standard deviation (\texttt{Volatility\_20D}) is also used as the underlying non-annualized volatility measure for regime classification in Section~\ref{subsec:regime_classification}, where expanding historical quantiles convert it into ex-post evaluation regime labels that are excluded from the predictor matrix.

The four VIX-related predictors capture different aspects of option-implied market uncertainty. \texttt{VIXCLS} represents the VIX level associated with date $t$, while \texttt{VIX\_lag\_1} captures its previous-session level. \texttt{VIX\_change\_1D} measures the absolute daily change in the index, whereas \texttt{VIX\_pct\_change\_1D} measures its relative daily change.

The intraday indicators and trading-activity predictor are calculated as follows:
\begin{equation}\label{eq:hl_range}
\text{HL\_range}_t = \frac{\text{High}_t - \text{Low}_t}{\text{Close}_t},
\end{equation}
\begin{equation}\label{eq:oc_return}
\text{OC\_return}_t = \frac{\text{Close}_t - \text{Open}_t}{\text{Open}_t},
\end{equation}
\begin{equation}\label{eq:volume_change}
\text{Volume\_change}_t = \frac{\text{Volume}_t - \text{Volume}_{t-1}}{\text{Volume}_{t-1}}.
\end{equation}
Here, $\text{High}_t$, $\text{Low}_t$, $\text{Open}_t$, and $\text{Close}_t$ denote the corresponding daily S\&P~500 index levels, while $\text{Volume}_t$ denotes the volume value reported in the downloaded dataset. The high--low range captures the relative intraday price range, and the open--close return measures the simple price change during the trading session. Volume change is included as a proxy for market trading activity rather than as a direct measure of aggregate constituent-stock liquidity.

A zero-volume observation on May~24, 2023, is treated as a data anomaly. It is replaced with the preceding valid volume observation before calculating percentage changes. This procedure prevents an invalid zero denominator from affecting the volume-change feature and preserves the previous valid observation when calculating the subsequent change.

\noindent\textbf{Information availability and temporal alignment.} All 16 predictors are constructed using market observations associated with date $t$, without incorporating future observations. Four predictors---\texttt{Return\_lag\_1}, \texttt{Return\_lag\_5}, \texttt{Return\_lag\_20}, and \texttt{VIX\_lag\_1}---use information from previous trading sessions and are therefore available before the current trading session begins. The remaining 12 predictors incorporate observations associated with date $t$, including its completed price, volume, and VIX observations. Accordingly, the predictor vector is interpreted as a post-close information set rather than as information necessarily available before or during the closing auction.

Following the alignment procedure described in \hyperref[subsec:data_sources]{Section~\ref*{subsec:data_sources}}, the VIX series is matched to the S\&P~500 trading calendar using a left join, with missing observations forward-filled using the most recent preceding available value. This procedure preserves temporal consistency by carrying forward historical information without introducing future VIX observations. On dates requiring forward-filling, the VIX level remains unchanged, resulting in zero daily changes for the affected VIX-derived change measures.

The forecasting framework therefore evaluates the statistical predictability of subsequent cumulative returns conditional on information associated with the end of date $t$. Since some predictors, particularly current-day volume and the completed VIX observation, are only available after the relevant market observations have been finalized, the framework does not assume that a trading position can be executed at the observed closing price $P_t$. An executable trading strategy would require a subsequent execution opportunity and would consequently involve additional risks, including overnight price movements.

\subsection{Volatility Regime Classification}\label{subsec:regime_classification}

To examine regime-dependent predictability, observations are classified into Low, Normal, and High volatility regimes using the non-annualized 20-day rolling sample standard deviation of daily market log returns (\texttt{Volatility\_20D}):
\begin{equation}\label{eq:volatility}
\sigma_t^{(20)} = \operatorname{Std}_{\mathrm{sample}}\left(R_{t-19}^{(1)}, \dots, R_t^{(1)}\right).
\end{equation}
Expanding historical quantiles are employed to determine the volatility thresholds using information available up to time $t$~\cite{feng2024}:
\begin{equation}\label{eq:quantiles}
Q_{p,t} = \operatorname{Quantile}_p\left(\left\{\sigma_\tau^{(20)} : \tau \le t\right\}\right), \qquad p \in \{0.33, 0.67\}.
\end{equation}
A minimum of 252 historical observations is required before calculating the thresholds, with earlier observations back-filled using the first available expanding quantile thresholds. Observations are classified as Low when $\sigma_t^{(20)} \le Q_{0.33,t}$, Normal when $Q_{0.33,t} < \sigma_t^{(20)} \le Q_{0.67,t}$, and High when $\sigma_t^{(20)} > Q_{0.67,t}$.

Across the final sample, 3,062 observations (33.75\%) are classified as Low, 2,753 (30.34\%) as Normal, and 3,258 (35.91\%) as High volatility. Although the continuous underlying series $\sigma_t^{(20)}$ matches the 20-day historical-volatility predictor in Section~\ref{subsec:targets_predictors}, the discretized categorical regime labels are excluded from the predictor matrix and are not supplied to the forecasting models. All forecasting models are estimated globally using the same 16 predictors, with regime labels used exclusively for ex-post conditional evaluation.

The 20-day rolling window serves as the baseline specification for volatility-regime classification. To assess whether the regime-dependent findings are sensitive to the choice of rolling-window length, additional robustness analyses are conducted using 10-day and 60-day rolling standard deviations. These alternative specifications retain the same expanding-quantile classification procedure and are evaluated on the same date-aligned holdout sample. The corresponding results are reported in Section~\ref{subsec:volatility_regimes}.

\subsection{Chronological Data Splitting and Walk-Forward Validation}\label{subsec:data_splitting}

The dataset is divided chronologically into three non-overlapping subsets: training (6,351 observations; March~27, 1990--June~9, 2015), validation (1,360 observations; June~10, 2015--October~30, 2020), and hold-out testing (1,362 observations; November~2, 2020--April~7, 2026).

The validation period is evaluated through five consecutive windows of 272 observations using an expanding-window walk-forward procedure, following the rolling-origin framework of Tashman~\cite{tashman2000}. For each forecast horizon $h$, the final $h$ training observations are purged to ensure that training targets do not extend into the corresponding evaluation period:
\begin{equation}\label{eq:train_purged}
N_{\mathrm{train,purged}}^{(h)} = N_{\mathrm{train}} - h.
\end{equation}
At the transition to hold-out testing, the training sample is expanded to include both the original training and validation subsets, followed by horizon-specific purging. \texttt{StandardScaler} is fitted exclusively on the corresponding training data and then applied to validation or test observations, preventing information from the evaluation periods from influencing feature scaling, as emphasized by Hewamalage, Ackermann, and Bergmeir~\cite{hewamalage2023}. The same set of 16 predictors defined in Section~\ref{subsec:targets_predictors} is used consistently across the forecasting models and evaluation horizons.

For RQ3, the validation period is further divided into three non-overlapping calendar-based subperiods: 2015--2016 (395 observations), 2017--2018 (502 observations), and 2019--2020 (463 observations). This validation-period decomposition is used to examine temporal variation in predictive performance during model development and validation. The hold-out test period is separately organized into four subperiods: November--December 2020 (42 observations), 2021--2022 (503 observations), 2023--2024 (502 observations), and 2025--2026 (315 observations), providing a separate basis for evaluating temporal stability on data not used for model selection or hyperparameter optimization. Although the calendar labels overlap for 2019--2020, the observations are assigned according to the original chronological split boundaries and do not overlap between subsets.

The hold-out test set is reserved exclusively for final out-of-sample evaluation. Model configurations and hyperparameters are determined using training and validation data, with no model selection, hyperparameter optimization, or feature-engineering decisions based on hold-out performance. This design supports the evaluation of horizon-dependent predictability (RQ1), volatility-regime differences (RQ2), and temporal stability (RQ3) across five forecast horizons and three volatility regimes.


\section{Forecasting Models and Evaluation Framework}\label{sec:methodology}

\subsection{Forecasting Models and Benchmark}\label{subsec:models_benchmark}

This study evaluates twelve forecasting models across five forecast horizons: 1, 5, 20, 60, and 120 trading days. The models represent four methodological categories: econometric models (ARIMA and ARIMAX), penalized linear models (Ridge Regression and Elastic Net), kernel- and tree-based models (Support Vector Regression (SVR), Random Forest (RF), and XGBoost), and deep learning architectures (Long Short-Term Memory (LSTM), Temporal Convolutional Network (TCN), Vanilla Transformer, PatchTST, and TSMixer).

The selected models provide a comparative assessment of linear, nonlinear, and neural network-based approaches to equity return forecasting. ARIMA captures temporal dependence in the target return series, whereas ARIMAX incorporates additional explanatory information through exogenous predictors. Ridge Regression and Elastic Net examine the predictive contribution of regularized linear relationships. SVR, RF, and XGBoost provide kernel-based and tree-based alternatives for modeling nonlinear patterns. The deep learning models represent recurrent, convolutional, attention-based, patch-based, and feature-temporal mixing approaches to sequential forecasting.

The econometric models use the target return series and, in the case of ARIMAX, ten exogenous predictors: \textit{VIXCLS}, \textit{VIX\_lag\_1}, \textit{VIX\_change\_1D}, \textit{VIX\_pct\_change\_1D}, \textit{Momentum\_5D}, \textit{Momentum\_20D}, \textit{Volatility\_20D}, \textit{HL\_range}, \textit{OC\_return}, and \textit{Volume\_change}. The remaining ten machine learning and deep learning models use the sixteen predictors described in \hyperref[sec:data]{Section~\ref*{sec:data}}. A separate forecasting model is estimated for each forecast horizon.

The Historical Mean (HM) serves as the benchmark for evaluating out-of-sample forecasting performance~\cite{welch2008}. For each forecast horizon $h$, the benchmark forecast is calculated as the arithmetic mean of the eligible training targets:
\begin{equation}\label{eq:historical_mean}
\widehat{\mu}_{\mathrm{HM}}^{(h)}
=
\frac{1}{N_{\mathrm{train,purged}}^{(h)}}
\sum_{t\in \mathcal{T}_{\mathrm{train,purged}}^{(h)}}
Y_t^{(h)}.
\end{equation}
The HM benchmark is re-estimated independently for each validation fold and forecast horizon. During the final holdout evaluation, it is estimated using the combined purged training and validation sample. The estimated mean is positive across all evaluation windows, implying that HM consistently generates positive return forecasts within each window.

\subsection{Model Configuration and Hyperparameter Optimization}\label{subsec:model_config}

The forecasting models are configured using model-specific specifications and optimization procedures. This section presents the hyperparameter selection strategies and training settings adopted for each model.

\paragraph{Econometric model configurations.} ARIMA considers three candidate specifications: ARIMA(1,0,0), ARIMA(0,0,1), and ARIMA(1,0,1). The final specification is selected according to the Akaike Information Criterion (AIC), without differencing. ARIMAX adopts a fixed ARMA(1,1) specification and incorporates the ten exogenous predictors described in \hyperref[subsec:models_benchmark]{Section~\ref*{subsec:models_benchmark}}. The final specifications selected for each forecasting horizon are summarized in \hyperref[tab:config_econ_ml]{Table~\ref*{tab:config_econ_ml}}. For ARIMA, the selected order is determined independently for each horizon using the AIC criterion. ARIMAX retains the fixed ARMA(1,1) specification with the ten exogenous predictors across all horizons.

\paragraph{Conventional machine learning configurations.} Ridge Regression searches over the following regularization strengths:
\begin{equation}
\alpha \in \{10^{-3},10^{-2},10^{-1},1,10,100\}.
\end{equation}
Elastic Net jointly considers the regularization strength and the mixing parameter, with the search space defined as:
\begin{equation}
\begin{aligned}
\alpha &\in \operatorname{LogUniform}(10^{-5},1),\\
\ell_1\text{-ratio} &\in \{0.1,0.3,0.5,0.7,0.9\}.
\end{aligned}
\end{equation}
For SVR, the candidate hyperparameters include the penalty parameter ($C$), the kernel coefficient ($\gamma$), and the insensitive loss parameter ($\varepsilon$):
\begin{equation}
\begin{aligned}
C &\in \{0.1,1,10,100\},\\
\gamma &\in \{\text{scale},0.001,0.01,0.1\},
\qquad
\varepsilon \in \{0.001,0.01,0.05,0.1\}.
\end{aligned}
\end{equation}
For RF, the search space covers the number of estimators, maximum tree depth, minimum samples required to split an internal node, minimum samples required at a leaf node, and the number of features considered at each split. The candidate values include 200 and 500 estimators; maximum depths of unrestricted, 5, 10, and 20; minimum split sizes of 2, 5, and 10; minimum leaf sizes of 1, 2, and 5; and feature-selection strategies based on the square root of the available features or a proportion of 0.5. XGBoost is optimized over its corresponding tree-based hyperparameter space using the optimization procedure described below. The final hyperparameter configurations of the conventional machine learning models are reported in \hyperref[tab:config_econ_ml]{Tables~\ref*{tab:config_econ_ml}} and~\hyperref[tab:config_trees]{\ref*{tab:config_trees}}. These values represent the configurations selected for the final holdout experiments rather than the candidate search spaces used during optimization.

\begin{table}[htbp]
\centering
\caption{Final configurations of econometric and conventional regression models.}\label{tab:config_econ_ml}
\setlength{\tabcolsep}{4.5pt}
\begin{tabular}{llccccc}
\hline
Model & Hyperparameter & 1D & 5D & 20D & 60D & 120D \\
\hline
ARIMA       & Order               & (1,0,1) & (1,0,1) & (1,0,0) & (1,0,1) & (1,0,1) \\
ARIMAX      & Order               & (1,0,1) & (1,0,1) & (1,0,1) & (1,0,1) & (1,0,1) \\
ARIMAX      & Exogenous variables & 10      & 10      & 10      & 10      & 10 \\
Ridge       & Alpha               & 100     & 100     & 100     & 100     & 100 \\
Elastic Net & Alpha               & 0.001   & 0.001   & 0.005   & 0.05    & 0.05 \\
Elastic Net & L1 ratio            & 0.5     & 0.9     & 0.9     & 0.3     & 0.5 \\
SVR         & Kernel              & RBF     & RBF     & RBF     & RBF     & RBF \\
SVR         & $C$                 & 0.1     & 0.1     & 0.1     & 0.1     & 0.1 \\
SVR         & $\gamma$            & 0.001   & 0.001   & 0.001   & 0.001   & 0.1 \\
SVR         & $\varepsilon$       & 0.01    & 0.01    & 0.01    & 0.001   & 0.1 \\
\hline
\end{tabular}
\par\vspace{4pt}\parbox{\linewidth}{\scriptsize\textit{Notes:} ARIMA specifications are selected using AIC from three candidate models. ARIMAX uses ten exogenous predictors. Elastic Net uses \texttt{max\_iter} = 10,000. SVR uses the RBF kernel.}
\end{table}

\begin{table}[htbp]
\centering
\caption{Final configurations of tree-based machine learning models.}\label{tab:config_trees}
\setlength{\tabcolsep}{5pt}
\begin{tabular}{lrrrrr}
\hline
Hyperparameter & 1D & 5D & 20D & 60D & 120D \\
\hline
\multicolumn{6}{l}{\textbf{(a) Random Forest}} \\
\texttt{n\_estimators}      & 200    & 500    & 200    & 500    & 500 \\
\texttt{max\_depth}         & 5      & 5      & 5      & 5      & 5 \\
\texttt{min\_samples\_split}& 5      & 10     & 10     & 2      & 2 \\
\texttt{min\_samples\_leaf} & 1      & 2      & 2      & 2      & 5 \\
\texttt{max\_features}      & sqrt   & 0.5    & sqrt   & sqrt   & sqrt \\
\hline
\multicolumn{6}{l}{\textbf{(b) XGBoost}} \\
\texttt{n\_estimators}      & 200    & 500    & 300    & 900    & 500 \\
\texttt{learning\_rate}     & 0.0935 & 0.0449 & 0.1915 & 0.0902 & 0.0403 \\
\texttt{max\_depth}         & 9      & 5      & 9      & 3      & 3 \\
\texttt{min\_child\_weight} & 1      & 4      & 4      & 6      & 1 \\
\texttt{subsample}          & 0.722  & 0.694  & 0.642  & 0.663  & 0.772 \\
\texttt{colsample\_bytree}  & 0.656  & 0.780  & 0.625  & 0.807  & 0.932 \\
\texttt{gamma}              & 0.0015 & 0.0382 & 0.0990 & 4.2589 & 3.6227 \\
\texttt{reg\_alpha}         & 0.1411 & 0.0083 & 0.0002 & 0.0469 & 0.1260 \\
\texttt{reg\_lambda}        & 2.8797 & 0.1562 & 16.5550 & 0.2765 & 0.0171 \\
\hline
\end{tabular}
\par\vspace{4pt}\parbox{\linewidth}{\scriptsize\textit{Notes:} The reported values are the final configurations selected for the holdout evaluation. The hyperparameters are selected independently for each forecasting horizon.}
\end{table}

\paragraph{Hyperparameter optimization.} Hyperparameter optimization is conducted using two complementary strategies: grid-based search for conventional statistical and machine learning models, and Optuna-based optimization for selected tree-based and deep learning models.

For Optuna-based optimization, a Tree-structured Parzen Estimator (TPE) sampler with a fixed random seed of 42 is employed~\cite{akiba2019}. The optimization objective is to minimize the validation root mean squared error (RMSE) on an internal purged validation split, consisting of the final 15\% of the training data within each fold. This split is used exclusively for hyperparameter selection and is separated from the corresponding training subset to reduce temporal information leakage. The optimization budget is 30 trials for Random Forest (RF), 50 trials for XGBoost, and 30 trials each for LSTM and Vanilla Transformer.

The candidate search spaces for the four Optuna-optimized models are summarized in \hyperref[tab:optuna_search_spaces]{Table~\ref*{tab:optuna_search_spaces}}. For RF, the search space includes the number of estimators, maximum tree depth, minimum split and leaf sizes, and feature-selection strategy. XGBoost explores tree complexity, learning rate, sampling ratios, and regularization parameters. For LSTM and Vanilla Transformer, the search spaces additionally cover sequence length, architectural dimensions, dropout rates, learning rates, and batch sizes. During the Optuna stage, LSTM and Vanilla Transformer are trained for a maximum of 15 epochs per trial to limit computational costs. Their selected configurations are subsequently used in the final training procedure.

For the deep learning models, the \texttt{weight\_decay} parameter is not included in the actual Optuna search because it is not passed to the trial suggestion procedure in the corresponding model implementations. Accordingly, in the final holdout evaluation, LSTM and Vanilla Transformer use the Adam optimizer without explicitly specifying weight decay. The reported search spaces therefore reflect the hyperparameters actually explored by the implementation rather than all parameters declared in the configuration file.

The remaining models follow their respective predefined selection procedures. Ridge Regression, Elastic Net, and SVR use the candidate values specified in the conventional machine learning configurations above. ARIMA selects among three candidate specifications using the Akaike Information Criterion (AIC), while ARIMAX retains its predefined ARMA(1,1) specification with ten exogenous predictors. TCN, PatchTST, and TSMixer adopt predefined standardized benchmark configurations rather than undergoing the same Optuna-based optimization procedure. Consequently, these architectures do not receive an equivalent hyperparameter search budget to LSTM and Vanilla Transformer, a distinction that should be considered when interpreting comparisons among the deep learning models.

\begin{table}[htbp]
\centering
\caption{Hyperparameter search spaces and optimization budgets for Optuna-based models.}\label{tab:optuna_search_spaces}
\setlength{\tabcolsep}{4.5pt}
\begin{tabular}{lllr}
\hline
Model & Hyperparameter & Search space & Trials \\
\hline
RF & \texttt{n\_estimators}      & $\{200, 500\}$                    & 30 \\
   & \texttt{max\_depth}         & $\{\text{None}, 5, 10, 20\}$      & \\
   & \texttt{min\_samples\_split}& $\{2, 5, 10\}$                    & \\
   & \texttt{min\_samples\_leaf} & $\{1, 2, 5\}$                     & \\
   & \texttt{max\_features}      & $\{\text{sqrt}, 0.5\}$            & \\
\hline
XGBoost & \texttt{n\_estimators}     & $[100, 1000]$, step 100           & 50 \\
        & \texttt{learning\_rate}    & $\operatorname{LogUniform}(0.01, 0.2)$ & \\
        & \texttt{max\_depth}        & $[2, 10]$                         & \\
        & \texttt{min\_child\_weight}& $[1, 10]$                         & \\
        & \texttt{subsample}         & $[0.6, 1.0]$                      & \\
        & \texttt{colsample\_bytree} & $[0.6, 1.0]$                      & \\
        & \texttt{gamma}             & $[0, 5]$                          & \\
        & \texttt{reg\_alpha}        & $\operatorname{LogUniform}(10^{-8}, 10)$ & \\
        & \texttt{reg\_lambda}       & $\operatorname{LogUniform}(10^{-3}, 100)$ & \\
\hline
LSTM & \texttt{lookback}      & $\{10, 20, 40, 60\}$                        & 30 \\
     & \texttt{hidden\_size}   & $\{32, 64, 128\}$                           & \\
     & \texttt{num\_layers}    & $\{1, 2, 3\}$                               & \\
     & \texttt{dropout}       & $[0, 0.5]$                                  & \\
     & \texttt{learning\_rate} & $\operatorname{LogUniform}(10^{-5}, 10^{-3})$ & \\
     & \texttt{batch\_size}    & $\{16, 32, 64\}$                            & \\
\hline
Vanilla Transformer & \texttt{lookback}           & $\{20, 40, 60\}$                                   & 30 \\
                    & \texttt{d\_model}           & $\{32, 64\}$                                       & \\
                    & \texttt{nhead}              & $\{2, 4\}$                                         & \\
                    & \texttt{num\_encoder\_layers}& $\{1, 2\}$                                        & \\
                    & \texttt{dim\_feedforward}   & $\{64, 128\}$                                      & \\
                    & \texttt{dropout}            & $[0.05, 0.3]$                                      & \\
                    & \texttt{learning\_rate}      & $\operatorname{LogUniform}(10^{-5}, 5\times 10^{-4})$ & \\
                    & \texttt{batch\_size}         & $\{16, 32\}$                                       & \\
\hline
\end{tabular}
\par\vspace{4pt}\parbox{\linewidth}{\scriptsize\textit{Notes:} Optuna optimization uses a TPE sampler with seed 42 and minimizes validation RMSE on an internal purged split containing the final 15\% of each fold's training data. LSTM and Vanilla Transformer are trained for 15 epochs per trial. $\operatorname{LogUniform}$ denotes logarithmic sampling over the specified interval. Weight decay is excluded from the actual LSTM and Vanilla Transformer search spaces because it is not suggested by their respective Optuna implementations.}
\end{table}

\paragraph{Deep learning training configuration.} The deep learning architectures employ optimizer configurations according to their respective training implementations. LSTM and Vanilla Transformer use the Adam optimizer without explicitly specifying weight decay~\cite{kingma2014}, whereas TCN, PatchTST, and TSMixer use AdamW with a weight decay of $10^{-4}$~\cite{loshchilov2018decoupled}.

The learning rates and batch sizes follow the selected or predefined configurations for each model and forecasting horizon. TCN, PatchTST, and TSMixer retain identical training configurations across all five horizons, whereas LSTM and Vanilla Transformer use horizon-specific configurations.

The maximum training duration, early stopping criteria, and optimization procedures are specified in \hyperref[tab:config_lstm_transformer]{Tables~\ref*{tab:config_lstm_transformer}} and~\hyperref[tab:config_dl_predefined]{\ref*{tab:config_dl_predefined}}. Early stopping with a patience of seven epochs is applied based on validation MSE. Gradient clipping with a maximum norm of 1.0 is used to improve optimization stability, and all deep learning models are trained using the MSE loss function. A fixed random seed of 42 is used for the baseline training procedure. The complete configurations implemented in the holdout evaluation are reported in \hyperref[tab:config_lstm_transformer]{Tables~\ref*{tab:config_lstm_transformer}} and~\hyperref[tab:config_dl_predefined]{\ref*{tab:config_dl_predefined}}.

\begin{table}[htbp]
\centering
\caption{Final configurations of LSTM and Vanilla Transformer.}\label{tab:config_lstm_transformer}
\setlength{\tabcolsep}{4.5pt}
\begin{tabular}{lrrrrr}
\hline
Hyperparameter & 1D & 5D & 20D & 60D & 120D \\
\hline
\multicolumn{6}{l}{\textbf{(a) LSTM}} \\
Lookback                & 20      & 40      & 40      & 10      & 40 \\
Hidden size             & 64      & 64      & 64      & 64      & 64 \\
Num.\ layers            & 3       & 1       & 3       & 3       & 3 \\
Dropout                 & 0.3539  & 0.4357  & 0.3861  & 0.4203  & 0.2136 \\
Learning rate           & 7.66e-4 & 4.05e-4 & 2.50e-5 & 5.34e-5 & 2.88e-4 \\
Batch size              & 16      & 32      & 32      & 32      & 16 \\
Optimizer               & Adam    & Adam    & Adam    & Adam    & Adam \\
Max.\ epochs            & 40      & 40      & 40      & 40      & 40 \\
Early stopping patience & 7       & 7       & 7       & 7       & 7 \\
\hline
\multicolumn{6}{l}{\textbf{(b) Vanilla Transformer}} \\
Lookback                & 60      & 40      & 20      & 20      & 60 \\
\texttt{d\_model}       & 32      & 32      & 64      & 64      & 64 \\
\texttt{nhead}          & 4       & 4       & 2       & 2       & 2 \\
\texttt{num\_layers}$^{\ast}$ & 2   & 2       & 2       & 2       & 2 \\
\texttt{dim\_feedforward} & 128   & 128     & 128     & 128     & 128 \\
Dropout                 & 0.2833  & 0.2963  & 0.2275  & 0.2767  & 0.1000 \\
Learning rate           & 2.02e-4 & 1.74e-4 & 4.80e-4 & 3.68e-4 & 4.86e-4 \\
Batch size              & 16      & 16      & 16      & 16      & 16 \\
Optimizer               & Adam    & Adam    & Adam    & Adam    & Adam \\
Max.\ epochs            & 40      & 40      & 40      & 40      & 40 \\
Early stopping patience & 7       & 7       & 7       & 7       & 7 \\
\hline
\end{tabular}
\par\vspace{4pt}\parbox{\linewidth}{\scriptsize\textit{Notes:} $^{\ast}$Important implementation note: although the tuning log stores the depth parameter as \texttt{num\_encoder\_layers} with values 1, 2, 2, 2, and 1, the holdout evaluation implementation retrieves \texttt{num\_layers} via \texttt{p.get('num\_layers', 2)}. Therefore, the effective number of encoder layers is 2 for all five horizons; the table reports the actual holdout configuration.}
\end{table}

\begin{table}[htbp]
\centering
\caption{Predefined configurations of TCN, PatchTST, and TSMixer.}\label{tab:config_dl_predefined}
\setlength{\tabcolsep}{6pt}
\begin{tabular}{lccc}
\hline
Hyperparameter & TCN & PatchTST & TSMixer \\
\hline
Lookback                & 20           & 20      & 20 \\
Architecture depth      & 3 levels     & 1 layer & 3 blocks \\
Hidden size / channels  & 64           & --      & -- \\
Number of channels      & [64, 64, 64] & --      & -- \\
Kernel size             & 3            & --      & -- \\
Patch length            & --           & 8       & -- \\
Stride                  & --           & 4       & -- \\
\texttt{d\_model}       & --           & 32      & -- \\
\texttt{nhead}          & --           & 2       & -- \\
\texttt{dim\_feedforward} & --         & 64      & -- \\
\texttt{ff\_dim}        & --           & --      & 64 \\
Dropout                 & 0.1          & 0.1     & 0.1 \\
Learning rate           & $10^{-4}$    & $10^{-4}$ & $10^{-4}$ \\
Batch size              & 32           & 32      & 32 \\
Optimizer               & AdamW        & AdamW   & AdamW \\
Weight decay            & $10^{-4}$    & $10^{-4}$ & $10^{-4}$ \\
Max.\ epochs            & 1,000        & 1,000   & 1,000 \\
Early stopping patience & 7            & 7       & 7 \\
\hline
\end{tabular}
\par\vspace{4pt}\parbox{\linewidth}{\scriptsize\textit{Notes:} These configurations are predefined and applied consistently across all five forecasting horizons. They are not selected through the same Optuna-based optimization procedure used for LSTM and Vanilla Transformer.}
\end{table}

\paragraph{Final configuration selection.} For models undergoing fold-based hyperparameter optimization, the configuration selected from Fold 5 is prioritized for the final holdout evaluation, as it corresponds to the longest training window. If no valid Fold 5 configuration is available, the configuration with the lowest recorded \texttt{Best\_Val\_Score} is used as a fallback. Models with predefined specifications or fixed benchmark configurations retain their original settings. All selected configurations are frozen before the final holdout evaluation, and no holdout observations are used for hyperparameter selection or model adjustment. The complete configurations implemented in the holdout experiments are reported in \hyperref[tab:config_econ_ml]{Tables~\ref*{tab:config_econ_ml}}--\hyperref[tab:config_dl_predefined]{\ref*{tab:config_dl_predefined}}.

\paragraph{Computational environment.} The computational experiments are conducted using Python 3.11.9 and PyTorch 2.6.0+cu124, with CUDA acceleration enabled on an NVIDIA GeForce RTX 4060 Laptop GPU with 8 GB of VRAM. The system has approximately 32 GB of RAM and runs Windows 11 Home. The principal software versions include scikit-learn 1.9.0, XGBoost 3.2.0, Optuna 5.0.0, statsmodels 0.15.0, and SciPy 1.17.1.

\paragraph{Multi-seed robustness evaluation.} To assess the sensitivity of deep learning forecasting performance to random initialization, an additional multi-seed evaluation is conducted on the final holdout set (Split 3: 2020--2026). The evaluation includes TSMixer with ten random seeds and LSTM, PatchTST, TCN, and Vanilla Transformer with five random seeds each. The seed values are specified before the experiments and are kept consistent across the evaluated architectures where applicable. The configurations reported in \hyperref[tab:config_lstm_transformer]{Tables~\ref*{tab:config_lstm_transformer}} and~\hyperref[tab:config_dl_predefined]{\ref*{tab:config_dl_predefined}} are kept fixed across all runs, with no additional hyperparameter optimization. For each model and horizon, a new model instance is initialized and trained independently under the corresponding random seed. The baseline seed (42) is included in the multi-seed evaluation, ensuring that the baseline result can be directly compared with the corresponding individual run in the multi-seed analysis. This procedure isolates the sensitivity of forecasting performance to stochastic training initialization while controlling for changes in model configuration and data partitioning.

All runs use the same training, validation, and holdout partitions, preprocessing procedures, and fixed historical-mean (HM) benchmark. The HM benchmark is calculated once from the development sample and remains unchanged across all random-seed runs. For each seed, OOS $R^2$ is calculated independently against this common benchmark. The reported mean OOS $R^2$ is then computed as the arithmetic average of the individual seed-level OOS $R^2$ values, rather than by calculating OOS $R^2$ from the average predictions across seeds:
\[
\frac{1}{S}\sum_{s=1}^{S} R^2_{\mathrm{OOS},s}.
\]
Accordingly, this statistic characterizes the average predictive performance of independently trained model instances rather than the performance of an ensemble formed by averaging their predictions. The resulting OOS $R^2$ values are summarized using the mean, standard deviation, median, minimum, maximum, and proportion of runs achieving positive OOS $R^2$. This supplementary evaluation is intended to assess initialization robustness rather than to select alternative model configurations or identify the best-performing random seed.

\subsection{Feature-Group Ablation Design}\label{subsec:ablation_design}

To examine how different predictor groups affect forecasting performance, a feature-group ablation analysis is conducted using three models: XGBoost, Ridge, and LSTM. These models represent tree-based, linear, and deep learning approaches, respectively. The analysis considers six feature groups: Lagged Returns, Momentum, Historical Volatility, VIX, Intraday, and Trading Volume. The full specification contains 16 predictors, as described in Section~\ref{sec:data}.

Following the dropped-variable importance framework proposed by Hooker, Mentch, and Zhou~\cite{hooker2021}, which evaluates predictive performance changes by comparing a full model with a model retrained after excluding specified predictors, this study extends the leave-one-variable-out approach to feature-group removal. This group-level ablation design is also consistent with established machine learning approaches that allow individual features or groups of features to be excluded from model configurations~\cite{sheikholeslami2021}. For each model and forecasting horizon, the effect of removing a feature group is evaluated by comparing the OOS $R^2$ of the corresponding ablated specification with that of the full model. The difference is defined as:
\begin{equation}\label{eq:delta_oos_r2}
\Delta \mathrm{OOS}\,R^2_g = \mathrm{OOS}\,R^2_{\text{Ablated},g} - \mathrm{OOS}\,R^2_{\text{Full}},
\end{equation}
where $g$ denotes the feature group removed from the model. A negative value indicates that removing the group reduces OOS $R^2$ relative to the full specification, whereas a positive value indicates that the ablated model achieves higher OOS $R^2$.

All ablated specifications retain the hyperparameter configurations selected for their corresponding full models at each forecasting horizon, without additional hyperparameter optimization after feature-group removal. Specifically, the Ridge regularization parameter, the XGBoost tree and learning parameters, and the LSTM architectural and optimization parameters are held fixed at their respective full-model values. Model parameters, including regression coefficients, tree structures, and neural network weights, are re-estimated from scratch using the reduced predictor sets. This fixed-hyperparameter design evaluates the effect of feature-group removal under a consistent model configuration, without allowing additional hyperparameter adaptation to compensate for the excluded predictors.

These differences reflect model- and horizon-specific performance effects of feature-group removal under the evaluated experimental configuration. They should not be interpreted as intrinsic feature importance, independent contributions, or causal effects~\cite{hooker2021}.

\subsection{Sequence Construction and Walk-Forward Implementation}\label{subsec:walk_forward}

The forecasting models are evaluated using the expanding-window walk-forward protocol and chronological data partitions described in \hyperref[subsec:data_splitting]{Section~\ref*{subsec:data_splitting}}. For sequence-based architectures, including LSTM, TCN, Vanilla Transformer, PatchTST, and TSMixer, historical observations are organized into input sequences according to the corresponding lookback length $L$.

When constructing validation and holdout inputs, the final $L-1$ observations from the preceding available data segment are retained as a historical lookback buffer. This allows the first observation in each evaluation segment to be forecast using the required historical context while preserving chronological ordering and avoiding the use of future information. The same sequence-construction procedure is applied consistently across validation folds and the final holdout evaluation.

\subsection{Forecast Evaluation Metrics}\label{subsec:eval_metrics}

Five metrics are used to evaluate forecasting performance: out-of-sample $R^2$, Root Mean Squared Error (RMSE), Mean Absolute Error (MAE), Directional Accuracy (DA), and Excess Directional Accuracy.

The primary metric is the out-of-sample $R^2$, which measures predictive performance relative to the Historical Mean (HM) benchmark \cite{campbell2008}:
\begin{equation}\label{eq:oos_r2}
R^2_{\mathrm{OOS},m,h} = 1 - \frac{\sum_{t=1}^{T_h} (y_t^{(h)} - \hat{y}_{m,t}^{(h)})^2}{\sum_{t=1}^{T_h} (y_t^{(h)} - \hat{y}_{\mathrm{HM},t}^{(h)})^2}.
\end{equation}
Positive values indicate improvements over HM, whereas negative values indicate inferior performance. RMSE and MAE are used as supplementary measures of forecast error magnitude and are calculated according to their standard definitions \cite{hodson2022}.

Directional Accuracy (DA) measures the proportion of forecasts that correctly predict the sign of realized returns, following the directional forecasting evaluation framework discussed by Pesaran and Timmermann \cite{pesaran1992}:
\begin{equation}\label{eq:da}
\text{DA}_{m,h} = \frac{1}{T_h}\sum_{t=1}^{T_h} \mathbf{1}\left[\operatorname{sgn}(\hat{y}_{m,t}^{(h)}) = \operatorname{sgn}(y_t^{(h)})\right].
\end{equation}
To account for differences in the frequency of positive returns across forecast horizons, Excess Directional Accuracy is defined as:
\begin{equation}\label{eq:excess_da}
\text{Excess DA}_{m,h} = \text{DA}_{m,h} - \text{BaseRate}(h),
\end{equation}
where $\text{BaseRate}(h)$ denotes the proportion of positive realized returns at horizon $h$. Since the HM benchmark issues positive forecasts throughout the evaluation sample, its directional accuracy equals this positive-return base rate. Positive Excess DA therefore indicates better directional accuracy than the always-positive HM benchmark, whereas negative values indicate inferior directional performance.

\paragraph{Statistical inference for Excess DA.} To assess whether a candidate model achieves statistically significant directional improvement over the always-positive benchmark, Excess DA is evaluated using a one-sided Newey--West~\cite{newey1987} heteroskedasticity and autocorrelation consistent (HAC) test. For model $m$ and forecast horizon $h$, define the directional performance differential as:
\[
q_{m,t,h}
=
\mathbf{1}\left[
\operatorname{sgn}(\hat{y}_{m,t}^{(h)})
=
\operatorname{sgn}(y_t^{(h)})
\right]
-
\mathbf{1}\left[y_t^{(h)}>0\right].
\]
Its sample mean corresponds to Excess DA:
\[
\overline{q}_{m,h}
=
\frac{1}{T_h}\sum_{t=1}^{T_h}q_{m,t,h}
=
\mathrm{Excess\ DA}_{m,h},
\]
where $T_h$ denotes the number of valid out-of-sample observations at horizon $h$.

The null and one-sided alternative hypotheses are:
\[
H_0:E[q_{m,t,h}]\leq 0,
\qquad
H_1:E[q_{m,t,h}]>0.
\]
The test statistic is:
\[
Z_{m,h}
=
\frac{\overline{q}_{m,h}}
{\sqrt{\widehat{\Omega}_{m,h}/T_h}},
\]
where $\widehat{\Omega}_{m,h}$ is the Newey--West estimate of the long-run variance:
\[
\widehat{\Omega}_{m,h}
=
\widehat{\gamma}_0
+
2\sum_{j=1}^{L_h}
\left(1-\frac{j}{L_h+1}\right)
\widehat{\gamma}_j,
\]
with $\widehat{\gamma}_j$ denoting the sample autocovariance of the directional differential at lag $j$. The HAC truncation lag is matched to the forecast horizon, $L_h=h-1$, to account for the serial dependence induced by overlapping $h$-day cumulative return targets under the corresponding forecasting framework. For the one-day horizon, $L_h=0$. The statistic is evaluated using a one-sided standard normal approximation.

The resulting $p$-values are adjusted using the Holm step-down procedure~\cite{holm1979} across all 60 model--horizon comparisons (12 candidate models and five forecast horizons). Statistical significance is assessed at the 5\% level using adjusted $p$-values.

Validation metrics are calculated as the arithmetic mean across the five walk-forward folds. Final performance is evaluated on the continuous hold-out sample of 1,362 observations. Results are further examined across forecast horizons, volatility regimes, and chronological subperiods.

As a benchmark robustness check, we additionally evaluate forecasting performance against a Zero-Return benchmark, which assumes zero expected cumulative returns over each forecast horizon, such that \(\hat{y}^{(h)}_{\mathrm{Zero},t}=0\). Unlike the primary Historical Mean benchmark, which uses the historical average of the target variable estimated from the training sample, the Zero-Return benchmark provides a fixed zero-return reference. Its corresponding out-of-sample \(R^2\) is calculated as:
\begin{equation}
R^2_{\mathrm{OOS,Zero},m,h} = 1 - \frac{\sum_{t=1}^{T_h}\left(y_t^{(h)}-\hat{y}_{m,t}^{(h)}\right)^2}{\sum_{t=1}^{T_h}\left(y_t^{(h)}\right)^2}.
\end{equation}
The Historical Mean benchmark remains the primary reference for assessing predictive performance, while the Zero-Return benchmark is used as an alternative specification to examine the sensitivity of the results to benchmark selection. Both benchmarks are evaluated over the same hold-out observations for each model and forecast horizon. The resulting comparisons are reported in Section~\ref{subsec:benchmark_robustness}.

\subsection{Statistical Testing Framework}\label{subsec:hypothesis_testing}

The statistical testing framework examines differences in predictive performance across forecast horizons, volatility regimes, and evaluation periods. It comprises three main tests corresponding to H1--H3, supplemented by Diebold--Mariano (DM) tests and binomial sign tests. Newey--West heteroskedasticity and autocorrelation consistent (HAC) estimators are employed to account for serial dependence where applicable.

\paragraph{Normalized loss differentials.} Let $e_{\mathrm{HM},t,h}$ and $e_{m,t,h}$ denote the forecast errors of the historical mean (HM) benchmark and model $m$, respectively. The normalized squared error differential is defined as:
\begin{equation}\label{eq:norm_loss_diff}
\widetilde d_{m,t,h}
=
\frac{e^2_{\mathrm{HM},t,h}-e^2_{m,t,h}}
{\mathrm{MSE}_{\mathrm{HM},h}},
\qquad
\mathrm{MSE}_{\mathrm{HM},h}
=
\frac{1}{T_h}\sum_{t=1}^{T_h}e^2_{\mathrm{HM},t,h}.
\end{equation}
A positive value indicates that model $m$ achieves a lower squared forecast error than the HM benchmark, whereas a negative value indicates inferior predictive performance relative to HM.

\paragraph{H1 -- Multi-horizon predictive performance.} H1 tests whether mean normalized predictive performance differs across the five forecast horizons $h\in\{1,5,20,60,120\}$. For each model, the horizon-specific differentials are arranged into a synchronized vector:
\[
X_{m,t}
=
\left(
\widetilde d_{m,t,1},
\widetilde d_{m,t,5},
\widetilde d_{m,t,20},
\widetilde d_{m,t,60},
\widetilde d_{m,t,120}
\right)^\top
\in\mathbb{R}^5,
\]
retaining only dates with valid observations at all five horizons. In the holdout evaluation, all five forecast horizons share the same 1,362 valid forecast origins. Therefore, the synchronized sample coincides with the full holdout sample, ensuring that the horizon-specific HM normalization terms in Equation~(\ref{eq:norm_loss_diff}) are computed over the same observations used to construct the synchronized differential vectors. Let $\theta_m=E[X_{m,t}]$. The null hypothesis is:
\begin{equation}
H_0:R\theta_m=0,
\end{equation}
where
\begin{equation}
R=
\begin{bmatrix}
1&-1&0&0&0\\
0&1&-1&0&0\\
0&0&1&-1&0\\
0&0&0&1&-1
\end{bmatrix}.
\end{equation}
This restriction tests whether the five horizon-specific mean normalized loss differentials are equal. The Wald statistic is:
\begin{equation}
W_{\mathrm{H1},m}
=
T(R\bar X_m)^\top
\left(R\widehat\Omega_mR^\top\right)^{-1}
(R\bar X_m),
\end{equation}
where $\bar X_m$ is the sample mean vector, $T$ is the synchronized sample size, and $\widehat\Omega_m$ is the Newey--West long-run covariance matrix estimated jointly from the five-dimensional differential vector. The baseline HAC lag truncation is $L=20$, applied consistently across models. This specification accounts for serial dependence in the loss differential series arising from overlapping multi-period forecast targets, following the overlapping-data inference framework established by Hansen and Hodrick~\cite{hansen1980}. In particular, overlapping $h$-period forecast targets can induce dependence extending up to $h-1$ lags under the corresponding forecasting framework. Hodrick~\cite{hodrick1992} further documents potential finite-sample distortions in conventional inference procedures at longer forecast horizons. Under the null, the statistic asymptotically follows a chi-squared distribution. The degrees of freedom are determined by the effective rank of the restriction covariance matrix: four for eleven models and two for Elastic Net, whose differential series are degenerate at three horizons.

Robustness is assessed using alternative HAC bandwidths $L\in\{5,60,119\}$, in addition to the baseline specification. A paired synchronized moving-block bootstrap is also conducted using five block lengths, $b\in\{20,40,60,80,120\}$, with 4,999 replications for each block length. The synchronized resampling scheme preserves the cross-horizon pairing of observations within each sampled block. The same resampled time indices are applied across all five horizon-specific differential series, preserving both their temporal dependence within blocks and their contemporaneous cross-horizon alignment. Bootstrap samples are recentered under the null hypothesis, and studentized Wald statistics are recalculated using the baseline HAC bandwidth $L=20$. Bootstrap $p$-values are subsequently adjusted for multiple testing using the Holm procedure separately for each block length. The bootstrap results are interpreted as a complementary finite-sample robustness assessment rather than as an equivalent replacement for the asymptotic HAC inference.

\paragraph{H2 -- Regime-Conditional Predictive Ability and Horizon--Regime Interaction.} H2 examines whether the predictive performance of candidate models relative to the Historical Mean (HM) benchmark exhibits regime-specific conditional behavior and whether the relative performance contrasts across volatility regimes vary with the forecast horizon. Using the volatility regime indicators defined in \hyperref[subsec:regime_classification]{Section~\ref*{subsec:regime_classification}}, the analysis considers Low, Normal, and High volatility regimes, with Normal volatility serving as the reference category.

The analysis comprises two complementary components. First, the Giacomini--White conditional predictive ability (GW-CPA) test examines whether the model's normalized loss differential relative to HM has zero conditional moments with respect to the volatility-regime instruments at each forecast horizon. Second, the horizon--regime interaction test examines whether the relative predictive performance contrasts between the Low and Normal regimes and between the High and Normal regimes vary across forecast horizons. Together, these tests assess conditional predictive performance relative to HM and the horizon dependence of regime-related performance contrasts, without directly establishing unconditional differences in predictive performance between the three volatility regimes.

\textit{Single-horizon conditional predictive ability (CPA).} At each forecast horizon $h$, the Giacomini--White~\cite{giacomini2006} conditional predictive ability (GW-CPA) test constructs the three-dimensional moment vector:
\begin{equation}
m_{m,t,h}=z_t\widetilde d_{m,t,h},
\end{equation}
using the regime instrument vector:
\begin{equation}
z_t=
\begin{bmatrix}
1\\
I(\mathrm{Low}_t)\\
I(\mathrm{High}_t)
\end{bmatrix},
\end{equation}
with Normal volatility serving as the reference regime absorbed into the intercept.

The joint null hypothesis is:
\begin{equation}
H_0: E[z_t\widetilde d_{m,t,h}]=0.
\end{equation}
This restriction is equivalent to requiring the expected normalized loss differential to equal zero conditionally within each of the three volatility regimes, provided that each regime has a positive probability of occurrence. Thus, rejection indicates that the model exhibits statistically significant conditional predictive ability relative to HM in at least one volatility regime. However, the test does not directly establish that predictive performance differs between Low, Normal, and High volatility regimes.

Letting
\begin{equation}
\bar m_{m,h} = \frac{1}{T_h}\sum_{t=1}^{T_h}z_t\widetilde d_{m,t,h},
\end{equation}
the Wald statistic is:
\begin{equation}
W_{\mathrm{CPA},m,h}
=
\bar m_{m,h}^{\top}
\left(
\frac{\widehat\Omega_{m,h}}{T_h}
\right)^{-1}
\bar m_{m,h}
\overset{a}{\sim}\chi^2(3),
\end{equation}
where $\widehat\Omega_{m,h}$ is the Newey--West HAC long-run covariance matrix of the moment series $z_t\widetilde d_{m,t,h}$. The HAC estimator uses a Bartlett kernel with a baseline lag truncation of $L=20$. Sensitivity analyses are conducted using $L\in\{5,60,119\}$. The joint null is evaluated separately for each model and horizon, with Holm adjustment applied across the twelve models within each horizon. For degenerate cases in which $\widetilde d_{m,t,h}\equiv 0$ identically (Elastic Net at $h\in\{20,60,120\}$), the corresponding moment vector is identically zero. These cases are treated as having zero effective degrees of freedom ($\mathrm{df}=0$).

\textit{Horizon--regime interaction.} To examine whether the relative predictive performance contrasts between volatility regimes vary across forecast horizons, observations are stacked across all five horizons $h\in H=\{1,5,20,60,120\}$ into the following horizon-specific projection:
\begin{equation}
\widetilde d_{m,t,h}
=
\alpha_h+\beta_{h,L}I(\mathrm{Low}_t)
+\beta_{h,H}I(\mathrm{High}_t)+u_{m,t,h},
\end{equation}
where Normal volatility serves as the reference regime. The coefficients $\beta_{h,L}$ and $\beta_{h,H}$ represent the Low--Normal and High--Normal differences in mean normalized loss differentials at horizon $h$, respectively.

The null hypothesis examines whether these regime contrasts remain unchanged relative to the 1-day baseline horizon:
\begin{equation}
H_0:
\begin{cases}
\beta_{h,L}-\beta_{1,L}=0,\\
\beta_{h,H}-\beta_{1,H}=0,
\end{cases}
\quad
\forall h\in\{5,20,60,120\}.
\end{equation}
The Wald statistic jointly tests these eight restrictions using the stacked five-horizon GMM/OLS covariance estimator:
\begin{equation}
W_{\mathrm{int}}
=
(R_{\mathrm{int}}\widehat\beta_{\mathrm{int}})^\top
\left(
R_{\mathrm{int}}\widehat V_{\mathrm{panel}}
R_{\mathrm{int}}^\top
\right)^{-1}
(R_{\mathrm{int}}\widehat\beta_{\mathrm{int}})
\overset{a}{\sim}\chi^2(8),
\end{equation}
where $\widehat V_{\mathrm{panel}}$ is estimated from the stacked moment conditions using a Newey--West HAC estimator with a Bartlett kernel and baseline lag truncation $L=20$. The covariance estimator accounts for contemporaneous cross-horizon dependence through the stacked moment matrix. Robustness is assessed using alternative bandwidths $L\in\{5,60,119\}$. Holm adjustment is applied across the twelve model-specific interaction tests. The test therefore determines whether the Low--Normal and High--Normal performance contrasts change across horizons relative to the 1-day baseline. It does not test whether the three volatility regimes have equal predictive performance at any individual horizon. For Elastic Net, the effective degrees of freedom reduce to four because of identical benchmark predictions at 20, 60, and 120 days.

\paragraph{H3 -- Temporal stability of predictive performance.} H3 evaluates whether predictive gains are stable over the evaluation sample or concentrated in particular calendar subperiods.

\textit{Giacomini--Rossi-type fluctuation test with a moving-block bootstrap approximation.} Motivated by the fluctuation-based approach of Giacomini and Rossi~\cite{giacomini2010}, we implement a modified trailing-window fluctuation procedure and use a moving-block bootstrap to obtain empirical critical probabilities under the null hypothesis. Unlike the original Giacomini and Rossi~\cite{giacomini2010} procedure, which derives its asymptotic reference distribution from a functional central limit theorem, the present implementation uses a fixed-length moving-block bootstrap approximation. For each model and horizon, a local standardized fluctuation statistic is computed over trailing windows of length $m_w = 252$ trading days:
\begin{equation}
t_{m,\tau,h}
=
\frac{\sqrt{m_w}}{\widehat\sigma_{\mathrm{HAC},m,h}}
\left(
\frac{1}{m_w}
\sum_{s=\tau-m_w+1}^{\tau}
\widetilde d_{m,s,h}
\right),
\qquad
\tau=m_w,\ldots,T_h,
\end{equation}
where $m_w = 252$ and $\widehat\sigma_{\mathrm{HAC},m,h}$ is the long-run standard deviation of $\widetilde d_{m,t,h}$ estimated over the full holdout sample using a Bartlett kernel with baseline lag truncation $L = 20$ separately for each model--horizon combination. This standardization factor is held fixed across all rolling windows. The trailing-window construction yields $T_h-m_w+1$ rolling statistics for each model--horizon combination. For $T_h=1,362$, this corresponds to 1,111 windows. The summary statistic is the two-sided supremum:
\begin{equation}
S_{m,h}
=
\max_{\tau=m_w,\ldots,T_h}
|t_{m,\tau,h}|.
\end{equation}
Empirical $p$-values are obtained using a fixed-length moving-block bootstrap with a baseline of $B=500$ replications and block length $b=60$. Before resampling, the loss-differential series is demeaned:
\begin{equation}
\widetilde d^*_{m,t,h} = \widetilde d_{m,t,h} - \bar{\widetilde d}_{m,h}.
\end{equation}
This recentering imposes a zero sample-mean restriction on the resampling series. For each bootstrap replication, the trailing-window fluctuation statistics are recomputed using the same full-sample HAC scaling factor as in the observed statistic. The bootstrap $p$-value is calculated as:
\begin{equation}
\widehat p = \frac{1}{B} \sum_{r=1}^{B} I(S_r^* \geq S_{\mathrm{obs}}), \qquad B=500.
\end{equation}
When none of the bootstrap replications exceeds the observed statistic, the empirical $p$-value is reported as $p<1/B$ rather than zero. Holm adjustment is applied across the twelve models separately within each forecast horizon.

Bootstrap sensitivity analyses are additionally conducted to examine the effects of replication count and block length on the fluctuation-test inference. Specifically, the number of replications is increased to $B\in\{1,999,4,999\}$ under the baseline block length $b=60$, while an alternative block length of $b=120$ is examined using $B\in\{500,4,999\}$. The same observed fluctuation statistic, null-centering procedure, trailing-window construction, and HAC standardization are maintained across these configurations. Only the bootstrap replication count and block length are varied. Holm adjustment is recalculated separately across the twelve models within each forecast horizon for each bootstrap configuration. These analyses assess Monte Carlo sensitivity and the dependence of the empirical inference on the selected resampling block length.

To assess sensitivity to the HAC bandwidth, the fluctuation statistic is recomputed using $L\in\{5,20,60,119\}$, while keeping the rolling-window length, block length, bootstrap replication count, and random seed fixed. Because the same bandwidth-specific HAC scaling factor is applied to both the observed and bootstrap statistics, the bootstrap exceedance probability is algebraically invariant to the bandwidth choice. Accordingly, this exercise evaluates the sensitivity of the standardized fluctuation statistic but does not constitute an independent robustness check of the bootstrap $p$-value.

\textit{Subperiod stability test.} To complement the continuous fluctuation test, calendar block stability is evaluated using a block-level HAC Wald test. Following the calendar-based subperiod definitions in \hyperref[subsec:data_splitting]{Section~\ref*{subsec:data_splitting}}, the analysis examines differences in mean normalized predictive performance across the predefined validation and holdout blocks. For the holdout sample, the November--December 2020 observations are excluded from the block-level analysis, while the remaining three calendar blocks are evaluated. The same three-block framework is applied to the validation sample. Let $\bar{\widetilde d}_k$ denote the mean normalized squared error differential in block $k$, scaled by the full-sample HM MSE rather than within-block OOS $R^2$. The null hypothesis of equal block-mean predictive performance is:
\begin{equation}
H_0:
\bar{\widetilde d}_1
=
\bar{\widetilde d}_2
=
\bar{\widetilde d}_3.
\end{equation}
The Wald test statistic is:
\begin{equation}
W_{\mathrm{sub}}
=
(R\widehat\mu)^\top
(R\widehat\Sigma_{\mathrm{HAC}}R^\top)^{-1}
(R\widehat\mu)
\overset{a}{\sim}\chi^2(2),
\end{equation}
where $\widehat\mu=(\bar{\widetilde d}_1,\bar{\widetilde d}_2,\bar{\widetilde d}_3)^\top$, $R$ imposes the two linear equality restrictions, and $\widehat\Sigma_{\mathrm{HAC}}$ is estimated using a Newey--West HAC estimator with a Bartlett kernel and baseline lag truncation $L=20$. The sensitivity of the subperiod Wald test to alternative bandwidths $L\in\{5,60,119\}$ is also examined. The test is conducted separately for the holdout and validation evaluation samples, with three calendar blocks in each sample. When a model's forecasts collapse identically to the HM benchmark across blocks (yielding zero differentials and $\mathrm{df}=0$, as observed for Elastic Net at longer horizons), the resulting statistic is identically zero.

\paragraph{Multiple-testing adjustment.} Multiple-testing adjustments are performed using the Holm step-down procedure. For the GW-CPA and fluctuation tests, Holm adjustment is applied separately within each forecast-horizon family across the twelve models. For the horizon--regime interaction tests, adjustment is performed across the twelve model-specific tests. Both raw and adjusted $p$-values are reported, and statistical significance is assessed at the 5\% level based on adjusted $p$-values.

\paragraph{Supplementary predictive accuracy tests.} In addition to the hypothesis-specific tests, two conventional statistical tests are conducted on the full out-of-sample evaluation sample to provide complementary assessments of predictive accuracy.

The Diebold--Mariano (DM) test evaluates whether a candidate model achieves a significantly lower mean squared forecast error than the Historical Mean (HM) benchmark~\cite{diebold1995}. For model $m$ and forecast horizon $h$, the loss differential is defined as:
\[
d_{m,t,h}
=
e^2_{\mathrm{HM},t,h}-e^2_{m,t,h}.
\]
A positive value indicates a forecasting improvement over HM. The test statistic is calculated as:
\[
\mathrm{DM}_{m,h}
=
\frac{\bar d_{m,h}}
{\sqrt{\widehat{\operatorname{Var}}(\bar d_{m,h})}},
\]
with the long-run variance estimated using a rectangular spectral window truncated at lag $h-1$. The Harvey--Leybourne--Newbold (HLN) small-sample correction is then applied~\cite{buturac2022,diebold1995,harvey1997}:
\begin{equation}\label{eq:dm_stat}
\mathrm{DM}^{*}_{m,h}
=
\mathrm{DM}_{m,h}
\sqrt{
\frac{T+1-2h+h(h-1)/T}{T}
}.
\end{equation}
The corrected statistic is evaluated against a Student-$t$ distribution with $T-1$ degrees of freedom~\cite{buturac2022,harvey1997}. A one-sided alternative is adopted to test whether the candidate model significantly outperforms HM.

\paragraph{Supplementary directional accuracy test.} As a conventional reference, the exact one-sided Binomial Test is additionally reported to examine whether directional accuracy exceeds the 50\% random-guessing benchmark:
\[
H_0:p\leq 0.5,
\qquad
H_1:p>0.5.
\]
Let $K$ denote the number of correct directional predictions among $n$ non-zero realized returns. Under the independent Bernoulli assumption, the corresponding $p$-value is:
\[
p_{\mathrm{bin}}
=
P(K_{\mathrm{bin}}\geq K)
=
\sum_{j=K}^{n}
\binom{n}{j}(0.5)^n.
\]
However, the independence assumption underlying the conventional Binomial Test is generally violated when cumulative $h$-day return targets overlap. Consequently, successive directional outcomes may exhibit serial dependence, making the nominal Binomial $p$-values potentially unreliable. Therefore, this test is treated exclusively as a supplementary descriptive reference and is not used as the primary statistical evidence of directional forecasting skill.

The primary assessment of directional improvement is instead based on Excess DA and horizon-matched Newey--West HAC inference, as defined in \hyperref[subsec:eval_metrics]{Section~\ref*{subsec:eval_metrics}}. Unlike the Binomial Test, Excess DA evaluates directional improvement relative to the realized positive-return frequency, which corresponds to the directional accuracy of the always-positive HM benchmark.

Both the DM and supplementary Binomial tests are conducted on the full out-of-sample evaluation sample for all twelve candidate models across the five forecast horizons. For the DM tests against the Historical Mean (HM) benchmark, raw $p$-values are calculated for each individual model--horizon comparison and Holm's step-down procedure is applied at two levels: separately within each forecast horizon across the twelve candidate models, and jointly across all 60 model--horizon comparisons. Both raw and Holm-adjusted $p$-values are reported, with statistical significance assessed at the 5\% level based on the adjusted $p$-values. For supplementary DM comparisons against alternative benchmarks, raw $p$-values are reported as nominal individual significance levels unless otherwise specified. The supplementary Binomial test $p$-values are also interpreted as nominal individual significance levels without multiple-testing adjustment. In contrast, Excess DA inference uses Holm-adjusted $p$-values across the 60 model--horizon comparisons.

\section{Empirical Results and Discussion}\label{sec:results}

\subsection{Out-of-Sample Forecast Performance}\label{subsec:oos_performance}

This section evaluates the out-of-sample forecasting performance of the twelve candidate models relative to the Historical Mean (HM) benchmark across five forecast horizons. Forecast performance is assessed using OOS $R^2$, root mean squared error (RMSE), mean absolute error (MAE), and directional accuracy (DA).

\hyperref[tab:oos_r2_main]{Table~\ref*{tab:oos_r2_main}} reports OOS $R^2$ relative to HM. Conventional machine-learning models generally exhibit more stable performance at short and intermediate horizons, whereas deep-learning results are more heterogeneous. TSMixer records the strongest OOS $R^2$ at 20D (+3.80\%), 60D (+4.98\%), and 120D (+17.53\%). At the 120D horizon, only TSMixer and ARIMA achieve positive OOS $R^2$, at +17.53\% and +1.95\%, respectively. However, these point estimates should not be interpreted as evidence of statistically significant improvements. TSMixer's DM test at the 120D horizon yields a raw (unadjusted) $p$-value of 0.0485, indicating nominal significance at the 5\% level. However, this result does not remain statistically significant after multiple-testing adjustment. The corresponding Holm-adjusted $p$-value is 0.5817 when correction is applied within the 12-model family at the 120D horizon, and 1.0000 when correction is applied jointly across all 60 model--horizon comparisons. Accordingly, the nominal significance of this individual comparison does not provide sufficient evidence of statistically reliable predictive improvement over HM after accounting for multiple testing. Overall, no candidate model demonstrates statistically significant forecast improvement over HM after multiple-testing adjustment.

Although TSMixer records positive OOS $R^2$ values at 20D, 60D, and 120D, with its largest gain occurring at 120D, none of these improvements is statistically significant after multiple-testing adjustment. Therefore, the favorable point estimates should be interpreted cautiously rather than as evidence of statistically reliable superiority over HM.

\begin{table}[htbp]
\centering
\caption{Out-of-sample $R^2$ by model and forecast horizon.}\label{tab:oos_r2_main}
\setlength{\tabcolsep}{4.5pt}
\begin{tabular}{lrrrrr}
\hline
Model & 1D & 5D & 20D & 60D & 120D \\
\hline
ARIMA & $-$0.49\% & $-$8.25\% & $-$12.70\% & $-$10.26\% & +1.95\% \\
ARIMAX & +1.18\% & $-$5.90\% & $-$4.61\% & +2.52\% & $-$0.35\% \\
Ridge & +1.50\% & +1.79\% & +1.66\% & +0.91\% & $-$4.95\% \\
Elastic Net & +0.13\% & +0.61\% & 0.00\% & 0.00\% & 0.00\% \\
SVR & +1.98\% & +1.48\% & +3.05\% & +4.64\% & $-$19.10\% \\
Random Forest & +1.77\% & $-$1.30\% & $-$0.39\% & +0.33\% & $-$6.42\% \\
XGBoost & +1.48\% & +0.11\% & +1.37\% & +0.13\% & $-$0.14\% \\
LSTM & $-$0.65\% & $-$0.50\% & +1.48\% & $-$0.96\% & $-$3.34\% \\
Transformer & $-$15.28\% & $-$39.72\% & $-$11.64\% & $-$21.28\% & $-$23.02\% \\
TCN & $-$14.16\% & $-$64.91\% & $-$123.43\% & $-$87.23\% & $-$142.55\% \\
PatchTST & $-$4.25\% & $-$6.85\% & $-$12.17\% & $-$21.86\% & $-$12.22\% \\
TSMixer & $-$2.72\% & $-$5.47\% & +3.80\% & +4.98\% & +17.53\% \\
\hline
\end{tabular}
\par\vspace{4pt}\parbox{0.85\linewidth}{\footnotesize\textit{Note:} OOS $R^2$ is calculated relative to the Historical Mean (HM) benchmark. See \hyperref[subsec:eval_metrics]{Section~\ref*{subsec:eval_metrics}} for its definition and interpretation.}
\end{table}

\hyperref[tab:rmse_main]{Table~\ref*{tab:rmse_main}} and \hyperref[tab:mae_main]{Table~\ref*{tab:mae_main}} present RMSE and MAE, respectively. The results indicate that the model with the lowest error varies across forecast horizons and evaluation metrics. SVR achieves the lowest RMSE at 1D (0.0104), while TSMixer records the lowest RMSE at 20D (0.0409), 60D (0.0639), and 120D (0.0840). For MAE, SVR performs best at 5D (0.0164) and 60D (0.0489), whereas TSMixer achieves the lowest values at 20D (0.0308) and 120D (0.0663). At 120D, TSMixer improves upon HM in both RMSE (0.0840 versus 0.0925) and MAE (0.0663 versus 0.0757). In contrast, TCN exhibits substantially larger errors at longer horizons, with its 120D RMSE reaching 0.1441. These results highlight differences in point-forecast performance across model architectures, although lower error values alone do not establish statistically significant gains over the benchmark.

\begin{table}[htbp]
\centering
\caption{RMSE by model and forecast horizon.}\label{tab:rmse_main}
\setlength{\tabcolsep}{4.5pt}
\begin{tabular}{lrrrrr}
\hline
Model & 1D & 5D & 20D & 60D & 120D \\
\hline
HM & 0.0106 & 0.0223 & 0.0417 & 0.0655 & 0.0925 \\
ARIMA & 0.0106 & 0.0232 & 0.0442 & 0.0688 & 0.0916 \\
ARIMAX & 0.0105 & 0.0230 & 0.0426 & 0.0647 & 0.0927 \\
Ridge & 0.0105 & 0.0221 & 0.0413 & 0.0652 & 0.0948 \\
Elastic Net & 0.0105 & 0.0223 & 0.0417 & 0.0655 & 0.0925 \\
SVR & 0.0104 & 0.0222 & 0.0410 & 0.0640 & 0.1010 \\
Random Forest & 0.0105 & 0.0225 & 0.0418 & 0.0654 & 0.0954 \\
XGBoost & 0.0105 & 0.0223 & 0.0414 & 0.0655 & 0.0926 \\
LSTM & 0.0106 & 0.0224 & 0.0414 & 0.0658 & 0.0941 \\
Transformer & 0.0113 & 0.0264 & 0.0440 & 0.0722 & 0.1026 \\
TCN & 0.0113 & 0.0287 & 0.0623 & 0.0897 & 0.1441 \\
PatchTST & 0.0108 & 0.0231 & 0.0441 & 0.0723 & 0.0980 \\
TSMixer & 0.0107 & 0.0229 & 0.0409 & 0.0639 & 0.0840 \\
\hline
\end{tabular}
\par\vspace{4pt}\parbox{0.85\linewidth}{\footnotesize\textit{Note:} RMSE is measured in cumulative log-return units. Lower values indicate better point-forecast accuracy.}
\end{table}

\begin{table}[htbp]
\centering
\caption{MAE by model and forecast horizon.}\label{tab:mae_main}
\setlength{\tabcolsep}{4.5pt}
\begin{tabular}{lrrrrr}
\hline
Model & 1D & 5D & 20D & 60D & 120D \\
\hline
HM & 0.0075 & 0.0166 & 0.0327 & 0.0527 & 0.0757 \\
ARIMA & 0.0076 & 0.0171 & 0.0338 & 0.0532 & 0.0690 \\
ARIMAX & 0.0075 & 0.0172 & 0.0338 & 0.0512 & 0.0732 \\
Ridge & 0.0075 & 0.0165 & 0.0327 & 0.0514 & 0.0765 \\
Elastic Net & 0.0075 & 0.0165 & 0.0327 & 0.0527 & 0.0757 \\
SVR & 0.0075 & 0.0164 & 0.0316 & 0.0489 & 0.0867 \\
Random Forest & 0.0075 & 0.0167 & 0.0328 & 0.0534 & 0.0784 \\
XGBoost & 0.0075 & 0.0165 & 0.0327 & 0.0526 & 0.0759 \\
LSTM & 0.0076 & 0.0166 & 0.0322 & 0.0531 & 0.0760 \\
Transformer & 0.0084 & 0.0196 & 0.0364 & 0.0559 & 0.0849 \\
TCN & 0.0084 & 0.0209 & 0.0440 & 0.0649 & 0.0998 \\
PatchTST & 0.0078 & 0.0168 & 0.0344 & 0.0553 & 0.0759 \\
TSMixer & 0.0077 & 0.0169 & 0.0308 & 0.0497 & 0.0663 \\
\hline
\end{tabular}
\par\vspace{4pt}\parbox{0.85\linewidth}{\footnotesize\textit{Note:} MAE is measured in cumulative log-return units. Lower values indicate better point-forecast accuracy.}
\end{table}

Directional performance is evaluated using DA and Excess DA, defined as the difference between a candidate model's DA and that of HM, in percentage points (pp). \hyperref[tab:da_main]{Table~\ref*{tab:da_main}} reports both measures. Because HM issues positive forecasts throughout all horizons, its DA equals the corresponding positive-return base rate, increasing from 53.74\% at 1D to 78.05\% at 120D. Consequently, the increase in raw DA across longer horizons should not, by itself, be interpreted as an improvement in directional forecasting ability, because it partly reflects the increasing frequency of positive cumulative returns rather than improved model-specific directional skill. Only TSMixer at 20D (+1.32 pp) and ARIMAX at 120D (+0.81 pp) numerically outperform HM. However, neither improvement is statistically significant under horizon-matched Newey--West HAC inference ($p=0.2758$ and $p=0.2035$, respectively). Furthermore, none of the 60 model--horizon comparisons demonstrates a statistically significant directional improvement after Holm correction. Nine comparisons produce exactly zero Excess DA because the corresponding models generate positive forecasts for all 1,362 observations, matching HM's directional predictions.

\begin{table}[htbp]
\centering
\caption{Directional accuracy and excess directional accuracy.}\label{tab:da_main}
{\small
\setlength{\tabcolsep}{2.2pt}
\begin{tabular}{lrrrrr@{\hspace{4pt}}c@{\hspace{4pt}}rrrrr}
\hline
 & \multicolumn{5}{c}{\textit{Panel A. DA (\%)}} & & \multicolumn{5}{c}{\textit{Panel B. Excess DA vs.\ HM (pp)}} \\
\cline{2-6}\cline{8-12}
Model & 1D & 5D & 20D & 60D & 120D & & 1D & 5D & 20D & 60D & 120D \\
\hline
HM & 53.74 & 59.40 & 66.89 & 74.96 & 78.05 & & -- & -- & -- & -- & -- \\
ARIMA & 50.37 & 53.01 & 59.10 & 68.21 & 76.87 & & $-$3.38 & $-$6.39 & $-$7.78 & $-$6.75 & $-$1.17 \\
ARIMAX & 52.94 & 50.00 & 53.45 & 72.10 & 78.85 & & $-$0.81 & $-$9.40 & $-$13.44 & $-$2.86 & \textbf{+0.81} \\
Ridge & 51.10 & 53.74 & 63.29 & 72.47 & 77.09 & & $-$2.64 & $-$5.65 & $-$3.60 & $-$2.50 & $-$0.95 \\
Elastic Net & 51.25 & 58.37 & 66.89 & 74.96 & 78.05 & & $-$2.50 & $-$1.03 & 0.00 & 0.00 & 0.00 \\
SVR & 52.72 & 58.08 & 66.81 & 74.96 & 65.12 & & $-$1.03 & $-$1.32 & $-$0.07 & 0.00 & $-$12.92 \\
Random Forest & 53.30 & 58.81 & 65.42 & 73.20 & 74.52 & & $-$0.44 & $-$0.59 & $-$1.47 & $-$1.76 & $-$3.52 \\
XGBoost & 53.16 & 59.25 & 66.52 & 74.96 & 78.05 & & $-$0.59 & $-$0.15 & $-$0.37 & 0.00 & 0.00 \\
LSTM & 50.07 & 51.69 & 66.89 & 74.96 & 78.05 & & $-$3.67 & $-$7.71 & 0.00 & 0.00 & 0.00 \\
Transformer & 48.53 & 51.98 & 42.14 & 62.56 & 72.03 & & $-$5.21 & $-$7.42 & $-$24.74 & $-$12.41 & $-$6.02 \\
TCN & 46.26 & 50.29 & 52.57 & 70.34 & 73.64 & & $-$7.49 & $-$9.10 & $-$14.32 & $-$4.63 & $-$4.41 \\
PatchTST & 50.88 & 56.98 & 50.88 & 74.30 & 74.89 & & $-$2.86 & $-$2.42 & $-$16.01 & $-$0.66 & $-$3.16 \\
TSMixer & 49.12 & 53.67 & 68.21 & 73.13 & 77.39 & & $-$4.63 & $-$5.73 & \textbf{+1.32} & $-$1.84 & $-$0.66 \\
\hline
\end{tabular}}
\par\vspace{4pt}\parbox{\linewidth}{\footnotesize\textit{Note:} Excess DA is calculated as candidate-model DA minus HM DA in percentage points (pp). Positive values indicate a numerical improvement over HM; they do not imply statistical significance. Values are rounded to two decimal places. The unrounded differences for TSMixer 20D and ARIMAX 120D are +1.3216 pp and +0.8076 pp, respectively.}
\end{table}

The comparison between magnitude-based and directional metrics reveals an important distinction. Although TSMixer achieves the highest OOS $R^2$ at 120D and lower RMSE and MAE than HM, its DA remains below the benchmark by 0.66 percentage points (pp). Conversely, ARIMAX achieves a small positive directional difference at 120D (+0.81 pp), despite recording a slightly higher RMSE than HM and a negative OOS $R^2$ ($-$0.35\%). These contrasting results indicate that improvements in forecast magnitude accuracy do not necessarily translate into better directional predictions. Overall, the evaluation reveals metric-dependent differences in model performance. Although several models achieve favorable point estimates under particular horizons and evaluation metrics, the multiple-testing-adjusted results provide no statistically significant evidence of improved forecast accuracy or directional performance over HM.

Importantly, statistical forecast improvements and economic value are distinct dimensions of predictive performance. Campbell and Thompson~\cite{campbell2008} show that relatively small improvements in out-of-sample forecasting can have economic implications when incorporated into investment decisions. Similarly, Gu et al.~\cite{gu2020} demonstrate that predictive information can translate into portfolio-level benefits, although such benefits must be evaluated alongside implementation costs and portfolio turnover. However, the present study evaluates statistical forecasting accuracy rather than portfolio returns or investor utility. Therefore, its OOS $R^2$ improvements should not be interpreted directly as evidence of economic profitability.

\subsection{Horizon-Dependent Predictability (RQ1)}

RQ1 examines whether equity return predictability varies across forecast horizons and whether the observed patterns are consistent across model classes. Rather than repeating the aggregate performance comparisons reported in \hyperref[subsec:oos_performance]{Section~\ref*{subsec:oos_performance}}, this subsection focuses on the shape of horizon-specific performance profiles and formally evaluates cross-horizon differences through H1.

\noindent\textbf{Descriptive Evidence of Horizon-Dependent Predictability.} The cross-horizon performance profiles reveal substantial heterogeneity rather than a uniform decline in predictive performance as the forecast horizon increases.

At short horizons (1D--5D), conventional machine-learning models exhibit relatively competitive performance, although their rankings are not consistent across the two horizons. This variation suggests that short-term predictive performance depends on the interaction between model specification and forecast horizon, rather than reflecting a systematic advantage of a particular model family.

At intermediate horizons (20D--60D), the relative performance profiles become more differentiated across architectures. Certain conventional machine-learning models retain comparatively stable performance, while TSMixer demonstrates stronger relative performance within the deep-learning group. This pattern suggests that predictive performance at intermediate horizons is model-dependent and cannot be characterized by a uniform advantage of either conventional machine learning or deep learning.

At the longest horizon (120D), TSMixer exhibits a pronounced improvement in relative point-forecast performance compared with its results at shorter horizons. This finding is particularly noteworthy because it contrasts with the assumption that predictive performance necessarily deteriorates as the forecast horizon extends. However, the observed pattern is specific to the evaluated TSMixer implementation and does not establish that longer horizons are generally more predictable. More broadly, the cross-horizon profiles indicate that model performance can vary non-monotonically, with certain architectures exhibiting relative advantages at particular horizons rather than maintaining consistent superiority across the entire forecasting range.

The observed non-monotonic pattern is broadly consistent with previous evidence that predictive performance varies across forecast horizons~\cite{quaedvlieg2021,xu2020}. However, the present findings concern aggregate S\&P~500 time-series forecasting and should be distinguished from evidence based on cross-sectional individual-stock prediction~\cite{nechvatalova2024}.

The cross-horizon results also caution against assuming that greater model complexity necessarily leads to stronger predictive performance. Transformer, TCN, and PatchTST do not consistently outperform simpler alternatives across the examined horizons. This observation suggests that architectural complexity alone is insufficient to ensure reliable financial forecasting performance, consistent with concerns raised in previous studies~\cite{gu2020,welch2008}.

\medskip
\noindent\textbf{Statistical Evidence for H1.} To determine whether the observed differences across horizons are statistically distinguishable, H1 tests the joint null hypothesis that a model's mean normalized loss differentials are equal across all five forecast horizons. Because the normalization denominator is the horizon-specific HM mean squared error, the sample mean of each normalized differential is equivalent to the corresponding OOS $R^2$. Thus, H1 formally evaluates whether the model's relative predictive performance differs across horizons. It does not test whether the model significantly outperforms HM at any individual horizon, nor does it impose a monotonic-decay restriction.

Under the baseline HAC specification with bandwidth $L=20$, ARIMA is the only model for which the equality restriction is rejected after Holm correction ($W=15.7438$, $\mathrm{df}=4$, adjusted $p=0.0406$). This rejection reflects substantial variation in ARIMA's relative performance across horizons, with its performance deteriorating from short to intermediate horizons before improving at the longest horizon. The result therefore indicates horizon-dependent performance, not strong predictive ability. In particular, rejecting equality across horizons should not be interpreted as evidence that ARIMA consistently outperforms the HM benchmark.

Other models show weaker statistical evidence. TSMixer, TCN, and SVR have unadjusted $p$-values below 0.05, but their Holm-adjusted $p$-values are 0.1620, 0.2520, and 0.4031, respectively. Consequently, their horizon-specific differences do not remain statistically significant after accounting for multiple comparisons.

The bandwidth sensitivity analysis provides additional evidence about the dependence of the results on the covariance specification. At $L=5$, five models---TCN, TSMixer, SVR, ARIMA, and Random Forest---remain significant after Holm correction. At $L=20$, $L=60$, and $L=119$, ARIMA is the only model that remains significant, with adjusted $p$-values of 0.0406, 0.0056, and 0.0108, respectively. Thus, ARIMA's rejection is consistent across the four examined bandwidths, whereas the evidence for the other models is sensitive to the selected HAC bandwidth.

The paired synchronized moving-block bootstrap provides a complementary finite-sample robustness assessment. Using 4,999 replications, five block lengths $b\in\{20,40,60,80,120\}$, and the baseline HAC bandwidth $L=20$, none of the twelve models is statistically significant after Holm correction at any of the five block lengths. For ARIMA, the raw bootstrap $p$-value increases from 0.0708 at $b=20$ to 0.2348 at $b=120$, while its Holm-adjusted $p$-value remains above the 5\% significance threshold throughout. At $b=60$, the bootstrap $p$-value is 0.1620, compared with the asymptotic HAC $p$-value of 0.0034. Thus, the bootstrap does not corroborate the baseline asymptotic rejection for ARIMA. The bootstrap-based decisions remain unchanged across the five examined block lengths. Nevertheless, the difference between the HAC and bootstrap results indicates that inference for ARIMA is sensitive to the statistical procedure employed. In particular, the block-length analysis provides evidence that the failure to reject the cross-horizon equality restrictions is not specific to the choice of $b=60$, although it does not establish that the asymptotic HAC result is invalid.

Overall, the H1 results provide limited evidence of horizon-dependent predictive performance. Although ARIMA exhibits statistically significant cross-horizon differences under the baseline HAC specification, this finding is not corroborated by the moving-block bootstrap across the five examined block lengths. Furthermore, the statistical significance of other models is sensitive to the selected HAC bandwidth: four additional models are significant only at $L=5$, whereas ARIMA remains significant across all four bandwidth specifications. Consequently, the evidence supports neither a general pattern of horizon-invariant performance nor a statistically robust monotonic decay in equity return predictability. Instead, the findings highlight the importance of model-specific horizon effects and the sensitivity of statistical inference to the covariance estimation and resampling procedures.

The limited statistical evidence also highlights the importance of distinguishing horizon-specific performance differences from a general monotonic relationship between forecast horizon and predictability. As discussed by Quaedvlieg~\cite{quaedvlieg2021}, evaluating multiple horizons separately can produce inconsistent conclusions because of multiple-testing concerns and increasing forecast uncertainty. Accordingly, the rejection of cross-horizon equality for ARIMA under the baseline HAC specification should not be interpreted as evidence of a systematic increase or decrease in predictability across horizons.

\FloatBarrier
\subsection{Volatility Regime-Dependent Predictability (RQ2)}\label{subsec:volatility_regimes}

This section examines whether out-of-sample (OOS) predictive performance varies across volatility regimes and forecast horizons. The analysis distinguishes descriptive regime-specific performance from formal statistical evidence, since differences in OOS $R^2$ across regimes do not necessarily imply statistically significant differences in predictive accuracy.

\medskip
\noindent\textbf{Regime-Specific Predictive Performance.}
Following the volatility classification procedure described in \hyperref[subsec:regime_classification]{Section~\ref*{subsec:regime_classification}}, regime-specific OOS $R^2$ is evaluated relative to the fixed HM benchmark within each volatility regime.

\hyperref[tab:selected_regime_oos_r2]{Table~\ref*{tab:selected_regime_oos_r2}} summarizes descriptive predictive performance across volatility regimes and forecast horizons. It reports the distribution of positive, zero, and negative OOS $R^2$ outcomes, together with the number of models achieving positive OOS $R^2$ at each horizon. For a comprehensive model-level comparison, \hyperref[tab:full_regime_oos_r2]{Table~\ref*{tab:full_regime_oos_r2}} presents the complete regime-specific OOS $R^2$ results for all twelve forecasting models across the five forecast horizons.

Across the 60 model--horizon combinations within each regime, the Normal regime records the largest number of positive OOS $R^2$ outcomes (25, 41.7\%), followed by the High regime (23, 38.3\%) and the Low regime (16, 26.7\%). Nevertheless, negative values remain prevalent across all three regimes, indicating that most model--horizon combinations do not outperform the historical-mean benchmark.

The horizon-specific distribution in \hyperref[tab:selected_regime_oos_r2]{Table~\ref*{tab:selected_regime_oos_r2}}, Panel~B, further illustrates this heterogeneity. At the one-day horizon, seven models achieve positive OOS $R^2$ in the High regime, compared with six in the Normal regime and four in the Low regime. At the five-day horizon, only one model achieves a positive value in the Low regime, while five models do so in each of the Normal and High regimes. At the 60-day horizon, the pattern changes: six models are positive in the Low regime, five in the Normal regime, and only two in the High regime. This variation suggests that the relationship between volatility conditions and predictive performance is horizon-dependent rather than monotonic.

Several model-specific results in \hyperref[tab:full_regime_oos_r2]{Table~\ref*{tab:full_regime_oos_r2}} further illustrate this heterogeneity. TSMixer is a notable example of positive performance across all three regimes at the 120-day horizon, achieving OOS $R^2$ values of 17.81\%, 20.34\%, and 15.27\% in the Low, Normal, and High regimes, respectively. At the 20-day horizon, SVR also achieves positive OOS $R^2$ in all three regimes, with values of 0.12\%, 7.08\%, and 1.39\%, respectively. At the one-day horizon, Random Forest and XGBoost likewise produce positive values across all three regimes.

However, the complete results in \hyperref[tab:full_regime_oos_r2]{Table~\ref*{tab:full_regime_oos_r2}} show that positive performance at one horizon does not imply sustained predictive ability at longer horizons. Previous research has documented that equity return predictability can be unstable over time and may disappear as market conditions evolve~\cite{harvey2021,paye2006}. SVR, for example, records positive OOS $R^2$ at the 60-day horizon in the Low and Normal regimes (3.49\% and 16.54\%, respectively), but its performance deteriorates at 120 days, reaching $-$33.02\% and $-$33.19\% in those regimes. Ridge exhibits a similar deterioration at 120 days, with OOS $R^2$ values of $-$17.17\%, $-$3.51\%, and $-$1.47\% across the Low, Normal, and High regimes.

As further illustrated by \hyperref[tab:full_regime_oos_r2]{Table~\ref*{tab:full_regime_oos_r2}}, the most pronounced adverse results occur for TCN under High volatility. Its OOS $R^2$ is $-$90.00\% at 5 days, $-$195.95\% at 20 days, $-$168.29\% at 60 days, and $-$289.23\% at 120 days. These results indicate substantial deterioration relative to the HM benchmark in this regime. To investigate the source of this deterioration, we further examine the distribution of TCN forecasts. At the 120-day horizon, the standard deviation of TCN predictions across the full holdout sample is 0.1053, compared with 0.0894 for realized returns, while the forecast--actual correlation is approximately zero ($-$0.005). The discrepancy becomes more pronounced under High volatility, where the forecast standard deviation reaches 0.1486, compared with 0.1119 for realized returns, and the mean predicted return increases to 0.1779. This combination of excessive forecast dispersion and negligible linear association with realized returns contributes to the substantial squared-error penalty relative to HM. For example, on April 21, 2022, TCN predicts a log return of 0.3937 when the realized value is $-$0.2056; on April 23, 2025, it predicts 0.7598 against a realized value of 0.2118. A similar pattern is observed at 20 days, where the forecast standard deviation is 1.16 times that of realized returns and High-regime OOS $R^2$ reaches $-$195.95\%.

This instability may be related to the model's architecture and training configuration. Unlike the regularized Ridge benchmark, which produces substantially less dispersed forecasts, the TCN implementation uses a fixed configuration with a dropout rate of 0.1 and residual convolutional blocks incorporating \texttt{BatchNorm1d}. Its final linear projection also imposes no explicit bound on forecast magnitude. These characteristics provide plausible explanations for the observed forecast dispersion, although the present results do not isolate the causal contribution of individual architectural components. Nevertheless, the evidence suggests that TCN's severe underperformance is associated not merely with reduced forecast accuracy under High volatility, but also with poorly calibrated forecast dispersion at longer horizons.

\begin{table}[htbp]
\centering
\caption{Summary of Regime-Specific Out-of-Sample Predictive Performance.}\label{tab:selected_regime_oos_r2}
\setlength{\tabcolsep}{6pt}
\begin{tabular}{llrrr}
\hline
\multicolumn{5}{l}{\textit{Panel A. Distribution of OOS $R^2$ outcomes}} \\
\hline
\multicolumn{2}{l}{Outcome} & Low & Normal & High \\
\hline
\multicolumn{2}{l}{Positive} & 16 (26.7\%) & 25 (41.7\%) & 23 (38.3\%) \\
\multicolumn{2}{l}{Zero} & 3 (5.0\%) & 3 (5.0\%) & 3 (5.0\%) \\
\multicolumn{2}{l}{Negative} & 41 (68.3\%) & 32 (53.3\%) & 34 (56.7\%) \\
\multicolumn{2}{l}{Total} & 60 (100\%) & 60 (100\%) & 60 (100\%) \\
\hline
\multicolumn{5}{l}{\textit{Panel B. Number of models with positive OOS $R^2$ by forecast horizon}} \\
\hline
\multicolumn{2}{l}{Forecast horizon} & Low & Normal & High \\
\hline
\multicolumn{2}{l}{1 day} & 4 & 6 & 7 \\
\multicolumn{2}{l}{5 days} & 1 & 5 & 5 \\
\multicolumn{2}{l}{20 days} & 1 & 4 & 6 \\
\multicolumn{2}{l}{60 days} & 6 & 5 & 2 \\
\multicolumn{2}{l}{120 days} & 4 & 5 & 3 \\
\hline
\end{tabular}
\par\vspace{4pt}\parbox{\linewidth}{\scriptsize\textit{Notes:} OOS $R^2$ is expressed as a percentage. Panel~A reports the distribution of outcomes across model--horizon combinations, while Panel~B reports the number of models with positive OOS $R^2$ at each horizon.}
\end{table}

\begin{table}[htbp]
\centering
\caption{Full Regime-Specific Out-of-Sample $R^2$ (\%) Across All Models and Forecast Horizons.}\label{tab:full_regime_oos_r2}
{\footnotesize
\setlength{\tabcolsep}{1.8pt}
\begin{tabular}{lrrr@{\hspace{3pt}}c@{\hspace{3pt}}rrr@{\hspace{3pt}}c@{\hspace{3pt}}rrr}
\hline
\multicolumn{12}{l}{\textit{Panel A. Short- and intermediate-horizon regimes (1D, 5D, and 20D)}} \\
\hline
 & \multicolumn{3}{c}{1D} & & \multicolumn{3}{c}{5D} & & \multicolumn{3}{c}{20D} \\
\cline{2-4}\cline{6-8}\cline{10-12}
Model & Low & Norm. & High & & Low & Norm. & High & & Low & Norm. & High \\
\hline
ARIMA         & $+$0.05  & $-$0.47  & $-$0.59  & & $-$2.74  & $-$7.60  & $-$9.49  & & $-$6.67  & $-$5.99  & $-$17.15 \\
ARIMAX        & $-$0.01  & $+$0.85  & $+$1.55  & & $-$9.19  & $-$6.26  & $-$5.18  & & $-$4.38  & $-$8.76  & $-$2.47 \\
Ridge         & $-$1.06  & $+$0.53  & $+$2.41  & & $-$1.30  & $+$1.09  & $+$2.64  & & $-$1.95  & $+$0.52  & $+$2.81 \\
Elastic Net   & $+$0.15  & $+$0.11  & $+$0.13  & & $-$0.19  & $+$0.70  & $+$0.70  & & 0.00     & 0.00     & 0.00 \\
SVR           & $-$0.21  & $+$0.43  & $+$3.10  & & $-$1.18  & $+$0.89  & $+$2.21  & & $+$0.12  & $+$7.08  & $+$1.39 \\
Random Forest & $+$0.07  & $+$0.11  & $+$2.85  & & $-$0.07  & $+$0.22  & $-$2.24  & & $-$1.88  & $-$1.68  & $+$0.51 \\
XGBoost       & $+$0.05  & $+$0.02  & $+$2.42  & & $+$0.00  & $-$0.33  & $+$0.34  & & $-$0.38  & $-$0.99  & $+$2.88 \\
LSTM          & $-$2.70  & $-$1.72  & $+$0.22  & & $-$6.75  & $-$0.03  & $+$0.31  & & $-$0.12  & $+$3.88  & $+$0.46 \\
Transformer   & $-$22.78 & $-$29.70 & $-$7.20  & & $-$63.97 & $-$70.18 & $-$20.80 & & $-$38.53 & $-$29.49 & $+$1.86 \\
TCN           & $-$18.44 & $-$17.19 & $-$11.97 & & $-$39.39 & $-$22.23 & $-$90.00 & & $-$27.69 & $-$13.20 & $-$195.95 \\
PatchTST      & $-$0.74  & $-$4.60  & $-$4.71  & & $-$1.24  & $-$1.27  & $-$10.51 & & $-$14.62 & $-$8.48  & $-$13.73 \\
TSMixer       & $-$2.92  & $-$1.32  & $-$3.34  & & $-$2.74  & $+$2.20  & $-$9.67  & & $-$0.10  & $+$15.14 & $-$1.55 \\
\hline
\end{tabular}

\vspace{6pt}
\setlength{\tabcolsep}{5pt}
\begin{tabular}{lrrr@{\hspace{8pt}}c@{\hspace{8pt}}rrr}
\hline
\multicolumn{8}{l}{\textit{Panel B. Longer-horizon regimes (60D and 120D)}} \\
\hline
 & \multicolumn{3}{c}{60D} & & \multicolumn{3}{c}{120D} \\
\cline{2-4}\cline{6-8}
Model & Low & Norm. & High & & Low & Norm. & High \\
\hline
ARIMA         & $-$1.30  & $-$10.76 & $-$12.41  & & $-$7.43  & $+$11.61 & $-$1.99 \\
ARIMAX        & $-$0.53  & $-$4.13  & $+$8.09   & & $-$15.27 & $+$1.28  & $+$3.99 \\
Ridge         & $+$1.12  & $+$4.23  & $-$1.51   & & $-$17.17 & $-$3.51  & $-$1.47 \\
Elastic Net   & 0.00     & 0.00     & 0.00      & & 0.00     & 0.00     & 0.00 \\
SVR           & $+$3.49  & $+$16.54 & $-$3.49   & & $-$33.02 & $-$33.19 & $-$2.98 \\
Random Forest & $-$1.99  & $-$4.22  & $+$4.21   & & $+$1.10  & $-$11.74 & $-$5.12 \\
XGBoost       & $+$0.13  & $+$0.33  & $-$0.01   & & $-$0.41  & $-$0.83  & $+$0.49 \\
LSTM          & $+$0.71  & $-$2.32  & $-$0.46   & & $+$1.90  & $+$4.19  & $-$11.12 \\
Transformer   & $-$23.95 & $-$9.91  & $-$28.61  & & $-$21.17 & $-$28.80 & $-$19.26 \\
TCN           & $-$9.80  & $-$3.43  & $-$168.29 & & $-$26.22 & $-$9.23  & $-$289.23 \\
PatchTST      & $+$2.21  & $+$3.52  & $-$46.59  & & $+$8.29  & $+$0.57  & $-$29.79 \\
TSMixer       & $+$1.73  & $+$16.42 & $-$2.22   & & $+$17.81 & $+$20.34 & $+$15.27 \\
\hline
\end{tabular}}
\par\vspace{4pt}\parbox{\linewidth}{\scriptsize\textit{Notes:} Regime-specific OOS $R^2$ is expressed as a percentage relative to the Historical Mean (HM) benchmark. Regime sample sizes are Low ($n=332$), Normal ($n=599$), and High ($n=431$). For XGBoost at 5D under Low volatility, the unrounded OOS $R^2$ is strictly positive ($+0.0007\%$), rounding to $+0.00\%$ at two decimal places, whereas $0.00$ for Elastic Net at 20D, 60D, and 120D reflects forecasts identical to HM.}
\end{table}

\noindent\textbf{Statistical Evaluation of Regime-Dependent Predictive Performance.}
The GW-CPA test is applied to evaluate regime-conditional predictive ability relative to the Historical Mean (HM) benchmark across Low, Normal, and High volatility regimes, following the procedure described in \hyperref[subsec:hypothesis_testing]{Section~\ref*{subsec:hypothesis_testing}}.

The results identify nine statistically significant GW-CPA rejections after Holm correction, as reported in \hyperref[tab:h2_gwcpa_sig]{Table~\ref*{tab:h2_gwcpa_sig}}. These GW-CPA results provide evidence of regime-conditional predictive ability relative to the Historical Mean benchmark: a statistically significant joint rejection indicates that the expected normalized loss differential is not jointly equal to zero across the three volatility regimes. However, this result does not independently establish statistically significant differences in predictive performance between individual volatility regimes. Importantly, all nine significant cases exhibit negative full-sample mean differentials, as well as negative mean differentials within each of the Low, Normal, and High volatility regimes. This consistent direction indicates that the significant joint rejections are associated with underperformance relative to the HM benchmark across all three volatility regimes. Nevertheless, the regime-specific mean differentials are descriptive and do not establish statistically significant underperformance within each regime individually.

The significant cases are concentrated in four models. TCN is significant at all five forecast horizons from 1 to 120 days. Transformer is significant at the 1-day and 5-day horizons, while ARIMA is significant at 5 days. SVR is significant at 120 days. No model demonstrates statistically significant outperformance over HM after the multiple-testing correction.

This distinction is particularly important for interpreting TSMixer at 120 days. Although TSMixer achieves positive OOS $R^2$ in all three regimes, its GW-CPA result does not remain statistically significant after Holm correction. Thus, the descriptive evidence of predictive gains should not be interpreted as statistically established superiority over the benchmark. This distinction is consistent with the broader forecasting literature, which emphasizes the importance of formal predictive accuracy testing rather than relying exclusively on descriptive forecast performance~\cite{diebold1995,giacomini2006,giacomini2010}.

\begin{table}[htbp]
\centering
\caption{Statistically Significant GW-CPA Test Results After Holm Correction.}\label{tab:h2_gwcpa_sig}
\setlength{\tabcolsep}{1.8pt}
{\scriptsize
\begin{tabular}{llrrrrrrrr}
\hline
 & & & & & & \multicolumn{4}{c}{Mean differential} \\
\cline{7-10}
Model & Horizon & Wald stat. & $\mathrm{df}$ & Raw $p$-value & Holm $p$-value & Full sample & Low & Normal & High \\
\hline
Transformer & 1D   & 30.9605 & 3 & 0.000001 & 0.000010 & $-$0.152796 & $-$0.100996 & $-$0.191921 & $-$0.138322 \\
TCN         & 1D   & 40.6554 & 3 & $7.66\times 10^{-9}$ & $<$0.000001 & $-$0.141560 & $-$0.081764 & $-$0.111082 & $-$0.229980 \\
ARIMA       & 5D   & 17.8814 & 3 & 0.000465 & 0.004653 & $-$0.082484 & $-$0.011318 & $-$0.050938 & $-$0.181146 \\
Transformer & 5D   & 23.9464 & 3 & 0.000026 & 0.000308 & $-$0.397171 & $-$0.264519 & $-$0.470621 & $-$0.397273 \\
TCN         & 5D   & 23.6924 & 3 & 0.000029 & 0.000319 & $-$0.649099 & $-$0.162898 & $-$0.149032 & $-$1.718609 \\
TCN         & 20D  & 21.2389 & 3 & 0.000094 & 0.001127 & $-$1.234276 & $-$0.104509 & $-$0.093654 & $-$3.689764 \\
TCN         & 60D  & 16.5454 & 3 & 0.000876 & 0.010517 & $-$0.872323 & $-$0.056391 & $-$0.027822 & $-$2.674518 \\
SVR         & 120D & 13.0872 & 3 & 0.004452 & 0.048970 & $-$0.190970 & $-$0.236622 & $-$0.271492 & $-$0.043895 \\
TCN         & 120D & 14.5199 & 3 & 0.002276 & 0.027318 & $-$1.425479 & $-$0.187887 & $-$0.075516 & $-$4.254965 \\
\hline
\end{tabular}
}
\par\vspace{4pt}\parbox{\linewidth}{\scriptsize\textit{Notes:} The tests use three degrees of freedom for non-degenerate cases. Positive normalized squared-error differentials indicate improvement over HM, whereas negative values indicate underperformance. The full-sample mean differential is calculated over the entire holdout sample, while the Low, Normal, and High columns report the corresponding regime-specific mean differentials. Holm--Bonferroni adjustment is conducted separately within each forecast horizon across the twelve models. Regime-specific mean differentials are reported descriptively and are not subject to separate significance tests.}
\end{table}

\paragraph{Sensitivity to alternative HAC bandwidths.} To assess whether the GW-CPA findings depend on the Newey--West bandwidth, the tests are repeated using $L \in \{5,20,60,119\}$. Six model--horizon combinations retain statistical significance after Holm correction across the examined specifications: ARIMA at 5 days, TCN at 1, 5, and 20 days, and Transformer at 1 and 5 days. In contrast, the significant baseline results for TCN and SVR at the 120-day horizon do not remain significant when the bandwidth is increased to $L=60$, with Holm-adjusted $p$-values exceeding 0.20. This loss of significance indicates that the long-horizon GW-CPA findings are sensitive to the bandwidth used to account for serial dependence. Overall, the results provide more consistent evidence at shorter and intermediate horizons than at the 120-day horizon.

\noindent\textbf{Horizon $\times$ Regime Interaction.}
While the GW-CPA test evaluates conditional predictive ability relative to the HM benchmark across volatility regimes, the horizon--regime interaction test examines whether the relative performance contrasts between volatility regimes vary across forecast horizons relative to the 1-day baseline.

The complete interaction-test results are presented in \hyperref[tab:h2_interaction]{Table~\ref*{tab:h2_interaction}}. Under the baseline specification, TCN is the only model with a statistically significant horizon--regime interaction after Holm correction ($p_{\mathrm{Holm}}=0.002514$). Specifically, this joint rejection suggests that at least one of the Low--Normal or High--Normal contrasts differs from its corresponding 1-day baseline contrast, providing evidence of horizon-dependent regime effects without establishing direct differences in predictive performance between volatility regimes at any individual horizon. Random Forest exhibits a relatively strong unadjusted result, but its adjusted $p$-value of 0.059990 does not meet the 5\% significance threshold.

To examine the sensitivity of this finding to the Newey--West bandwidth, the interaction tests are repeated using $L \in \{5,20,60,119\}$. TCN remains statistically significant after Holm correction across all four bandwidth specifications, with adjusted $p$-values below 0.03. It is therefore the only model whose horizon--regime interaction remains significant across the examined bandwidth choices. This consistency suggests that the evidence of horizon-dependent regime effects for TCN is not driven solely by the baseline bandwidth of 20. Nevertheless, the result does not establish direct differences in predictive performance between volatility regimes at any individual horizon, because the interaction test evaluates changes in regime contrasts relative to the 1-day baseline.

The TCN results illustrate the economic magnitude of this interaction. Its normalized High--Normal differential gap changes from $-$0.1189 at the 1-day horizon to $-$1.5696 at 5 days and $-$4.1795 at 120 days. At the 120-day horizon, its regime-specific \texttt{Mean\_Skill} is $-$0.0755 in the Normal regime and $-$4.2550 in the High regime. These results indicate that TCN's relative performance deteriorates disproportionately under High volatility as the forecast horizon increases.

\begin{table}[htbp]
\centering
\caption{Horizon $\times$ Regime Interaction Test Results.}\label{tab:h2_interaction}
\setlength{\tabcolsep}{6pt}
\begin{tabular}{lrrrr}
\hline
Model & Wald statistic & $\mathrm{df}$ & Raw $p$-value & Holm $p$-value \\
\hline
ARIMA         & 17.1180 & 8 & 0.028904 & 0.289037 \\
ARIMAX        &  9.2826 & 8 & 0.319019 & 1.000000 \\
Ridge         &  6.4439 & 8 & 0.597638 & 1.000000 \\
Elastic Net   &  3.4982 & 4 & 0.478153 & 1.000000 \\
SVR           &  9.4479 & 8 & 0.305929 & 1.000000 \\
Random Forest & 21.7240 & 8 & 0.005454 & 0.059990 \\
XGBoost       & 14.1973 & 8 & 0.076766 & 0.690897 \\
LSTM          & 12.1211 & 8 & 0.145880 & 1.000000 \\
Transformer   & 10.2816 & 8 & 0.245815 & 1.000000 \\
TCN           & 30.0223 & 8 & 0.000209 & 0.002514 \\
PatchTST      &  9.5971 & 8 & 0.294449 & 1.000000 \\
TSMixer       &  8.6626 & 8 & 0.371553 & 1.000000 \\
\hline
\end{tabular}
\par\vspace{4pt}\parbox{\linewidth}{\scriptsize\textit{Notes:} The interaction test examines whether the Low--Normal and High--Normal contrasts at the four non-baseline horizons differ from the corresponding contrasts at the 1-day horizon. The test has eight degrees of freedom for non-degenerate models. Elastic Net has four degrees of freedom because its forecast differentials are identically zero at the 20-, 60-, and 120-day horizons. Holm--Bonferroni correction is applied across the twelve models. Statistical significance is assessed at the 5\% level.}
\end{table}

\noindent\textbf{Distinguishing Regime-Specific OOS $R^2$ from \texttt{Mean\_Skill}.}
An important methodological distinction concerns the normalization used in regime-specific OOS $R^2$ and the \texttt{Mean\_Skill} statistic used in the GW-CPA analysis.

Regime-specific OOS $R^2$ uses the HM mean squared error calculated within the corresponding volatility regime. In contrast, \texttt{Mean\_Skill} uses the full-holdout HM mean squared error as a common normalization denominator for all regimes. Consequently, although both measures reflect the direction of relative squared-error performance, their magnitudes are not directly interchangeable because they use different benchmark normalization scales.

The relationship between the two measures has been verified across all 180 model--horizon--regime combinations. The difference arises because the regime-specific HM benchmark error can be substantially higher or lower than the full-holdout benchmark error.

This distinction is illustrated by TSMixer at the 120-day horizon. Its OOS $R^2$ is highest in the Normal regime (20.34\%), followed by Low (17.81\%) and High (15.27\%). However, its \texttt{Mean\_Skill} is highest in the High regime ($+$0.224579), followed by Normal ($+$0.166346) and Low ($+$0.127670). The difference arises because the High-regime HM MSE is approximately 1.471 times the full-holdout HM MSE, compared with ratios of approximately 0.818 in Normal and 0.717 in Low.

Accordingly, OOS $R^2$ is used to describe regime-specific predictive performance relative to the benchmark within each regime, whereas \texttt{Mean\_Skill} is used in the conditional predictive ability tests with a common full-holdout normalization. The two statistics serve different analytical purposes and should not be compared as though they were expressed on the same scale.

\medskip
\noindent\textbf{Overall Findings for RQ2.}
Overall, the findings indicate that the relationship between volatility regimes and equity return predictability is heterogeneous and dependent on both forecasting models and investment horizons. Previous studies have documented regime-dependent differences in equity return predictability~\cite{baltas2018,hammerschmid2018}, as well as the temporal instability of predictive relationships~\cite{giacomini2010,harvey2021,paye2006}. Although the Normal regime records the largest number of positive OOS $R^2$ outcomes, the descriptive results do not establish a universally superior volatility regime. The formal statistical tests provide further qualification: GW-CPA identifies significant underperformance relative to the historical-mean benchmark, while the horizon $\times$ regime interaction test identifies statistically significant variation in regime differences only for TCN under the baseline specification. Taken together, these findings suggest that volatility conditions alone do not systematically determine forecasting performance; their association with predictive accuracy depends on the interaction between model characteristics and forecast horizons.

\subsubsection{Robustness to Alternative Volatility Window Definitions}\label{subsec:volatility_robustness}

To assess whether the findings for RQ2 depend on the 20-day rolling window used in the baseline regime classification, we repeat the regime-dependent evaluation using 10-day and 60-day rolling standard deviations. Both alternative specifications retain the expanding historical-quantile procedure described in Section~\ref{subsec:regime_classification} and use the same date-aligned holdout sample of 1,362 observations. The forecasting models, predictor sets, and evaluation framework remain unchanged.

The resulting regime distributions are broadly comparable across the three specifications, although the relative proportions of observations assigned to each regime vary with the window length. Under the 10-day specification, the holdout sample contains 354 Low-volatility observations (25.99\%), 567 Normal-volatility observations (41.63\%), and 441 High-volatility observations (32.38\%). The baseline 20-day specification contains 332, 599, and 431 observations in the respective regimes, while the 60-day specification contains 315, 662, and 385 observations. Thus, the longer window assigns a larger proportion of observations to the Normal regime and smaller proportions to both tail regimes.

Pairwise classification agreement further illustrates the sensitivity of regime assignments to the rolling-window length. The 10-day and 20-day specifications agree on 949 of 1,362 observations (69.68\%), while the agreement between the 20-day and 60-day specifications is 67.33\% (917 observations). Agreement between the two alternative specifications is lower, at 58.00\%. Importantly, direct switches between Low and High classifications are rare when either alternative is compared with the baseline: only 3 observations switch directly between these two regimes for the 10-day versus 20-day comparison, and 12 observations do so for the 20-day versus 60-day comparison. Most classification differences therefore involve reassignment between adjacent volatility regimes rather than a direct reversal between the two extremes.

Table~\ref{tab:volatility_window_robustness} summarizes the regime distributions and the main statistical results across the three rolling-window specifications. The number of statistically significant GW-CPA results after Holm correction is relatively similar, with 8, 9, and 6 significant model--horizon combinations under the 10-day, 20-day, and 60-day specifications, respectively. Under all three specifications, the significant cases indicate underperformance relative to the HM benchmark. The recurring results are concentrated in ARIMA at the 5-day horizon and Transformer at the 1-day and 5-day horizons. TCN also accounts for several significant cases, although the number of significant TCN horizons decreases under the 60-day specification. SVR at 120 days is significant under the 20-day and 60-day specifications, but not under the 10-day specification. Overall, the GW-CPA results do not indicate statistically significant predictive gains over HM under any of the three regime definitions.

\begin{table}[htbp]
\centering
\caption{Summary of Robustness Results Across Volatility Window Specifications.}\label{tab:volatility_window_robustness}
\setlength{\tabcolsep}{6pt}
\begin{tabular}{lccc}
\hline
\textbf{Measure} & \textbf{10-day} & \textbf{20-day} & \textbf{60-day} \\
\hline
\multicolumn{4}{l}{\textit{Regime distribution}} \\
Low volatility & 354 (25.99\%) & 332 (24.38\%) & 315 (23.13\%) \\
Normal volatility & 567 (41.63\%) & 599 (43.98\%) & 662 (48.60\%) \\
High volatility & 441 (32.38\%) & 431 (31.64\%) & 385 (28.27\%) \\
\hline
GW-CPA significant cases & 8/60 & 9/60 & 6/60 \\
Significant horizon--regime interactions & 0/12 & 1/12 & 2/12 \\
\hline
\end{tabular}
\par\vspace{4pt}\parbox{\linewidth}{\scriptsize\textit{Notes:} All specifications use the same 1,362-observation holdout sample. GW-CPA counts represent model--horizon combinations with statistically significant results after Holm correction. Interaction counts represent models with statistically significant horizon--regime interactions after Holm correction at the 5\% level. The 20-day specification is the baseline.}
\end{table}

The horizon--regime interaction tests provide a more differentiated picture. Under the 10-day specification, none of the twelve models exhibits a statistically significant interaction after Holm correction. Under the baseline 20-day specification, TCN is the only model with a significant interaction ($p_{\mathrm{Holm}}=0.0025$). Under the 60-day specification, both TCN ($p_{\mathrm{Holm}}=0.0324$) and Random Forest ($p_{\mathrm{Holm}}=0.0373$) exhibit significant interactions. Random Forest is close to, but does not reach, the 5\% significance threshold under the 10-day and 20-day specifications, with adjusted $p$-values of 0.0633 and 0.0600, respectively. TCN also approaches the threshold under the 10-day specification ($p_{\mathrm{Holm}}=0.0573$). These results suggest that evidence of horizon-dependent regime differences is sensitive to the volatility-window definition for some models, although TCN remains the most recurrent source of significant interaction effects.

The regime-specific OOS $R^2$ results provide additional descriptive evidence. Although the exact values and the number of positive model--horizon combinations vary across the three specifications, the main model-level patterns remain recognizable. TSMixer maintains positive OOS $R^2$ across all three regimes at the 120-day horizon under each window definition. Conversely, TCN continues to exhibit particularly severe deterioration under High volatility at longer horizons. These recurring patterns suggest that the principal descriptive findings are not solely artifacts of the baseline 20-day classification, even though the magnitude of regime-specific performance and the statistical evidence for horizon--regime interactions are not identical across specifications.

Overall, the alternative-window analyses provide qualified support for the robustness of the RQ2 findings. The main descriptive patterns remain broadly consistent across the three specifications, whereas the statistical evidence for horizon--regime interactions shows some sensitivity to the rolling-window length. Accordingly, the 20-day specification is retained as the baseline, with the 10-day and 60-day results serving as robustness checks.

\FloatBarrier
\subsection{Temporal Stability of Predictive Performance (RQ3)}\label{subsec:temporal_stability}

RQ3 examines whether out-of-sample (OOS) predictive performance remains stable across calendar time or exhibits model- and horizon-specific differences. Temporal stability is evaluated using two complementary approaches: block-level HAC Wald tests and the modified Giacomini--Rossi-type fluctuation test specified in \hyperref[subsec:hypothesis_testing]{Section~\ref*{subsec:hypothesis_testing}}.

The block-level analysis examines whether average predictive skill differs across predefined calendar periods, whereas the fluctuation test evaluates whether local predictive ability relative to the historical-mean (HM) benchmark differs from zero over rolling windows. This distinction follows the broader econometric literature emphasizing that predictive performance can vary over time and that full-sample evaluation may conceal important local changes in relative forecasting accuracy~\cite{giacomini2010}. These analyses focus on chronological variation and are conceptually distinct from the volatility-regime analysis in RQ2, which classifies observations into Low, Normal, and High volatility states.

\medskip
\noindent\textbf{Holdout Block Stability.}
Following the calendar-based subperiod definitions in \hyperref[subsec:data_splitting]{Section~\ref*{subsec:data_splitting}}, the primary holdout block-level analysis examines three non-overlapping calendar periods: 2021--2022, 2023--2024, and 2025--2026. As the holdout sample was not used for model selection, hyperparameter optimization, or feature-engineering decisions, these results provide the primary evidence for assessing temporal stability on unseen data. The November--December 2020 observations are excluded from this block-level analysis but retained in the aggregate holdout evaluation and rolling fluctuation test.

The HAC Wald tests examine whether mean normalized predictive skill is equal across the three calendar blocks. Predictive skill is measured using the normalized squared-error differential, with the full-holdout HM mean squared error serving as the common normalization denominator. Positive values indicate lower squared forecast errors than HM, whereas negative values indicate inferior predictive performance.

Under the baseline Newey--West bandwidth of $L=20$, \hyperref[tab:h3_holdout_sig]{Table~\ref*{tab:h3_holdout_sig}} reports five model--horizon combinations with statistically significant differences across blocks after Holm correction: TCN at 20D, 60D, and 120D; LSTM at 60D; and PatchTST at 60D.

To examine the sensitivity of these block-level results to the HAC bandwidth, the Wald tests are repeated using $L \in \{5,20,60,119\}$. The number of statistically significant model--horizon combinations after Holm correction decreases from 15 at $L=5$ to 5 at $L=20$, and to zero at both $L=60$ and $L=119$. This pattern indicates that the evidence of differences in average predictive skill across calendar blocks is sensitive to the bandwidth used in the HAC covariance estimator. In particular, the baseline findings should not be interpreted as uniformly robust to longer bandwidth choices.

For TCN at 120D, mean predictive skill varies substantially across blocks ($-$3.6284, 0.0679, and $-$0.1418, respectively; adjusted $p=0.0047$ under the baseline $L=20$ specification), indicating considerable temporal heterogeneity in the baseline analysis. However, the broader bandwidth sensitivity results qualify the statistical strength of this evidence.

In contrast, TSMixer at 120D shows no statistically significant block-level variation (Wald = 0.6653; adjusted $p=1.0000$). Its mean skill values are 0.2113, 0.0860, and 0.1472, with within-block OOS $R^2$ values of 15.90\%, 11.53\%, and 18.12\%, respectively. Although the mean skill values differ descriptively, the Wald test does not reject equality across blocks, suggesting comparatively consistent predictive performance over the three evaluation periods.

The interpretation of TSMixer's block-level stability at 120D requires an additional qualification. As discussed in \hyperref[subsec:multiseed_robustness]{Section~\ref*{subsec:multiseed_robustness}}, its favorable performance under the baseline seed-42 configuration is not consistently reproduced across independent random initializations. Therefore, the absence of statistically significant differences across calendar blocks should be interpreted as temporal consistency conditional on a particular fitted model, rather than evidence of initialization-robust predictive performance.

\begin{table}[htbp]
\centering
\caption{Significant Holdout Block HAC Wald Test Results.}\label{tab:h3_holdout_sig}
\setlength{\tabcolsep}{4.5pt}
\begin{tabular}{llrrrr}
\hline
Model & Horizon & Wald statistic & $\mathrm{df}$ & Raw $p$-value & Holm-adjusted $p$-value \\
\hline
TCN & 20D & 10.9722 & 2 & 0.0041 & 0.0497 \\
TCN & 60D & 17.0106 & 2 & 0.0002 & 0.0024 \\
TCN & 120D & 15.6699 & 2 & 0.0004 & 0.0047 \\
LSTM & 60D & 14.3866 & 2 & 0.0008 & 0.0083 \\
PatchTST & 60D & 10.6183 & 2 & 0.0049 & 0.0495 \\
\hline
\end{tabular}
\par\vspace{4pt}\parbox{\linewidth}{\scriptsize\textit{Notes:} Holm correction is applied separately within each forecasting horizon across the twelve competing models. Statistical significance is assessed at the 5\% level.}
\end{table}

\noindent\textbf{Validation Block Stability.}
As a supplementary analysis, the validation sample is evaluated across the three predefined calendar periods to examine temporal variation during model development and validation. The HAC Wald tests examine equality of mean normalized predictive skill across these periods. Because the validation sample contributes to model selection and hyperparameter optimization, these results are not treated as independent evidence of out-of-sample generalization. Instead, they provide contextual evidence on temporal variation within the model development process.

Although some model--horizon combinations exhibit nominal significance, none remains significant after Holm correction. The smallest adjusted $p$-value is 0.1184 for PatchTST at 120D. Thus, the validation results provide no statistically significant evidence of systematic differences in average predictive skill across the three calendar periods.

\medskip
\noindent\textbf{Rolling Predictive Performance: Giacomini--Rossi-Type Fluctuation Test.}
To complement the block-level analysis, the modified Giacomini--Rossi-type fluctuation test specified in \hyperref[subsec:hypothesis_testing]{Section~\ref*{subsec:hypothesis_testing}} is applied to examine temporal variations in local predictive ability relative to the HM benchmark.

The analysis uses trailing windows of 252 trading days and Newey--West HAC standardization estimated over the full holdout sample. The baseline specification uses a bandwidth of 20 observations, while sensitivity analyses examine $L\in\{5,20,60,119\}$. For each bandwidth, the HAC scaling factor is held fixed across the observed rolling statistic and all bootstrap replications. Baseline statistical inference is conducted using 500 moving-block bootstrap replications with a block length of 60 observations. In each replication, the rolling fluctuation statistics and the corresponding two-sided supremum statistic are recalculated from the resampled demeaned differential series. The supremum statistic is evaluated over all 1,111 rolling windows in the holdout sample, and Holm correction is applied across the twelve forecasting models separately for each forecast horizon.

Additional bootstrap sensitivity analyses examine two dimensions of the resampling procedure. First, the replication count is increased to 1,999 and 4,999 while retaining the baseline block length of 60. Second, an alternative block length of 120 is examined using 500 and 4,999 replications. These configurations retain the same observed statistic, null-centering procedure, rolling-window construction, and HAC standardization as the baseline implementation. Holm-adjusted $p$-values are recalculated within each forecast-horizon family for every configuration.

An important feature of this sensitivity analysis is that the same positive HAC scaling factor is applied to both the observed supremum statistic and the bootstrap supremum statistics. Consequently, changing the bandwidth rescales both quantities proportionally and does not change the bootstrap exceedance event. The resulting empirical $p$-values and rejection decisions are therefore algebraically invariant to the examined HAC bandwidths under this implementation. This invariance should not be interpreted as independent empirical confirmation that the fluctuation-test results are robust to alternative dependence structures; rather, it follows from the common scaling factor used in the observed and bootstrap statistics.

\hyperref[tab:h3_fluctuation_sig]{Table~\ref*{tab:h3_fluctuation_sig}} reports the baseline fluctuation statistics and corresponding $p$-values after multiple-testing adjustment. Under the baseline specification ($B=500$, $b=60$, and $L=20$), the null hypothesis is rejected after Holm correction for Transformer at 1D, ARIMA at 20D, and TCN at 60D. These findings indicate statistically significant departures from zero local predictive ability relative to the HM benchmark in at least one 252-observation evaluation window under the baseline specification. However, supplementary bootstrap sensitivity analyses reveal that the statistical evidence for some model--horizon combinations depends on the selected block length and replication count.

The same three rejections are obtained under all four HAC bandwidth specifications, $L\in\{5,20,60,119\}$. The corresponding empirical bootstrap $p$-values remain unchanged because the HAC scaling factor is common to the observed and bootstrap supremum statistics. Specifically, the baseline raw empirical $p$-values are $p<0.002$ for Transformer at 1D, $p=0.002$ for ARIMA at 20D, and $p=0.004$ for TCN at 60D; the corresponding Holm-adjusted $p$-values are below 0.002, 0.024, and 0.048, respectively. These unchanged $p$-values reflect the algebraic structure of the bootstrap comparison rather than an independent test of sensitivity to alternative HAC bandwidths.

The additional bootstrap sensitivity analyses provide a more direct assessment of the dependence of these findings on the resampling configuration. For TCN at 60D, increasing the number of replications while retaining $b=60$ produces Holm-adjusted $p$-values of 0.0600 and 0.0480 for $B=1,999$ and $B=4,999$, respectively. Thus, the result remains close to the 5\% significance threshold and does not exhibit a uniformly significant outcome across replication counts. More importantly, increasing the block length to $b=120$ yields a raw $p$-value of 0.0606 and a Holm-adjusted $p$-value of 0.7274 with $B=4,999$. Consequently, the baseline rejection for TCN at 60D is sensitive to block-length selection and should be interpreted as suggestive rather than robust evidence of local predictive instability.

A similar pattern is observed for ARIMA at 20D. Under $b=60$, its Holm-adjusted $p$-value is 0.0480 with $B=4,999$, whereas increasing the block length to $b=120$ raises the adjusted $p$-value to 0.1056. This indicates that its baseline evidence of local predictive instability is also sensitive to the resampling block length.

In contrast, Transformer at 1D produces no bootstrap exceedances across the examined replication counts and both block lengths. Its empirical $p$-value therefore remains below the corresponding resolution threshold, providing substantially more consistent evidence of departure from zero local predictive ability across the tested bootstrap configurations. Nevertheless, the absence of observed exceedances should not be interpreted as a zero probability or as proof of robustness under all possible resampling specifications.

Importantly, rejection does not necessarily imply that predictive performance alternates between superior and inferior outcomes over time. The rolling statistics provide further context for interpreting these results.

For Transformer at 1D, the rolling statistics remain negative across all 1,111 windows, ranging from $-$4.8926 to $-$0.2339, indicating persistently inferior predictive performance relative to HM. Its significant result therefore reflects sustained underperformance rather than alternating periods of relative predictive advantage.

The persistent underperformance of Transformer at 1D can be further understood from its forecast distribution. The standard deviation of its predictions is approximately 0.0042, only 0.40 times the standard deviation of realized returns (0.0105), while the forecast--actual correlation is just 0.031. Thus, despite producing relatively low-variance forecasts, Transformer captures almost none of the linear variation in realized returns. Its regime-specific OOS $R^2$ is negative across all three volatility regimes, as reported in \hyperref[tab:full_regime_oos_r2]{Table~\ref*{tab:full_regime_oos_r2}}, indicating that its persistent rolling underperformance is not confined to a particular volatility state. Rather, its forecasts fail to provide sufficient information about realized return movements to improve upon the HM benchmark. The low forecast dispersion should therefore not be interpreted as evidence of effective regularization or reliable predictive behavior.

For ARIMA at 20D and TCN at 60D, the rolling statistics range from $-$3.3391 to 0.4237 and from $-$7.5441 to 0.1435, respectively. Negative statistics account for 88.39\% of ARIMA's windows and 77.77\% of TCN's windows, indicating that their local predictive performance is predominantly inferior to HM. These descriptive patterns are consistent with predominantly negative local predictive performance, although the bootstrap sensitivity analyses indicate that the statistical evidence of temporal instability for both models depends on the selected block length. In particular, neither ARIMA at 20D nor TCN at 60D retains Holm-adjusted significance under the alternative block length of 120.

\begin{table}[htbp]
\centering
\caption{Significant Giacomini--Rossi-Type Fluctuation Test Results.}\label{tab:h3_fluctuation_sig}
\setlength{\tabcolsep}{5pt}
\begin{tabular}{llrrr}
\hline
Model & Horizon & Supremum statistic & Raw $p$-value & Holm-adjusted $p$-value \\
\hline
Transformer & 1D  & 4.8926 & $<$0.002 & $<$0.002 \\
ARIMA       & 20D & 3.3391 & 0.002    & 0.024 \\
TCN         & 60D & 7.5441 & 0.004    & 0.048 \\
\hline
\end{tabular}
\par\vspace{4pt}\parbox{\linewidth}{\scriptsize\textit{Notes:} Results are based on the baseline specification with 500 moving-block bootstrap replications and a block length of 60. Empirical $p$-values are calculated as the proportion of bootstrap supremum statistics greater than or equal to the observed statistic. The minimum non-zero empirical $p$-value resolution is 0.002; zero-exceedance results are reported as $p<0.002$. Holm correction is applied separately within each forecast horizon across the twelve candidate models using unrounded empirical $p$-values. Supplementary sensitivity analyses using alternative replication counts ($B=1,999$ and $4,999$) and block length ($b=120$) are discussed in Section 5.4.}
\end{table}

Descriptively, TSMixer at 120D records an average 252-day rolling normalized predictive skill of 0.138 (SD = 0.174), compared with $-$1.670 (SD = 2.414) for TCN. These statistics are calculated as the mean and sample standard deviation of rolling normalized predictive skill across the 1,111 windows. They are not within-window OOS $R^2$ measures, as their normalization uses the full-holdout HM MSE rather than the HM MSE calculated separately within each rolling window. The substantially larger dispersion for TCN indicates pronounced temporal variation in its local predictive performance.

The rolling instability observed for TCN complements the forecast-distribution diagnostics reported in \hyperref[subsec:volatility_regimes]{Section~\ref*{subsec:volatility_regimes}}, which identify excessive forecast dispersion and negligible forecast--actual correlation at longer horizons. In the present analysis, the large dispersion of its rolling statistics indicates substantial temporal variation in forecast-error performance rather than necessarily reflecting alternating periods of strong predictive ability. Transformer at 1D exhibits a different failure pattern: its forecast distribution is considerably less dispersed than realized returns, yet its predictive correlation is also negligible. These results illustrate distinct forms of underperformance: pronounced forecast dispersion and temporal variation for TCN, versus insufficient predictive information despite low forecast dispersion for Transformer.

The comparatively favorable rolling statistics of TSMixer at 120D should be interpreted in conjunction with the multi-seed robustness analysis in \hyperref[subsec:multiseed_robustness]{Section~\ref*{subsec:multiseed_robustness}}. These statistics characterize the temporal behavior of the baseline seed-42 configuration only and should not be interpreted as evidence that the architecture consistently maintains favorable local predictive performance across independent random initializations.

\medskip
\noindent\textbf{Overall Assessment of RQ3.}
The results provide qualified support for H3, although the evidence depends on the temporal dimension being examined. Block-level HAC Wald tests identify significant differences in average predictive skill across calendar periods for five model--horizon combinations in the holdout sample, whereas validation block tests yield no significant results after Holm correction.

The Giacomini--Rossi-type fluctuation test identifies three model--horizon combinations with statistically significant departures from zero local predictive ability relative to HM under the baseline bootstrap specification. However, supplementary sensitivity analyses qualify the strength of this evidence. Transformer at 1D exhibits no bootstrap exceedances across the examined replication counts and block lengths, providing consistent evidence of persistent underperformance. In contrast, the statistical significance of ARIMA at 20D and TCN at 60D does not persist when the block length is increased to 120. Their baseline rejections should therefore be interpreted as specification-sensitive evidence rather than robust findings of local predictive instability. Furthermore, rejection does not necessarily imply alternating periods of superior and inferior predictive performance; Transformer at 1D, for example, exhibits persistent underperformance across all rolling windows rather than alternating predictive outcomes.

Overall, temporal performance characteristics are model- and horizon-dependent rather than universal. The holdout block-level and rolling-window analyses provide complementary evidence on temporal stability, while the validation block results offer supplementary context on temporal variation during model development. These findings highlight the importance of evaluating forecasting models through both chronological block comparisons and local rolling-window assessments, rather than relying exclusively on aggregate OOS performance.

The supplementary multi-seed results in \hyperref[subsec:multiseed_robustness]{Section~\ref*{subsec:multiseed_robustness}} further qualify these findings by demonstrating that temporal consistency within a particular fitted model does not necessarily imply reproducibility across independent random initializations. Taken together, the results provide qualified support for H3: predictive performance exhibits model- and horizon-specific temporal variation, while initialization sensitivity represents an additional and distinct dimension of forecasting robustness. Therefore, chronological stability and reproducibility across random initializations should be considered jointly when evaluating financial forecasting models.


\subsection{Multi-Seed Robustness of Deep Learning Models}\label{subsec:multiseed_robustness}

\noindent\textbf{Multi-Seed Robustness of Deep Learning Models.} Following the multi-seed evaluation protocol described in \hyperref[subsec:model_config]{Section~\ref*{subsec:model_config}}, this subsection examines the robustness of deep learning forecasting performance across independent random initializations. The analysis covers five architectures: TSMixer, LSTM, PatchTST, TCN, and Vanilla Transformer. TSMixer is evaluated using ten independent random seeds, with its results summarized separately in \hyperref[tab:tsmixer_multiseed]{Table~\ref*{tab:tsmixer_multiseed}}, whereas the remaining four architectures are evaluated using five random seeds each, with their results reported in \hyperref[tab:dl_multiseed_robustness]{Table~\ref*{tab:dl_multiseed_robustness}}. Panel~A presents the mean, standard deviation, and median OOS $R^2$, while Panel~B reports the minimum and maximum OOS $R^2$, together with MAE, RMSE, and directional accuracy (DA).

\begin{table}[htbp]
\centering
\caption{Multi-seed robustness of TSMixer across forecast horizons.}\label{tab:tsmixer_multiseed}
\setlength{\tabcolsep}{4.5pt}
\begin{tabular}{lrrrrrr}
\hline
Horizon & Mean OOS $R^2$ & SD & Median & Min & Max & Positive seeds \\
\hline
1D   & $-$1.67\%  & 0.98\%  & $-$1.55\%  & $-$3.52\%  & $-$0.28\% & 0/10 \\
5D   & $-$2.73\%  & 2.15\%  & $-$2.57\%  & $-$6.34\%  & $-$0.11\% & 0/10 \\
20D  & $-$6.96\%  & 10.10\% & $-$4.08\%  & $-$25.34\% & +4.53\%   & 2/10 \\
60D  & $-$21.49\% & 15.15\% & $-$23.45\% & $-$38.26\% & +4.91\%   & 1/10 \\
120D & $-$18.93\% & 19.09\% & $-$20.10\% & $-$44.09\% & +17.48\%  & 1/10 \\
\hline
\end{tabular}
\par\vspace{4pt}\parbox{\linewidth}{\scriptsize\textit{Notes:} OOS $R^2$ is expressed as a percentage. Positive seeds indicate the number of runs with OOS $R^2$ greater than zero. TSMixer was evaluated using ten random seeds with fixed hyperparameters.}
\end{table}

The multi-seed results for TSMixer reveal substantial initialization sensitivity, particularly at longer forecast horizons. As reported in \hyperref[tab:tsmixer_multiseed]{Table~\ref*{tab:tsmixer_multiseed}}, TSMixer achieves negative mean OOS $R^2$ at every horizon, ranging from $-$1.67\% at 1D to $-$21.49\% at 60D. Performance dispersion also increases with the forecast horizon, with the standard deviation rising from 0.98\% at 1D to 19.09\% at 120D. The proportion of runs achieving positive OOS $R^2$ remains low: none of the ten runs produces positive predictive skill at either 1D or 5D, while only two runs at 20D and one run at each of 60D and 120D outperform the historical-mean (HM) benchmark. These results indicate that the predictive performance observed under a particular initialization is not consistently reproduced across independent training runs.

The discrepancy is particularly evident at 120D. The baseline TSMixer model reported in \hyperref[subsec:oos_performance]{Section~\ref*{subsec:oos_performance}} achieves an aggregate OOS $R^2$ of +17.53\% at the 120-day horizon. In the independent multi-seed experiment, the separately retrained seed-42 run achieves +17.48\%. Although both runs use the same nominal random seed, minor differences in GPU-based training execution lead to slightly different model predictions. Accordingly, the original baseline result is retained in the main performance analysis, while the multi-seed statistics are calculated from the independently retrained runs without replacing their original results. Although the baseline configuration exhibits positive within-block OOS $R^2$ across all three holdout periods in \hyperref[subsec:temporal_stability]{Section~\ref*{subsec:temporal_stability}}, the ten-seed evaluation produces a mean OOS $R^2$ of $-$18.93\% (SD = 19.09\%), with only one of ten runs achieving positive predictive skill. The median of $-$20.10\% further indicates that the negative mean is not driven solely by an isolated unfavorable run. Moreover, the mean OOS $R^2$ across the ten seeds is negative in all three calendar blocks ($-$6.18\%, $-$14.61\%, and $-$64.87\%, respectively). Thus, the favorable performance observed under seed 42 does not represent the typical outcome across alternative initializations. This discrepancy highlights the distinction between the performance of an individual training run and the expected performance across repeated initializations. As demonstrated by Coakley and Gundersen~\cite{coakley2026}, deep learning performance can exhibit substantial variability across random seeds, with the magnitude and distribution of this variability differing across architectures and tasks. Consequently, an isolated positive OOS result should not be interpreted as sufficient evidence of reproducible predictive performance.

\begin{table}[htbp]
\centering
\caption{Multi-seed robustness of deep learning models across forecast horizons.}\label{tab:dl_multiseed_robustness}
\renewcommand{\arraystretch}{0.96}
\setlength{\tabcolsep}{4.5pt}
\textbf{Panel A: Out-of-sample $R^2$ distribution across random seeds}\\[3pt]
\begin{tabular}{llrrr}
\hline
Model & Horizon & Mean OOS $R^2$ & SD & Median \\
\hline
LSTM & 1D   & $+$0.80\% & 1.85\% & $-$0.10\% \\
LSTM & 5D   & $+$0.19\% & 1.09\% & $+$0.55\% \\
LSTM & 20D  & $-$0.75\% & 2.03\% & $-$0.49\% \\
LSTM & 60D  & $-$2.14\% & 2.13\% & $-$1.64\% \\
LSTM & 120D & $-$5.36\% & 5.46\% & $-$3.13\% \\
\hline
PatchTST & 1D   & $-$8.98\% & 6.77\% & $-$6.05\% \\
PatchTST & 5D   & $-$10.03\% & 4.62\% & $-$8.81\% \\
PatchTST & 20D  & $-$23.97\% & 10.50\% & $-$26.36\% \\
PatchTST & 60D  & $-$23.78\% & 14.03\% & $-$27.21\% \\
PatchTST & 120D & $-$23.43\% & 12.84\% & $-$20.26\% \\
\hline
TCN & 1D   & $-$14.62\% & 4.77\% & $-$14.44\% \\
TCN & 5D   & $-$54.69\% & 19.34\% & $-$57.72\% \\
TCN & 20D  & $-$83.30\% & 54.60\% & $-$59.17\% \\
TCN & 60D  & $-$75.42\% & 51.96\% & $-$85.88\% \\
TCN & 120D & $-$83.54\% & 52.63\% & $-$70.73\% \\
\hline
Vanilla Transformer & 1D   & $-$18.32\% & 11.31\% & $-$19.61\% \\
Vanilla Transformer & 5D   & $-$26.26\% & 14.35\% & $-$21.10\% \\
Vanilla Transformer & 20D  & $-$7.83\% & 3.78\% & $-$8.93\% \\
Vanilla Transformer & 60D  & $-$25.60\% & 18.62\% & $-$18.90\% \\
Vanilla Transformer & 120D & $-$13.13\% & 18.44\% & $-$6.28\% \\
\hline
\end{tabular}

\vspace{8pt}
\setlength{\tabcolsep}{3.2pt}
\textbf{Panel B: Range bounds, forecast errors, and directional accuracy}\\[3pt]
\begin{tabular}{llrrrrr}
\hline
Model & Horizon & Min $R^2$ & Max $R^2$ & MAE & RMSE & DA \\
\hline
LSTM & 1D   & $-$0.86\% & $+$3.06\% & 0.0076 & 0.0105 & 49.16\% \\
LSTM & 5D   & $-$1.27\% & $+$1.37\% & 0.0165 & 0.0223 & 54.61\% \\
LSTM & 20D  & $-$3.95\% & $+$1.48\% & 0.0330 & 0.0418 & 64.45\% \\
LSTM & 60D  & $-$5.85\% & $-$0.51\% & 0.0536 & 0.0662 & 74.76\% \\
LSTM & 120D & $-$13.81\% & $+$0.47\% & 0.0759 & 0.0949 & 78.05\% \\
\hline
PatchTST & 1D   & $-$19.12\% & $-$2.24\% & 0.0080 & 0.0110 & 49.77\% \\
PatchTST & 5D   & $-$16.84\% & $-$5.74\% & 0.0172 & 0.0234 & 54.73\% \\
PatchTST & 20D  & $-$32.66\% & $-$7.19\% & 0.0352 & 0.0463 & 56.27\% \\
PatchTST & 60D  & $-$40.97\% & $-$7.42\% & 0.0548 & 0.0728 & 72.00\% \\
PatchTST & 120D & $-$41.36\% & $-$11.86\% & 0.0792 & 0.1027 & 74.80\% \\
\hline
TCN & 1D   & $-$21.83\% & $-$10.30\% & 0.0083 & 0.0113 & 48.21\% \\
TCN & 5D   & $-$79.12\% & $-$29.41\% & 0.0209 & 0.0277 & 47.75\% \\
TCN & 20D  & $-$162.60\% & $-$34.06\% & 0.0422 & 0.0559 & 48.68\% \\
TCN & 60D  & $-$129.96\% & $-$21.27\% & 0.0644 & 0.0860 & 64.95\% \\
TCN & 120D & $-$139.12\% & $-$26.02\% & 0.0906 & 0.1243 & 72.66\% \\
\hline
Vanilla Transformer & 1D   & $-$34.81\% & $-$6.24\% & 0.0085 & 0.0115 & 49.85\% \\
Vanilla Transformer & 5D   & $-$47.03\% & $-$11.52\% & 0.0189 & 0.0251 & 51.09\% \\
Vanilla Transformer & 20D  & $-$11.79\% & $-$3.75\% & 0.0341 & 0.0433 & 53.67\% \\
Vanilla Transformer & 60D  & $-$58.01\% & $-$12.12\% & 0.0563 & 0.0733 & 66.68\% \\
Vanilla Transformer & 120D & $-$43.34\% & $+$4.69\% & 0.0752 & 0.0982 & 74.90\% \\
\hline
\end{tabular}
\par\vspace{4pt}\parbox{\linewidth}{\scriptsize\textit{Notes:} OOS $R^2$ values in Panel A, minimum and maximum bounds in Panel B, and directional accuracy (DA) are expressed as percentages. SD denotes the standard deviation across random seeds. All metrics are computed across five independent random-seed training runs for each architecture and forecast horizon.}
\end{table}

The results in \hyperref[tab:dl_multiseed_robustness]{Table~\ref*{tab:dl_multiseed_robustness}} reveal substantial differences in initialization robustness across the remaining four architectures. LSTM exhibits the most favorable aggregate OOS $R^2$ performance among these models, with slightly positive mean values at 1D and 5D (+0.80\% and +0.19\%, respectively). However, its mean OOS $R^2$ becomes negative at longer horizons, declining to $-$5.36\% at 120D. The standard deviation also increases from 1.85\% at 1D to 5.46\% at 120D, indicating greater variation across initializations at the longest horizon. The median and minimum values further show that the positive short-horizon means do not imply uniformly positive performance across individual runs.

PatchTST produces negative mean OOS $R^2$ at all five horizons. Its mean performance is particularly weak at 20D, 60D, and 120D, with values of $-$23.97\%, $-$23.78\%, and $-$23.43\%, respectively. The maximum OOS $R^2$ remains negative for every horizon, indicating that none of the five evaluated initializations achieves positive aggregate predictive skill relative to the HM benchmark. The standard deviations, especially at 60D and 120D, also indicate non-negligible variation in performance across runs. Thus, the underperformance of PatchTST is not attributable to a single unfavorable initialization in this experiment.

TCN exhibits the weakest and most variable results among the evaluated architectures. Its mean OOS $R^2$ is negative across all horizons, ranging from $-$14.62\% at 1D to $-$83.54\% at 120D. The deterioration is particularly pronounced at 5D, 20D, 60D, and 120D. At 20D, for example, the mean OOS $R^2$ is $-$83.30\%, while the standard deviation reaches 54.60\%. The minimum value of $-$162.60\% further illustrates the severity of the model's sensitivity to initialization at this horizon. Although the standard deviations are large, the maximum OOS $R^2$ remains negative at every horizon. This suggests that the observed weakness is systematic across the evaluated runs rather than being driven solely by one extreme seed.

Vanilla Transformer also records negative mean OOS $R^2$ at all horizons. Its weakest average performance occurs at 5D and 60D, at $-$26.26\% and $-$25.60\%, respectively. At 20D, the mean OOS $R^2$ improves to $-$7.83\%, accompanied by a comparatively smaller standard deviation of 3.78\%. Nevertheless, the maximum OOS $R^2$ remains negative at this horizon. At 120D, the maximum becomes positive (+4.69\%), but the mean remains negative ($-$13.13\%) and the standard deviation is relatively large (18.44\%). This indicates that a favorable result can occur under a particular initialization without representing the typical outcome across runs.

The error and directional metrics provide additional context for interpreting these results across the evaluated architectures. Across all models, MAE and RMSE generally increase with the forecast horizon, reflecting the greater difficulty of predicting returns over longer periods. Directional accuracy also tends to increase at longer horizons, particularly for LSTM, PatchTST, and TCN. However, higher DA does not necessarily correspond to positive OOS $R^2$. For example, LSTM achieves a mean DA of 78.05\% at 120D while its mean OOS $R^2$ remains $-$5.36\%. This distinction indicates that directional performance and relative squared-error performance capture different aspects of forecasting quality and should not be treated as interchangeable measures.

Comparing across the five architectures highlights the divergence between baseline point estimates and multi-seed distributions. TSMixer and Vanilla Transformer both show that an isolated run with positive performance (such as the retrained seed-42 outcome of +17.48\% for TSMixer at 120D or a maximum $R^2$ of +4.69\% for Transformer at 120D) can coexist with negative mean predictive performance across repeated runs. Conversely, architectures with tighter dispersion, such as LSTM, provide greater stability across initializations, although their relative advantage remains horizon-dependent.

Overall, the multi-seed evaluation demonstrates substantial differences in initialization robustness across deep learning architectures and forecasting horizons. LSTM exhibits comparatively stronger aggregate OOS performance, although its predictive skill remains sensitive to the forecast horizon and random initialization. In contrast, TSMixer, PatchTST, TCN, and Vanilla Transformer generally produce negative mean OOS $R^2$ across the evaluated horizons, with varying degrees of performance dispersion. TSMixer provides a particularly illustrative case: its favorable baseline performance at 120D is not representative of its typical outcome across ten independent runs. These findings indicate that conclusions drawn from a single training initialization may overstate the reproducibility of deep learning forecasting performance. They also complement the temporal stability analysis in \hyperref[subsec:temporal_stability]{Section~\ref*{subsec:temporal_stability}} by distinguishing chronological consistency within a fitted model from reproducibility across independent training runs.

These findings are consistent with Coakley and Gundersen~\cite{coakley2026}, who demonstrate that deep learning performance exhibits substantial variability across random initializations and that the resulting performance distributions can differ considerably across architectures and tasks. Their empirical analysis, based on 18,600 training runs across 186 experiments, further shows that performance distributions frequently deviate from normality, indicating that mean performance and standard deviation alone may not fully characterize initialization sensitivity. The present study extends this perspective to equity return forecasting by examining initialization sensitivity across multiple forecast horizons and evaluating its implications for OOS predictability relative to the historical-mean benchmark. In particular, the results demonstrate that favorable performance under an individual seed may substantially overstate the reproducibility of predictive skill when evaluated against alternative initializations.

\subsection{Feature-Group Ablation Analysis}\label{subsec:ablation_study}

\begin{table}[htbp]
\centering
\caption{Feature ablation results for XGBoost ($\Delta \mathrm{OOS}\,R^2$, percentage points).}\label{tab:ablation_xgboost}
{\footnotesize
\setlength{\tabcolsep}{6pt}
\begin{tabular}{lrrrrr}
\hline
Feature specification & 1D & 5D & 20D & 60D & 120D \\
\hline
Full Model & 0.00 & 0.00 & 0.00 & 0.00 & 0.00 \\
Without Lagged Returns & $-$0.32 & $+$0.06 & $+$0.67 & 0.00 & $-$0.05 \\
Without Momentum & $+$0.05 & $-$0.11 & $+$0.51 & 0.00 & $-$0.05 \\
Without Historical Volatility & $-$0.67 & $-$0.54 & $+$1.20 & 0.00 & $-$0.05 \\
Without VIX & $+$0.01 & $+$0.53 & $+$0.65 & 0.00 & $+$0.18 \\
Without Intraday & $-$1.20 & $+$0.09 & $+$0.75 & 0.00 & $-$0.05 \\
Without Trading Volume & $+$0.60 & $+$0.21 & $+$0.37 & 0.00 & 0.00 \\
\hline
\end{tabular}
}
\par\vspace{4pt}\parbox{\linewidth}{\scriptsize\textit{Notes:} Values are calculated as the OOS $R^2$ of the ablated specification minus that of the full model. Positive values indicate an improvement after removing the corresponding feature group; negative values indicate a reduction. All values are expressed in percentage points.}
\end{table}

The XGBoost results indicate that the contribution of feature groups varies across forecast horizons. At 1D, removing Intraday produces the largest negative change in OOS $R^2$ ($-$1.20 percentage points), followed by removing Historical Volatility ($-$0.67 percentage points) and Lagged Returns ($-$0.32 percentage points). These results indicate that removing Intraday, Historical Volatility, and Lagged Returns leads to lower short-horizon predictive performance in this XGBoost specification. In particular, the deterioration following the removal of Intraday is larger than that associated with any other feature group at 1D.

At 5D, removing Historical Volatility also reduces OOS $R^2$ by 0.54 percentage points, while removing Momentum produces a smaller negative change of 0.11 percentage points. In contrast, removing VIX or Trading Volume improves performance by 0.53 and 0.21 percentage points, respectively. At 20D, all six ablated specifications produce non-negative changes, ranging from $+$0.37 to $+$1.20 percentage points. The largest improvement occurs when Historical Volatility is removed ($+$1.20 percentage points), followed by Intraday ($+$0.75 percentage points) and Lagged Returns ($+$0.67 percentage points). This pattern suggests that the full set of predictors does not uniformly improve medium-horizon forecasting performance for XGBoost.

At 120D, the changes are comparatively small. Removing VIX improves OOS $R^2$ by 0.18 percentage points, whereas removing Lagged Returns, Momentum, or Historical Volatility produces a small deterioration of 0.05 percentage points. The removal of Trading Volume produces no reported change.

At 60D, all six ablated specifications produce identical out-of-sample predictions to the full specification within the ablation experiment, resulting in an unchanged OOS $R^2$ ($\Delta \mathrm{OOS}\,R^2 = 0.00$ percentage points). Diagnostic refitting using the selected 60D hyperparameters further confirms that the XGBoost model produces constant predictions across the test observations: all 900 trees consist exclusively of root leaf nodes, with no feature-based splits. This structural result was reproduced across repeated runs under both single-threaded and multi-threaded configurations. Although a small difference in the constant prediction level exists between the original evaluation and the diagnostic refitting, it does not affect the within-experiment ablation comparisons. Consequently, the unchanged performance should be interpreted as a property of the fitted XGBoost specification in this experiment, rather than as evidence that the six feature groups contain no predictive information at the 60D horizon.

\begin{table}[htbp]
\centering
\caption{Feature ablation results for Ridge ($\Delta \mathrm{OOS}\,R^2$, percentage points).}\label{tab:ablation_ridge}
{\footnotesize
\setlength{\tabcolsep}{6pt}
\begin{tabular}{lrrrrr}
\hline
Feature specification & 1D & 5D & 20D & 60D & 120D \\
\hline
Full Model & 0.00 & 0.00 & 0.00 & 0.00 & 0.00 \\
Without Lagged Returns & $+$0.10 & $-$0.04 & $-$0.03 & $-$0.01 & $+$0.05 \\
Without Momentum & $+$0.04 & $+$0.29 & $-$0.08 & $-$1.27 & $-$0.76 \\
Without Historical Volatility & $-$0.37 & $-$0.38 & $-$0.55 & $+$1.82 & $+$4.30 \\
Without VIX & $-$0.50 & $-$0.44 & $-$1.13 & $+$0.93 & $+$6.53 \\
Without Intraday & $-$0.07 & $-$0.20 & $-$0.07 & $-$0.25 & $-$0.02 \\
Without Trading Volume & $-$0.07 & 0.00 & 0.00 & 0.00 & $-$0.01 \\
\hline
\end{tabular}
}
\par\vspace{4pt}\parbox{\linewidth}{\scriptsize\textit{Notes:} Values are calculated as the OOS $R^2$ of the ablated specification minus that of the full model. Positive values indicate an improvement after removing the corresponding feature group; negative values indicate a reduction. All values are expressed in percentage points.}
\end{table}

Ridge exhibits a different feature-contribution pattern from XGBoost, particularly at longer horizons. At 1D, removing VIX reduces OOS $R^2$ by 0.50 percentage points, while removing Historical Volatility produces a reduction of 0.37 percentage points. These are the largest negative changes at this horizon, suggesting that both volatility-related feature groups contribute positively to the short-horizon performance of Ridge. Removing Intraday or Trading Volume also results in small negative changes, whereas removing Lagged Returns or Momentum produces slight improvements.

At 5D and 20D, the removal of VIX and Historical Volatility continues to reduce OOS $R^2$. The largest deterioration at 20D occurs when VIX is removed ($-$1.13 percentage points), followed by Historical Volatility ($-$0.55 percentage points). This pattern indicates that removing volatility-related predictors leads to larger reductions in Ridge performance at the shorter and intermediate horizons considered in the experiment.

The results change substantially at longer horizons. At 60D, removing Historical Volatility improves OOS $R^2$ by 1.82 percentage points, while removing VIX improves it by 0.93 percentage points. At 120D, these improvements become considerably larger, reaching $+$4.30 percentage points without Historical Volatility and $+$6.53 percentage points without VIX. The latter is the largest positive change in the Ridge ablation results. Given that the full Ridge model has a negative OOS $R^2$ of $-$5.01\% at 120D, the positive changes indicate that the corresponding ablated specifications perform better than the full model, although the magnitude of the improvement should be interpreted relative to the weak full-model baseline.

Momentum also displays horizon-dependent effects. Removing it produces a small improvement at 1D and a larger improvement at 5D, but reduces OOS $R^2$ at 60D and 120D. By comparison, removing Intraday results in small negative changes across all five horizons, suggesting a relatively consistent positive contribution within the Ridge specification. The removal of Trading Volume produces changes close to zero throughout, indicating limited measured contribution to Ridge performance in this experiment.

\begin{table}[htbp]
\centering
\caption{Feature ablation results for LSTM ($\Delta \mathrm{OOS}\,R^2$, percentage points).}\label{tab:ablation_lstm}
{\footnotesize
\setlength{\tabcolsep}{6pt}
\begin{tabular}{lrrrrr}
\hline
Feature specification & 1D & 5D & 20D & 60D & 120D \\
\hline
Full Model & 0.00 & 0.00 & 0.00 & 0.00 & 0.00 \\
Without Lagged Returns & $+$1.28 & $+$0.92 & $-$2.95 & $+$3.09 & $+$3.94 \\
Without Momentum & $+$0.75 & $+$0.44 & $-$0.92 & $+$4.91 & $+$3.08 \\
Without Historical Volatility & $+$2.15 & $+$1.62 & $-$0.33 & $+$1.99 & $+$6.19 \\
Without VIX & $+$0.96 & $+$0.99 & $-$1.52 & $+$0.56 & $+$9.64 \\
Without Intraday & $+$1.57 & $+$2.01 & $-$0.60 & $+$2.08 & $+$6.44 \\
Without Trading Volume & $+$1.69 & $+$1.62 & $-$1.73 & $+$1.18 & $-$2.86 \\
\hline
\end{tabular}
}
\par\vspace{4pt}\parbox{\linewidth}{\scriptsize\textit{Notes:} Values are calculated as the OOS $R^2$ of the ablated specification minus that of the full model. Positive values indicate an improvement after removing the corresponding feature group; negative values indicate a reduction. All values are expressed in percentage points.}
\end{table}

The LSTM ablation results show a pronounced horizon-dependent pattern. At 1D and 5D, removing any of the six feature groups improves OOS $R^2$. At 1D, the largest improvement is obtained by removing Historical Volatility ($+$2.15 percentage points), followed by Trading Volume ($+$1.69 percentage points) and Intraday ($+$1.57 percentage points). At 5D, removing Intraday yields the largest improvement ($+$2.01 percentage points), while removing Historical Volatility or Trading Volume improves OOS $R^2$ by 1.62 percentage points in each case. These results suggest that, for the evaluated LSTM configuration, the full predictor set does not provide the most favorable specification at the two shortest horizons.

The pattern differs at 20D. Removing any of the six feature groups reduces OOS $R^2$, with the largest deterioration occurring when Lagged Returns are removed ($-$2.95 percentage points). Removing Trading Volume and VIX also produces negative changes of $-$1.73 and $-$1.52 percentage points, respectively. This indicates that removing these feature groups has a more detrimental effect on the LSTM's intermediate-horizon performance than on its performance at 1D and 5D.

At 60D, removing any feature group except Trading Volume improves OOS $R^2$. The largest improvement occurs when Momentum is removed ($+$4.91 percentage points), followed by Lagged Returns ($+$3.09 percentage points). Removing Historical Volatility, Intraday, and VIX also produces positive changes. At 120D, the improvements associated with removing VIX ($+$9.64 percentage points), Intraday ($+$6.44 percentage points), and Historical Volatility ($+$6.19 percentage points) are particularly pronounced. However, removing Trading Volume reduces OOS $R^2$ by 2.86 percentage points, indicating that this group retains a positive contribution at the longest horizon in the evaluated LSTM specification.

Taken together, the ablation results show that the performance effects of feature-group removal are model- and horizon-dependent rather than uniform across forecasting configurations. XGBoost shows relatively small changes in most specifications, with more pronounced effects at 1D and 20D. Ridge exhibits substantial performance differences following the removal of VIX and Historical Volatility at longer horizons, particularly at 120D. LSTM displays the strongest variation across feature groups and horizons, with widespread improvements following feature removal at short and long horizons but deteriorations at 20D.

These findings do not imply that a feature group is universally useful or harmful. The observed performance effects depend on the forecasting model, forecast horizon, and other predictors retained in the specification. As discussed by Hooker, Mentch, and Zhou~\cite{hooker2021}, correlations among predictors can produce a compensation effect, whereby variables containing overlapping information partially offset the loss associated with perturbing or removing another predictor. Consequently, the predictive contribution measured for an individual variable depends partly on the information retained in the remaining predictors. A positive $\Delta \mathrm{OOS}\,R^2$ indicates that removing the corresponding feature group improves performance under the evaluated experimental setup; it does not establish that the removed variables are intrinsically uninformative. Similarly, a negative change indicates that the full specification outperforms its corresponding ablated version, but does not establish a causal relationship between the removed predictors and forecasting performance.

An additional consideration is the potential overlap in predictive information among feature groups. When predictors are correlated, the information retained in the remaining features may partially offset the loss associated with removing a particular group. This issue is closely related to the masking effect documented in permutation feature importance research. Kaneko~\cite{kaneko2022} shows that correlations among predictors can lead to underestimated individual feature importance because other correlated variables may preserve information relevant to the prediction task. Similarly, Hooker, Mentch, and Zhou~\cite{hooker2021} demonstrate that unrestricted permutation can introduce extrapolation bias by generating feature combinations outside the observed data distribution. These findings highlight two distinct limitations of conventional permutation-based importance measures: masking due to information overlap and bias arising from unrealistic perturbed observations. Although these studies examine permutation-based importance rather than the model-retraining procedure employed in the present feature-group ablation analysis, they provide complementary evidence that predictive contributions should be interpreted in the context of the dependence structure among retained predictors.

The present analysis differs from conventional permutation-based importance measures in that each feature group is removed from the model specification and the corresponding model is retrained using the reduced predictor set. Therefore, the reported $\Delta \mathrm{OOS}\,R^2$ measures the change in predictive performance between separately fitted model specifications rather than the immediate performance deterioration caused by perturbing inputs to a fixed model. This distinction is important because permutation-based measures evaluate the model under modified input distributions, whereas the present ablation procedure evaluates how predictive performance changes after model re-estimation with fewer predictors. Accordingly, the results are interpreted as the incremental predictive contribution of each feature group within a particular model--horizon configuration, rather than as an intrinsic or causal measure of feature importance.

Overall, the ablation analysis provides complementary evidence to the horizon- and regime-dependent results in Sections~\ref{subsec:oos_performance} and~\ref{subsec:volatility_regimes}. It suggests that the effectiveness of predictor groups depends not only on the forecast horizon but also on the model architecture. Consequently, the results support evaluating predictor contributions at the model--horizon level rather than assuming that a single feature configuration is optimal across all forecasting tasks.

\FloatBarrier

\subsection{Robustness to Alternative Benchmark Specifications: Zero-Return Benchmark}\label{subsec:benchmark_robustness}

To examine whether the evaluation of predictive performance is sensitive to benchmark selection, this section compares the primary Historical Mean (HM) benchmark with an alternative Zero-Return benchmark. Both specifications are evaluated using the same holdout observations for each model--horizon pair. The HM-based results are reported in \hyperref[tab:oos_r2_main]{Table~\ref*{tab:oos_r2_main}} (\hyperref[subsec:oos_performance]{Section~\ref*{subsec:oos_performance}}), while this section focuses on the alternative Zero-Return specification and its implications for benchmark sensitivity.

\paragraph{Comparison of out-of-sample performance across benchmarks.} Table~\ref{tab:benchmark_comparison} reports out-of-sample \(R^2\) values relative to the Zero-Return benchmark. For comparison, the corresponding results relative to the primary Historical Mean (HM) benchmark are presented in \hyperref[tab:oos_r2_main]{Table~\ref*{tab:oos_r2_main}} (\hyperref[subsec:oos_performance]{Section~\ref*{subsec:oos_performance}}).

Under the Zero-Return benchmark, more model--horizon combinations exhibit positive $\mathrm{OOS}\,R^2$ values, particularly at 60D and 120D. At 120D, for example, TSMixer, ARIMA, and Elastic Net achieve $\mathrm{OOS}\,R^2$ values of 38.631\%, 27.033\%, and 25.582\%, respectively. Nevertheless, these values measure performance relative to a zero-return forecast and should not be interpreted as evidence that these models outperform the primary HM benchmark. As the benchmark determines the reference forecast error in the denominator of $\mathrm{OOS}\,R^2$, a change in benchmark can alter the resulting performance measure without changing the model's underlying predictions or absolute forecast errors~\cite{campbell2008,gu2020,welch2008}. Therefore, positive $\mathrm{OOS}\,R^2$ values against Zero-Return do not necessarily indicate positive predictive performance against HM. The comparison illustrates that conclusions about predictive performance can depend materially on the benchmark specification.

\begin{table}[htbp]
\centering
\caption{Out-of-sample $R^2$ relative to the zero-return benchmark.}\label{tab:benchmark_comparison}
{\small
\setlength{\tabcolsep}{5pt}
\begin{tabular}{lrrrrr}
\hline
Model & 1D & 5D & 20D & 60D & 120D \\
\hline
ARIMA          & $-$0.286 & $-$7.102 &  $-$7.414 &   6.408 & 27.033 \\
ARIMAX         &  1.379 & $-$4.782 &   0.298 &  17.257 & 25.322 \\
Ridge          &  1.695 &  2.826 &   6.270 &  15.888 & 21.900 \\
Elastic Net    &  0.327 &  1.662 &   4.691 &  15.119 & 25.582 \\
SVR            &  2.181 &  2.523 &   7.596 &  19.057 & 11.370 \\
Random Forest  &  1.968 & $-$0.224 &   4.316 &  15.400 & 20.806 \\
XGBoost        &  1.676 &  1.167 &   6.001 &  15.229 & 25.476 \\
LSTM           & $-$0.449 &  0.560 &   6.097 &  14.305 & 23.098 \\
Transformer    & $-$15.049 & $-$38.237 &  $-$6.400 &  $-$2.945 &  8.447 \\
TCN            & $-$13.927 & $-$63.163 & $-$112.947 & $-$58.924 & $-$80.500 \\
PatchTST       & $-$4.042 & $-$5.722 &  $-$6.909 &  $-$3.438 & 16.491 \\
TSMixer        & $-$2.513 & $-$4.355 &   8.308 &  19.348 & 38.631 \\
\hline
\end{tabular}}
\par\vspace{4pt}\parbox{\linewidth}{\linespread{1.2}\selectfont\footnotesize\textit{Notes:} Values are expressed as percentages and calculated relative to the Zero-Return benchmark. The corresponding HM-based results are reported in \hyperref[tab:oos_r2_main]{Table~\ref*{tab:oos_r2_main}}.\par}
\end{table}

\paragraph{Differences between benchmark-based evaluations.} Table~\ref{tab:benchmark_difference} reports the difference between the out-of-sample \(R^2\) values obtained using the Zero-Return and Historical Mean (HM) benchmarks. The difference is positive for all 60 model--horizon combinations. Averaged across models, the gap increases from 0.205 percentage points at 1D to 1.173 at 5D, 5.291 at 20D, 16.733 at 60D, and 29.687 at 120D. The overall average difference is 10.618 percentage points.

This widening gap at longer horizons indicates that the choice of benchmark has a particularly substantial effect when evaluating cumulative returns. Importantly, the difference does not represent an improvement in the forecasts themselves. It reflects the change in the benchmark error used in the denominator of the out-of-sample $R^2$ calculation, consistent with the relative nature of benchmark-based predictive evaluation~\cite{campbell2008,welch2008}. Consequently, a positive $\mathrm{OOS}\,R^2$ against Zero-Return does not necessarily imply positive predictive performance against HM. This distinction is particularly relevant when comparing forecasting models across different horizons, as the benchmark error can vary substantially with the forecast horizon. The HM benchmark is therefore retained as the primary reference throughout the study, as reported in \hyperref[tab:oos_r2_main]{Table~\ref*{tab:oos_r2_main}}, while the Zero-Return specification provides complementary evidence on the sensitivity of the evaluation results to alternative benchmark assumptions.

\begin{table}[htbp]
\centering
\caption{Difference in out-of-sample $R^2$ between zero-return and historical mean benchmarks.}\label{tab:benchmark_difference}
{\small
\setlength{\tabcolsep}{6pt}
\begin{tabular}{lrrrrr}
\hline
Model & 1D & 5D & 20D & 60D & 120D \\
\hline
ARIMA          & 0.201 & 1.146 &  5.287 & 16.671 & 25.083 \\
ARIMAX         & 0.198 & 1.122 &  4.907 & 14.738 & 25.671 \\
Ridge          & 0.197 & 1.040 &  4.613 & 14.982 & 26.847 \\
Elastic Net    & 0.200 & 1.053 &  4.691 & 15.119 & 25.582 \\
SVR            & 0.196 & 1.043 &  4.548 & 14.418 & 30.467 \\
Random Forest  & 0.197 & 1.073 &  4.709 & 15.069 & 27.223 \\
XGBoost        & 0.197 & 1.058 &  4.627 & 15.100 & 25.618 \\
LSTM           & 0.202 & 1.064 &  4.622 & 15.264 & 26.435 \\
Transformer    & 0.231 & 1.480 &  5.237 & 18.337 & 31.472 \\
TCN            & 0.229 & 1.746 & 10.481 & 28.308 & 62.048 \\
PatchTST       & 0.209 & 1.132 &  5.262 & 18.425 & 28.706 \\
TSMixer        & 0.206 & 1.117 &  4.513 & 14.366 & 21.096 \\
\hline
Average        & 0.205 & 1.173 &  5.291 & 16.733 & 29.687 \\
\hline
\end{tabular}}
\par\vspace{4pt}\parbox{\linewidth}{\linespread{1.2}\selectfont\footnotesize\textit{Notes:} Entries represent percentage-point differences calculated as $R^2_{\mathrm{Zero}} - R^2_{\mathrm{HM}}$ relative to the respective benchmarks. The overall average difference across all 60 model--horizon combinations is 10.618 percentage points.\par}
\end{table}

\paragraph{Statistical comparison.} The Diebold--Mariano tests with the Harvey--Leybourne--Newbold correction provide additional evidence on whether individual models significantly outperform each benchmark (summarized in Table~\ref{tab:dm_benchmark_summary}). Based on raw (unadjusted) $p$-values at the nominal 5\% significance level, 1 of the 60 model--horizon comparisons is significant against HM (1.7\%), corresponding to TSMixer at 120D ($p = 0.0485$). This result does not remain statistically significant after Holm adjustment, as reported in \hyperref[subsec:oos_performance]{Section~\ref*{subsec:oos_performance}}. Against Zero-Return, 8 of the 60 comparisons are nominally significant based on raw $p$-values (13.3\%). These include ARIMAX at 60D, Ridge at 5D and 20D, Random Forest at 60D, XGBoost at 20D, LSTM at 60D, TSMixer at 120D, and Elastic Net at 5D. These counts represent nominal significance before multiple-testing adjustment and should not be interpreted as evidence of statistically significant improvements after controlling for multiple comparisons.

The greater number of nominally significant comparisons against Zero-Return suggests that the statistical conclusions are also sensitive to benchmark selection. However, these raw $p$-value counts remain descriptive and do not establish statistical significance after multiple-testing adjustment. The findings do not establish uniform predictive superiority across models or forecast horizons. The HM-based statistical results are consistent with the primary evaluation reported in \hyperref[subsec:oos_performance]{Section~\ref*{subsec:oos_performance}}, where the nominally significant TSMixer result does not remain significant after Holm correction.

\begin{table}[htbp]
\centering
\caption{Summary of Diebold--Mariano test results across benchmarks.}\label{tab:dm_benchmark_summary}
{\small
\setlength{\tabcolsep}{8pt}
\begin{tabular}{lrr}
\hline
Benchmark & Nominally significant comparisons & Share \\
\hline
Historical Mean & 1/60 & 1.7\% \\
Zero-Return     & 8/60 & 13.3\% \\
\hline
\end{tabular}}
\par\vspace{4pt}\parbox{0.8\linewidth}{\linespread{1.2}\selectfont\footnotesize\textit{Notes:} One-sided Diebold--Mariano tests with the Harvey--Leybourne--Newbold correction. Reported counts are based on raw (unadjusted) $p$-values at the nominal 5\% significance level. They represent individual significance before multiple-testing adjustment and should not be interpreted as evidence of statistical significance after controlling for multiple comparisons.\par}
\end{table}

\paragraph{Benchmark sensitivity across volatility regimes.} Table~\ref{tab:benchmark_regimes} further examines the differences between the two benchmark specifications across low-, normal-, and high-volatility regimes. At shorter horizons, the average differences are relatively small. At longer horizons, however, the gaps become more pronounced, particularly under low- and normal-volatility conditions. At 120D, the difference reaches 30.79 percentage points in the low-volatility regime and 37.93 percentage points in the normal-volatility regime, compared with 16.28 percentage points in the high-volatility regime.

The regime-level results reinforce the finding that benchmark selection has a greater influence on measured performance at longer horizons. They also indicate that the magnitude of this effect varies across market conditions rather than remaining constant across volatility regimes.

\begin{table}[htbp]
\centering
\caption{Benchmark sensitivity across volatility regimes.}\label{tab:benchmark_regimes}
{\small
\setlength{\tabcolsep}{7pt}
\begin{tabular}{llrrr}
\hline
Horizon & Volatility Regime & HM (\%) & Zero-Return (\%) & Difference (pp) \\
\hline
1D   & Low    & $-$4.04  & $-$3.98  &  0.07 \\
1D   & Normal & $-$4.41  & $-$4.07  &  0.34 \\
1D   & High   & $-$1.26  & $-$1.10  &  0.17 \\
\hline
5D   & Low    & $-$10.73 & $-$8.91  &  1.82 \\
5D   & Normal & $-$8.57  & $-$6.63  &  1.93 \\
5D   & High   & $-$11.81 & $-$11.13 &  0.68 \\
\hline
20D  & Low    & $-$8.01  &  0.03  &  8.04 \\
20D  & Normal & $-$3.50  &  7.71  & 11.20 \\
20D  & High   & $-$18.41 & $-$17.63 &  0.78 \\
\hline
60D  & Low    & $-$2.35  & 17.28  & 19.63 \\
60D  & Normal &  0.52  & 26.54  & 26.02 \\
60D  & High   & $-$20.94 & $-$16.67 &  4.28 \\
\hline
120D & Low    & $-$7.63  & 23.16  & 30.79 \\
120D & Normal & $-$4.11  & 33.82  & 37.93 \\
120D & High   & $-$28.43 & $-$12.15 & 16.28 \\
\hline
\end{tabular}}
\par\vspace{4pt}\parbox{\linewidth}{\linespread{1.2}\selectfont\footnotesize\textit{Notes:} HM denotes the Historical Mean benchmark. Difference is calculated as the Zero-Return OOS $R^2$ minus the HM OOS $R^2$, in percentage points. Values are averaged across the 12 candidate models within each horizon--regime combination.\par}
\end{table}

Overall, the robustness analysis demonstrates that measured predictive performance is sensitive to the benchmark specification, especially at longer forecast horizons. Although several models achieve positive OOS $R^2$ values and nominally significant DM test results relative to Zero-Return based on raw $p$-values, these results should not be interpreted as statistically significant improvements after multiple-testing adjustment. Their performance is generally less favorable when assessed against HM. The primary conclusions of this study therefore continue to rely on the Historical Mean benchmark, as reported in \hyperref[subsec:oos_performance]{Section~\ref*{subsec:oos_performance}}, while the Zero-Return results serve as a complementary robustness check.

\section{Conclusion}\label{sec:conclusion}

This study investigates how out-of-sample (OOS) predictability of S\&P 500 price-index returns varies across forecast horizons, market volatility regimes, and chronological subperiods. Using daily S\&P 500 and VIX data from 1990 to 2026, the analysis compares twelve econometric, conventional machine-learning, and deep-learning models across five cumulative-return horizons, with the Historical Mean (HM) serving as the primary benchmark. By integrating multi-horizon evaluation, volatility-regime analysis, chronological stability tests, feature-group ablation, and multi-seed robustness assessment, the study provides a unified empirical examination of the conditions under which forecasting performance varies and the extent to which observed gains are statistically reliable and reproducible.

First, the findings indicate substantial model-specific and non-monotonic variation in predictive performance across forecast horizons. TSMixer achieves positive OOS $R^2$ values at 20, 60, and 120 days, recording its largest improvement of 17.53\% at the 120-day horizon. Nevertheless, no candidate model significantly outperforms HM after Holm correction across the 60 model--horizon comparisons. The formal cross-horizon equality tests provide only limited evidence of horizon-dependent performance. ARIMA is the only model whose equality restriction remains statistically significant under the examined HAC bandwidth specifications. However, the synchronized moving-block bootstrap does not confirm this rejection at any of the tested block lengths. Consequently, the results suggest that predictive performance can differ across horizons in a model-specific manner, but they do not establish a general monotonic relationship between forecast horizon and equity return predictability. These findings also indicate that favorable point estimates at individual horizons should not be interpreted as evidence of statistically reliable predictive superiority.

Second, the relationship between predictive performance and market volatility is heterogeneous across models and forecast horizons. The Normal-volatility regime contains the largest number of positive model--horizon OOS $R^2$ outcomes (25 of 60), followed by the High-volatility regime (23) and the Low-volatility regime (16). However, these descriptive differences do not establish that the Normal regime is statistically superior. The Giacomini--White conditional predictive ability analysis identifies nine statistically significant model--horizon combinations after Holm correction under the baseline specification, with all nine exhibiting negative reported full-sample mean differentials. Thus, these rejections indicate conditional departures from the HM benchmark rather than predictive improvements over it. Furthermore, the horizon--regime interaction test identifies TCN as the only statistically significant model under the baseline 20-day classification. Its interaction remains significant across the examined HAC bandwidth specifications, although the results of alternative volatility-window analyses indicate that the broader evidence of horizon-dependent regime differences is sensitive to the regime definition. Overall, volatility conditions are associated with model- and horizon-specific differences in forecasting performance rather than a universally advantageous market regime.

Third, the temporal stability analysis provides qualified evidence that predictive performance varies across chronological subperiods. Under the baseline HAC bandwidth of 20, block-level tests identify statistically significant differences in average predictive skill for five model--horizon combinations in the holdout sample: TCN at 20, 60, and 120 days, LSTM at 60 days, and PatchTST at 60 days. However, the number of significant results decreases substantially as the HAC bandwidth increases, with no statistically significant block-level results remaining at bandwidths of 60 and 119. The modified Giacomini--Rossi-type fluctuation analysis further identifies local departures from the HM benchmark for Transformer at 1 day, ARIMA at 20 days, and TCN at 60 days under the baseline bootstrap specification. Nevertheless, supplementary bootstrap analyses show that some findings are sensitive to the selected block length and replication configuration. In particular, the significance of ARIMA at 20 days and TCN at 60 days does not persist when the block length increases to 120. Transformer at 1 day exhibits more consistent evidence of persistent underperformance across the examined bootstrap configurations. Taken together, these findings suggest that temporal variation is model- and horizon-dependent, while the strength of statistical evidence depends on the inference procedure and its underlying specifications.

The supplementary robustness analyses further qualify the interpretation of the main results. Feature-group ablation demonstrates that the effects of removing predictor groups depend on both the forecasting model and the prediction horizon; no feature group provides a uniformly beneficial contribution across the evaluated configurations. The alternative Zero-Return benchmark produces more favorable OOS $R^2$ values than HM, particularly at longer horizons, illustrating that measured predictive performance depends on benchmark selection. Moreover, the multi-seed evaluation reveals that favorable results from an individual deep-learning initialization do not necessarily translate into reproducible performance. In particular, TSMixer achieves negative mean OOS $R^2$ across all five horizons over ten independent runs, with the mean reaching $-$18.93\% at 120 days and only one run producing a positive OOS $R^2$ at that horizon. This contrasts with the positive 17.53\% OOS $R^2$ obtained by the baseline seed-42 configuration at the same horizon. The difference highlights the importance of distinguishing the performance of a particular fitted model from the reproducibility of an architecture across independent training runs.

Several limitations should be acknowledged. First, the use of overlapping cumulative-return targets induces serial dependence, making statistical inference sensitive to covariance estimation and resampling procedures. Although multiple HAC bandwidths and synchronized moving-block bootstrap specifications are examined, the results demonstrate that some statistical conclusions remain sensitive to these choices. Second, the analysis focuses on the S\&P 500 price index and excludes dividend payments; consequently, the findings do not directly represent total-return predictability. Third, the study considers a single aggregate equity market index, limiting the generalizability of the results to other markets and asset classes. Fourth, the deep-learning architectures do not receive identical hyperparameter optimization budgets, which should be considered when interpreting comparisons across model architectures. Finally, the use of post-close information for certain predictors means that the forecasts should not be interpreted as directly executable trading signals at the corresponding closing prices. Future research could extend the analysis to total-return indices and additional equity markets, examine alternative volatility definitions, and employ more extensive dependence-robust inference procedures. Repeated-seed experiments and more balanced hyperparameter optimization protocols could also help distinguish architecture-level performance from initialization effects. In addition, incorporating realistic execution prices, transaction costs, portfolio turnover, and risk-adjusted performance measures would provide a more comprehensive assessment of whether statistical forecasting differences translate into economically meaningful investment value.

Overall, this study does not establish a universally superior forecasting model, a monotonic relationship between forecast horizon and predictability, or a consistently advantageous volatility regime. Instead, the empirical evidence demonstrates that OOS equity return forecasting performance is conditional on model specification, forecast horizon, market volatility, and evaluation period. More importantly, descriptive improvements, statistical significance, temporal consistency, and reproducibility represent distinct dimensions of forecasting performance. Evaluating these dimensions jointly is therefore essential for drawing reliable conclusions about equity return predictability.


\begin{thebibliography}{40}
\bibitem{akiba2019}
Akiba, T., Sano, S., Yanase, T., Ohta, T., Koyama, M.: Optuna: A next-generation hyperparameter optimization framework. In: Proceedings of the 25th ACM SIGKDD International Conference on Knowledge Discovery \& Data Mining, pp. 2623--2631 (2019).

\bibitem{alizadeh2002}
Alizadeh, S., Brandt, M.W., Diebold, F.X.: Range-based estimation of stochastic volatility models. The Journal of Finance \textbf{57}(3), 1047--1091 (2002).

\bibitem{balduzzi1999}
Balduzzi, P., Lynch, A.W.: Transaction costs and predictability: some utility cost calculations. Journal of Financial Economics \textbf{52}(1), 47--78 (1999).

\bibitem{baltas2018}
Baltas, N., Karyampas, D.: Forecasting the equity risk premium: The importance of regime-dependent evaluation. Journal of Financial Markets \textbf{38}, 83--102 (2018).

\bibitem{banerjee2007}
Banerjee, P.S., Doran, J.S., Peterson, D.R.: Implied volatility and future portfolio returns. Journal of Banking \& Finance \textbf{31}(10), 3183--3199 (2007).

\bibitem{bekaert2014}
Bekaert, G., Hoerova, M.: The VIX, the variance premium and stock market volatility. Journal of Econometrics \textbf{183}(2), 181--192 (2014).

\bibitem{buturac2022}
Buturac, G.: Measurement of economic forecast accuracy: A systematic overview of the empirical literature. Journal of Risk and Financial Management \textbf{15}(1), 1 (2022).

\bibitem{cakici2025}
Cakici, N., Zaremba, A.: Accounting vs technical information: what matters more for stock return predictability? Journal of International Financial Markets, Institutions and Money \textbf{104}, 102207 (2025).

\bibitem{coakley2026}
Coakley, K.L., Gundersen, O.E.: Measuring deep learning performance - an empirical study of performance distributions across architectures and tasks. Scientific Reports \textbf{16}(1), 20474 (2026).

\bibitem{campbell1993}
Campbell, J.Y., Grossman, S.J., Wang, J.: Trading volume and serial correlation in stock returns. The Quarterly Journal of Economics \textbf{108}(4), 905--939 (1993).

\bibitem{campbell1998}
Campbell, J.Y., Lo, A.W., MacKinlay, A.C.: The Econometrics of Financial Markets. Princeton University Press, Princeton (1997).

\bibitem{campbell2008}
Campbell, J.Y., Thompson, S.B.: Predicting excess stock returns out of sample: Can anything beat the historical average? The Review of Financial Studies \textbf{21}(4), 1509--1531 (2008).

\bibitem{cenesizoglu2012}
Cenesizoglu, T., Timmermann, A.: Do return prediction models add economic value? Journal of Banking \& Finance \textbf{36}(11), 2974--2987 (2012).

\bibitem{diebold1995}
Diebold, F.X., Mariano, R.S.: Comparing predictive accuracy. Journal of Business \& Economic Statistics \textbf{13}(3), 253--263 (1995).

\bibitem{fredvix}
Federal Reserve Bank of St. Louis: CBOE Volatility Index: VIX (2026), \url{https://fred.stlouisfed.org/series/VIXCLS}. Accessed 28 September 2026

\bibitem{feng2024}
Feng, Y., Zhang, Y., Wang, Y.: Out-of-sample volatility prediction: Rolling window, expanding window, or both? Journal of Forecasting \textbf{43}(3), 567--582 (2024).

\bibitem{french1987}
French, K.R., Schwert, G.W., Stambaugh, R.F.: Expected stock returns and volatility. Journal of Financial Economics \textbf{19}(1), 3--29 (1987).

\bibitem{giacomini2006}
Giacomini, R., White, H.: Tests of conditional predictive ability. Econometrica \textbf{74}(6), 1545--1578 (2006).

\bibitem{giacomini2010}
Giacomini, R., Rossi, B.: Forecast comparisons in unstable environments. Journal of Applied Econometrics \textbf{25}(4), 595--620 (2010).

\bibitem{gu2020}
Gu, S., Kelly, B., Xiu, D.: Empirical asset pricing via machine learning. The Review of Financial Studies \textbf{33}(5), 2223--2273 (2020).

\bibitem{hammerschmid2018}
Hammerschmid, R., Lohre, H.: Regime shifts and stock return predictability. International Review of Economics \& Finance \textbf{56}, 138--160 (2018).

\bibitem{hansen1980}
Hansen, L.P., Hodrick, R.J.: Forward exchange rates as optimal predictors of future spot rates: An econometric analysis. Journal of Political Economy \textbf{88}(5), 829--853 (1980).

\bibitem{harvey2021}
Harvey, D.I., Leybourne, S.J., Sollis, R., Taylor, A.M.R.: Real-time detection of regimes of predictability in the US equity premium. Journal of Applied Econometrics \textbf{36}(1), 45--70 (2021).

\bibitem{harvey1997}
Harvey, D., Leybourne, S., Newbold, P.: Testing the equality of prediction mean squared errors. International Journal of Forecasting \textbf{13}(2), 281--291 (1997).

\bibitem{hewamalage2023}
Hewamalage, H., Ackermann, K., Bergmeir, C.: Forecast evaluation for data scientists: Common pitfalls and best practices. Data Mining and Knowledge Discovery \textbf{37}(2), 788--832 (2023).

\bibitem{hodrick1992}
Hodrick, R.J.: Dividend yields and expected stock returns: Alternative procedures for inference and measurement. The Review of Financial Studies \textbf{5}(3), 357--386 (1992).

\bibitem{hodson2022}
Hodson, T.O.: Root-mean-square error (RMSE) or mean absolute error (MAE): when to use them or not. Geoscientific Model Development \textbf{15}(14), 5481--5487 (2022).

\bibitem{hooker2021}
Hooker, G., Mentch, L., Zhou, S.: Unrestricted permutation forces extrapolation: variable importance requires at least one more model, or there is no free variable importance. Statistics and Computing \textbf{31}(6), 82 (2021).

\bibitem{holm1979}
Holm, S.: A simple sequentially rejective multiple test procedure. Scandinavian Journal of Statistics \textbf{6}(2), 65--70 (1979).

\bibitem{jegadeesh1993}
Jegadeesh, N., Titman, S.: Returns to buying winners and selling losers: Implications for stock market efficiency. The Journal of Finance \textbf{48}(1), 65--91 (1993).

\bibitem{kaneko2022}
Kaneko, H.: Cross-validated permutation feature importance considering correlation between features. Analytical Science Advances \textbf{3}(9--10), 278--287 (2022).

\bibitem{kingma2014}
Kingma, D.P., Ba, J.: Adam: A method for stochastic optimization. arXiv preprint arXiv:\allowbreak 1412.6980 (2014).

\bibitem{li2017}
Li, J., Tsiakas, I.: Equity premium prediction: The role of economic and statistical constraints. Journal of Financial Markets \textbf{36}, 56--75 (2017).

\bibitem{loshchilov2018decoupled}
Loshchilov, I., Hutter, F.: Decoupled weight decay regularization. In: International Conference on Learning Representations (2019). \url{https://openreview.net/forum?id=Bkg6RiCqY7}

\bibitem{moskowitz2012}
Moskowitz, T.J., Ooi, Y.H., Pedersen, L.H.: Time series momentum. Journal of Financial Economics \textbf{104}(2), 228--250 (2012).

\bibitem{nechvatalova2024}
Nechv{\'a}talov{\'a}, L.: Multi-horizon equity returns predictability via machine learning. Czech Journal of Economics and Finance \textbf{74}(2), 142--190 (2024).

\bibitem{newey1987}
Newey, W.K., West, K.D.: A simple, positive semi-definite, heteroskedasticity and autocorrelation consistent covariance matrix. Econometrica \textbf{55}(3), 703--708 (1987).

\bibitem{perez2026}
Nu{\~n}ez Perez, F.A., Aguilar Mosqueda, F.J., Ramos Cuevas, A., Mu{\~n}oz Beltran, J., Nu{\~n}ez Perez, J.C.: Leakage-controlled horizon-specific model selection for daily equity forecasting: An automated multi-model pipeline. Forecasting \textbf{8}(2), 34 (2026).

\bibitem{paye2006}
Paye, B.S., Timmermann, A.: Instability of return prediction models. Journal of Empirical Finance \textbf{13}(3), 274--315 (2006).

\bibitem{pettenuzzo2011}
Pettenuzzo, D., Timmermann, A.: Predictability of stock returns and asset allocation under structural breaks. Journal of Econometrics \textbf{164}(1), 60--78 (2011).

\bibitem{pesaran1992}
Pesaran, M.H., Timmermann, A.: A simple nonparametric test of predictive performance. Journal of Business \& Economic Statistics \textbf{10}(4), 461--465 (1992).

\bibitem{quaedvlieg2021}
Quaedvlieg, R.: Multi-horizon forecast comparison. Journal of Business \& Economic Statistics \textbf{39}(1), 40--53 (2021).

\bibitem{sheikholeslami2021}
Sheikholeslami, S., Meister, M., Wang, T., Payberah, A.H., Vlassov, V., Dowling, J.: AutoAblation: Automated parallel ablation studies for deep learning. In: Proceedings of the 1st Workshop on Machine Learning and Systems (EuroMLSys '21), pp. 55--61. ACM, New York (2021).

\bibitem{tashman2000}
Tashman, L.J.: Out-of-sample tests of forecasting accuracy: An analysis and review. International Journal of Forecasting \textbf{16}(4), 437--450 (2000).

\bibitem{welch2008}
Welch, I., Goyal, A.: A comprehensive look at the empirical performance of equity premium prediction. The Review of Financial Studies \textbf{21}(4), 1455--1508 (2008).

\bibitem{xu2020}
Xu, K.-L.: Testing for multiple-horizon predictability: Direct regression based versus implication based. The Review of Financial Studies \textbf{33}(9), 4403--4443 (2020).

\bibitem{yahoosandp}
Yahoo Finance: S\&P 500 (\textasciicircum GSPC) historical data (2026), \url{https://finance.yahoo.com/quote/%5EGSPC/history/}. Accessed 29 September 2026
\end{thebibliography}

\end{document}
